\clearpage
\section{Parametricity family diagrams: declarations and proofs}
\label{sec:parametricity-families-concordance}
These declarations construct the original complete family copying and beta diagrams, their original source isomorphisms and directed corrections. They also establish search-root transport and retained allocation and numeric branches. They do not yet constitute universal parametricity robustness.

\subsection{Allocation Enabled}
Allocation enabledness follows from the actual search history and the unchanged local stable-value condition. The operand may be a boundary lambda whose box belongs to an arbitrary larger template component.

The actual ordered unary argument is a stable primitive when a rooted search has reached this ref request. The complete two-edge match is constructed.
For the proof, constructor and boundary cases are separated explicitly; the remaining natural-number index bounds follow by linear arithmetic; the defined maps and incidence equations are substituted.
\lean{Copying/AllocationEnabled}{15}{61}{Allocation.rooted_enabled}
\par\penalty0

The enabled original rule reaches exactly its canonical ref-edge replacement, retaining the whole argument and every surrounding edge.
For the proof, the preceding component lemmas are composed or specialized at the stated parameters.
\lean{Copying/AllocationEnabled}{65}{70}{Allocation.rooted_canonical}
\par\penalty0

\subsection{Boundary Reboundary}
Literal equality transport of boundary isomorphisms, retaining the actual port cast rather than identifying differently ordered interfaces.

Interface equality changes only dependent port indices.
\lean{Copying/BoundaryReboundary}{11}{14}{BoundaryIso.interfaceEquality}
\par\penalty0

The interface Equality port assertion is the compatibility law for the indicated component of this construction. It relates the actual finite coordinates or certificates, with the hypotheses stated in the linked declaration.
For the proof, constructor and boundary cases are separated explicitly; the defined maps and incidence equations are substituted.
\lean{Copying/BoundaryReboundary}{16}{19}{BoundaryIso.interfaceEquality_port}
\par\penalty0

Simultaneous source/target interface equalities preserve the explicit boundary action by the corresponding input/output casts.
\lean{Copying/BoundaryReboundary}{23}{26}{BoundaryIso.reboundaryBoth}
\par\penalty0

\subsection{Boundary Source Reboundary}
Exact source interface casts for family boundary isomorphisms.

The interface Equality inverse assertion is the compatibility law for the indicated component of this construction. It relates the actual finite coordinates or certificates, with the hypotheses stated in the linked declaration.
For the proof, constructor and boundary cases are separated explicitly; the defined maps and incidence equations are substituted.
\lean{Copying/BoundarySourceReboundary}{9}{15}{BoundaryIso.interfaceEquality_inverse}
\par\penalty0

Change only the source interface presentation; the full original target family and its ordered coordinates remain literal.
\lean{Copying/BoundarySourceReboundary}{19}{21}{BoundaryIso.reboundarySource}
\par\penalty0

The reboundary Source inputs assertion computes the indicated ordered boundary coordinate in the assembled graph.
For the proof, the defined maps and incidence equations are substituted.
\lean{Copying/BoundarySourceReboundary}{23}{32}{BoundaryIso.reboundarySource_inputs}
\par\penalty0

The reboundary Source outputs assertion computes the indicated ordered boundary coordinate in the assembled graph.
For the proof, the defined maps and incidence equations are substituted.
\lean{Copying/BoundarySourceReboundary}{34}{43}{BoundaryIso.reboundarySource_outputs}
\par\penalty0

\subsection{Boundary Synchrony}
Synchronous boundary actions of private-output and closed families. The left outputs track precisely the common external input permutation.

The trans inputs assertion computes the indicated ordered boundary coordinate in the assembled graph.
For the proof, the defined maps and incidence equations are substituted.
\lean{Copying/BoundarySynchrony}{9}{17}{BoundaryIso.trans_inputs}
\par\penalty0

The trans outputs assertion computes the indicated ordered boundary coordinate in the assembled graph.
For the proof, the defined maps and incidence equations are substituted.
\lean{Copying/BoundarySynchrony}{19}{27}{BoundaryIso.trans_outputs}
\par\penalty0

Shared boundary action, including the left private-output attachment.
\lean{Copying/BoundarySynchrony}{30}{33}{BoundaryIso.Synchronous}
\par\penalty0

The synchronous refl assertion is the compatibility law for the indicated component of this construction. It relates the actual finite coordinates or certificates, with the hypotheses stated in the linked declaration.
For the proof, the defined maps and incidence equations are substituted.
\lean{Copying/BoundarySynchrony}{35}{36}{BoundaryIso.synchronous_refl}
\par\penalty0

Composing synchronized family changes retains both required boundary agreements; no graph equality or semantic closure is assumed.
For the proof, the defined maps and incidence equations are substituted.
\lean{Copying/BoundarySynchrony}{40}{46}{BoundaryIso.Synchronous.trans}
\par\penalty0

Independent synchronized families can be assembled in the literal tensor, preserving the same complete input/private-output order.
For the proof, constructor and boundary cases are separated explicitly; the defined maps and incidence equations are substituted.
\lean{Copying/BoundarySynchrony}{50}{93}{BoundaryIso.Synchronous.tensor}
\par\penalty0

\subsection{Boundary Tensor Ports}
Tensor boundary actions retain the actual port maps of each complete component, including internal ports and unequal interface presentations.

The tensor left port assertion is the compatibility law for the indicated component of this construction. It relates the actual finite coordinates or certificates, with the hypotheses stated in the linked declaration.
For the proof, the preceding component lemmas are composed or specialized at the stated parameters.
\lean{Copying/BoundaryTensorPorts}{11}{16}{BoundaryIso.tensor_left_port}
\par\penalty0

The tensor right port assertion is the compatibility law for the indicated component of this construction. It relates the actual finite coordinates or certificates, with the hypotheses stated in the linked declaration.
For the proof, the preceding component lemmas are composed or specialized at the stated parameters.
\lean{Copying/BoundaryTensorPorts}{18}{23}{BoundaryIso.tensor_right_port}
\par\penalty0

\subsection{Expansion Target Origin}
The target-side counterpart of actual expansion source ownership.

An expanded edge targeting an original context vertex must target one of its component's outputs; private internal vertices cannot escape.
For the proof, constructor and boundary cases are separated explicitly.
\lean{Copying/ExpansionTargetOrigin}{9}{16}{Search.Expansion.target_output}
\par\penalty0

Every source of a retained complete corolla is an original context vertex, even when the retained corolla labels a recursively filled box.
For the proof, constructor and boundary cases are separated explicitly; the defined maps and incidence equations are substituted.
\lean{Copying/ExpansionTargetOrigin}{20}{29}{Search.Expansion.Retained.source_origin}
\par\penalty0

\subsection{Lambda Entry Allocation}
Complete allocation overlap for lambda-entry templates. A stable lambda may belong to a multi-input shallow filling: only the retained ref edge is rewritten, while that whole original filling remains in the residual.

A non-hole primitive in the original context has literal corolla copies on both sides of an original recursive filling certificate.
For the proof, constructor and boundary cases are separated explicitly; the defined maps and incidence equations are substituted.
\lean{Copying/LambdaEntryAllocation}{14}{21}{DeepFillingEdge.primitive_copies}
\par\penalty0

An unchanged original ref head gives the complete one-step/one-step matching diagram. Stability on each side is proved from its actual rooted search history; the original multi-edge arguments need not be retained.
For the proof, constructor and boundary cases are separated explicitly; the defined maps and incidence equations are substituted.
\lean{Copying/LambdaEntryAllocation}{33}{54}{LambdaEntryOverlap.retained_allocation_matching}
\par\penalty0

The actual ref edge at an original host token is retained, although its stable lambda operand may lie inside any shallow lambda-entry family.
For the proof, constructor and boundary cases are separated explicitly; the defined maps and incidence equations are substituted.
\lean{Copying/LambdaEntryAllocation}{59}{69}{LambdaEntryOverlap.allocation_head}
\par\penalty0

Every actual allocation overlap has an opposite original step and a genuine original residual, including shallow lambda operands and deep boxes.
For the proof, the defined maps and incidence equations are substituted.
\lean{Copying/LambdaEntryAllocation}{73}{90}{LambdaEntryOverlap.allocation_matching}
\par\penalty0

The exact robust case quantifies over every independently supplied successor of the original machine.
For the proof, the preceding component lemmas are composed or specialized at the stated parameters.
\lean{Copying/LambdaEntryAllocation}{94}{110}{LambdaEntryOverlap.allocation_robust_case}
\par\penalty0

\subsection{Lambda Entry Numeric Overlap}
Arithmetic footprints cannot cross a lambda-family entry. These actual matching diagrams apply to both directions of Figure 20 parametricity.

Every edge of a sourced non-lambda ancestry avoids shallow family holes. The proof follows real component incidences, without rootedness.
For the proof, induction follows the original derivation or finite structure.
\lean{Copying/LambdaEntryNumericOverlap}{17}{30}{LambdaEntryOverlap.ancestor_not_hole}
\par\penalty0

The complete actual unary arithmetic match is retained.
For the proof, constructor and boundary cases are separated explicitly; the defined maps and incidence equations are substituted.
\lean{Copying/LambdaEntryNumericOverlap}{33}{49}{LambdaEntryOverlap.unary_avoids}
\par\penalty0

Both ordered numeric operands remain in the unchanged context.
For the proof, constructor and boundary cases are separated explicitly; the defined maps and incidence equations are substituted.
\lean{Copying/LambdaEntryNumericOverlap}{52}{71}{LambdaEntryOverlap.binary_avoids}
\par\penalty0

\subsection{Lambda Entry Numeric Residuals}
Actual quantitative arithmetic branches for templates with lambda-value entries, including both directions of the original parametricity template.

Every actual unary arithmetic successor has an opposite original compute step and a genuine residual specimen, with one transition per side.
For the proof, the defined maps and incidence equations are substituted.
\lean{Copying/LambdaEntryNumericResiduals}{16}{30}{LambdaEntryOverlap.unary_robust_case}
\par\penalty0

Both binary arithmetic operations preserve a genuine original specimen and its exact one-step comparison, for every actual represented successor.
For the proof, the defined maps and incidence equations are substituted.
\lean{Copying/LambdaEntryNumericResiduals}{34}{48}{LambdaEntryOverlap.binary_robust_case}
\par\penalty0

\subsection{Lambda Entry Search}
Original search across complete lambda families preserves every common host coordinate and every transition count. No robustness theorem is used.

A search step at an original host coordinate is matched in the other filled graph. Entering a family returns at the same lambda input; exiting one is impossible because its actual output interface is empty.
For the proof, constructor and boundary cases are separated explicitly.
\lean{Copying/LambdaEntrySearch}{18}{38}{Search.Expansion.lambda_step}
\par\penalty0

Every finite original search run remains at host coordinates and has an equally long actual search run through the opposite filled family.
For the proof, induction follows the original derivation or finite structure; the defined maps and incidence equations are substituted.
\lean{Copying/LambdaEntrySearch}{42}{54}{Search.Expansion.lambda_run}
\par\penalty0

Original search rootedness is equivalent at every common host position, including internal tokens and independently represented recursive families.
For the proof, the preceding component lemmas are composed or specialized at the stated parameters.
\lean{Copying/LambdaEntrySearch}{58}{78}{Search.Expansion.lambda_rooted_iff}
\par\penalty0

The complete witness transports rootedness by actual search alone; this lemma does not depend on parametricity robustness or equivalence.
For the proof, the preceding component lemmas are composed or specialized at the stated parameters.
\lean{Copying/LambdaEntrySearch}{87}{96}{FillingWitness.lambda_rooted_iff}
\par\penalty0

\subsection{No Output Overlap}
A closed-output template component cannot leak into the unchanged context. This is a structural fact about genuine recursive fillings.

Every pair in the pre-template has an empty ordered output interface.
\lean{Copying/NoOutputOverlap}{11}{12}{PreTemplate.NoOutputs}
\par\penalty0

The converse assertion is the compatibility law for the indicated component of this construction. It relates the actual finite coordinates or certificates, with the hypotheses stated in the linked declaration.
For the proof, the preceding component lemmas are composed or specialized at the stated parameters.
\lean{Copying/NoOutputOverlap}{14}{15}{PreTemplate.NoOutputs.converse}
\par\penalty0

A real target incidence at an original vertex forces a retained edge; it cannot be supplied by a component with no outputs.
For the proof, constructor and boundary cases are separated explicitly.
\lean{Copying/NoOutputOverlap}{24}{30}{NoOutputOverlap.target_retained}
\par\penalty0

Reflecting an adjacency backwards uses the actual source coordinates of the retained successor and retains the whole predecessor.
For the proof, the preceding component lemmas are composed or specialized at the stated parameters.
\lean{Copying/NoOutputOverlap}{34}{41}{NoOutputOverlap.predecessor_retained}
\par\penalty0

Every predecessor along a finite path ending at an unchanged context edge is unchanged too. No rootedness or acyclicity assumption is used.
For the proof, induction follows the original derivation or finite structure.
\lean{Copying/NoOutputOverlap}{45}{53}{NoOutputOverlap.ancestor_retained}
\par\penalty0

Hole freedom also excludes an unfilled original hole at a target incidence. This is the direct no-leakage condition used by overlap proofs.
For the proof, the preceding component lemmas are composed or specialized at the stated parameters.
\lean{Copying/NoOutputOverlap}{57}{64}{NoOutputOverlap.target_not_hole}
\par\penalty0

\subsection{Occurrence Boundary Cover Iso}
A complete edge-covering occurrence is a boundary isomorphism whenever its input and output attachments remain external. Boundary order is recorded by the actual occurrence, never silently identified.

A boundary-preserving complete occurrence has no input/output aliasing; its incidence laws already give injectivity on each remaining polarity.
For the proof, constructor and boundary cases are separated explicitly; the defined maps and incidence equations are substituted.
\lean{Copying/OccurrenceBoundaryCoverIso}{19}{44}{Opening.Occurrence.boundary_ports_injective}
\par\penalty0

Edge coverage and actual incidence coverage give port surjectivity. This assertion alone does not forbid input/output aliasing at internal wires.
For the proof, constructor and boundary cases are separated explicitly; the defined maps and incidence equations are substituted.
\lean{Copying/OccurrenceBoundaryCoverIso}{50}{62}{Opening.Occurrence.ports_surjective_of_edge_surjective}
\par\penalty0

The actual boundary action and full recursive labels determine a genuine isomorphism, with no separate carrier-count or vertex-surjectivity premise.
\lean{Copying/OccurrenceBoundaryCoverIso}{73}{96}{Opening.Occurrence.boundaryIso_of_edge_surjective}
\par\penalty0

\subsection{Occurrence Cut Boundary}
Complete disjoint occurrences retain their original external input and output attachments under an actual computed cut replacement.

The across Cut input assertion computes the indicated ordered boundary coordinate in the assembled graph.
For the proof, the defined maps and incidence equations are substituted.
\lean{Copying/OccurrenceCutBoundary}{12}{19}{Opening.Occurrence.acrossCut_input}
\par\penalty0

The across Cut output assertion computes the indicated ordered boundary coordinate in the assembled graph.
For the proof, the defined maps and incidence equations are substituted.
\lean{Copying/OccurrenceCutBoundary}{21}{28}{Opening.Occurrence.acrossCut_output}
\par\penalty0

\subsection{Occurrence Tensor Boundary}
Actual tensor assembly preserves external boundary roles componentwise. These facts retain the boundary permutation of the chosen occurrences.

Ports in the actual first tensor component retain their supplied map.
For the proof, constructor and boundary cases are separated explicitly; the defined maps and incidence equations are substituted.
\lean{Copying/OccurrenceTensorBoundary}{13}{18}{Opening.Occurrence.tensorJoin_left}
\par\penalty0

Ports in the actual second tensor component retain their supplied map.
For the proof, constructor and boundary cases are separated explicitly; the defined maps and incidence equations are substituted.
\lean{Copying/OccurrenceTensorBoundary}{21}{26}{Opening.Occurrence.tensorJoin_right}
\par\penalty0

An assembled external input remains external when both complete components have external input attachments.
For the proof, constructor and boundary cases are separated explicitly; the defined maps and incidence equations are substituted.
\lean{Copying/OccurrenceTensorBoundary}{30}{45}{Opening.Occurrence.tensorJoin_inputs}
\par\penalty0

The identical incidence statement for ordered external outputs.
For the proof, constructor and boundary cases are separated explicitly; the defined maps and incidence equations are substituted.
\lean{Copying/OccurrenceTensorBoundary}{48}{63}{Opening.Occurrence.tensorJoin_outputs}
\par\penalty0

\subsection{Parametricity Apply First}
Apply the selected first member of an arbitrary original finite family to an arbitrary value graph. All remaining family inputs and all argument resource outputs are preserved in their ordered original interfaces.

This defines the Vertex data for the construction just described. Boundary positions are kept distinct from newly inserted internal vertices.
\lean{Copying/ParametricityApplyFirst}{10}{10}{Parametricity.ApplyFirst.Vertex}
\par\penalty0

This defines the active data for the construction just described. Boundary positions are kept distinct from newly inserted internal vertices.
\lean{Copying/ParametricityApplyFirst}{12}{12}{Parametricity.ApplyFirst.active}
\par\penalty0

This defines the remaining data for the construction just described. Boundary positions are kept distinct from newly inserted internal vertices.
\lean{Copying/ParametricityApplyFirst}{13}{13}{Parametricity.ApplyFirst.remaining}
\par\penalty0

This defines the label data for the construction just described. Boundary positions are kept distinct from newly inserted internal vertices.
\lean{Copying/ParametricityApplyFirst}{14}{17}{Parametricity.ApplyFirst.label}
\par\penalty0

This defines the sources data for the construction just described. Boundary positions are kept distinct from newly inserted internal vertices.
\lean{Copying/ParametricityApplyFirst}{18}{21}{Parametricity.ApplyFirst.sources}
\par\penalty0

This defines the targets data for the construction just described. Boundary positions are kept distinct from newly inserted internal vertices.
\lean{Copying/ParametricityApplyFirst}{22}{25}{Parametricity.ApplyFirst.targets}
\par\penalty0

The remaining ne assertion is the compatibility law for the indicated component of this construction. It relates the actual finite coordinates or certificates, with the hypotheses stated in the linked declaration.
For the proof, the remaining natural-number index bounds follow by linear arithmetic; the defined maps and incidence equations are substituted.
\lean{Copying/ParametricityApplyFirst}{27}{31}{Parametricity.ApplyFirst.remaining_ne}
\par\penalty0

This defines the common graph frame for the construction just described. Boundary positions are kept distinct from newly inserted internal vertices.
\lean{Copying/ParametricityApplyFirst}{33}{157}{Parametricity.ApplyFirst.frame}
\par\penalty0

This defines the pattern data for the construction just described. Boundary positions are kept distinct from newly inserted internal vertices.
\lean{Copying/ParametricityApplyFirst}{159}{162}{Parametricity.ApplyFirst.pattern}
\par\penalty0

This defines the filling data for the construction just described. Boundary positions are kept distinct from newly inserted internal vertices.
\lean{Copying/ParametricityApplyFirst}{164}{172}{Parametricity.ApplyFirst.filling}
\par\penalty0

This defines the finite typed net for the construction just described. Boundary positions are kept distinct from newly inserted internal vertices.
\lean{Copying/ParametricityApplyFirst}{174}{178}{Parametricity.ApplyFirst.net}
\par\penalty0

The head map is a complete occurrence, preserving original labels and every ordered source and target attachment.
\lean{Copying/ParametricityApplyFirst}{179}{180}{Parametricity.ApplyFirst.head}
\par\penalty0

The family map is a complete occurrence, preserving original labels and every ordered source and target attachment.
\lean{Copying/ParametricityApplyFirst}{181}{182}{Parametricity.ApplyFirst.family}
\par\penalty0

The argument map is a complete occurrence, preserving original labels and every ordered source and target attachment.
\lean{Copying/ParametricityApplyFirst}{183}{184}{Parametricity.ApplyFirst.argument}
\par\penalty0

The head input assertion computes the indicated ordered boundary coordinate in the assembled graph.
For the proof, constructor and boundary cases are separated explicitly; the defined maps and incidence equations are substituted.
\lean{Copying/ParametricityApplyFirst}{186}{186}{Parametricity.ApplyFirst.head_input}
\par\penalty0

The head function assertion is the compatibility law for the indicated component of this construction. It relates the actual finite coordinates or certificates, with the hypotheses stated in the linked declaration.
For the proof, constructor and boundary cases are separated explicitly; the defined maps and incidence equations are substituted.
\lean{Copying/ParametricityApplyFirst}{188}{189}{Parametricity.ApplyFirst.head_function}
\par\penalty0

The head argument assertion is the compatibility law for the indicated component of this construction. It relates the actual finite coordinates or certificates, with the hypotheses stated in the linked declaration.
For the proof, constructor and boundary cases are separated explicitly; the defined maps and incidence equations are substituted.
\lean{Copying/ParametricityApplyFirst}{191}{192}{Parametricity.ApplyFirst.head_argument}
\par\penalty0

The family succ assertion is the compatibility law for the indicated component of this construction. It relates the actual finite coordinates or certificates, with the hypotheses stated in the linked declaration.
For the proof, constructor and boundary cases are separated explicitly; the defined maps and incidence equations are substituted.
\lean{Copying/ParametricityApplyFirst}{194}{200}{Parametricity.ApplyFirst.family_succ}
\par\penalty0

The argument output assertion computes the indicated ordered boundary coordinate in the assembled graph.
For the proof, constructor and boundary cases are separated explicitly; the defined maps and incidence equations are substituted.
\lean{Copying/ParametricityApplyFirst}{202}{206}{Parametricity.ApplyFirst.argument_output}
\par\penalty0

The witness data give an actual recursive paired filling, with its common context and endpoint representation maps.
\lean{Copying/ParametricityApplyFirst}{208}{216}{Parametricity.ApplyFirst.witness}
\par\penalty0

\subsection{Parametricity Apply Join}
Complete actual app/family/argument occurrences assemble the original first-input application context, including all remaining external inputs.

Original frame attachments, with the selected application at input0.
\lean{Copying/ParametricityApplyJoin}{15}{18}{Parametricity.ApplyFirst.Join.host}
\par\penalty0

The components map is a complete occurrence, preserving original labels and every ordered source and target attachment.
\lean{Copying/ParametricityApplyJoin}{20}{28}{Parametricity.ApplyFirst.Join.components}
\par\penalty0

The inputs assertion computes the indicated ordered boundary coordinate in the assembled graph.
For the proof, constructor and boundary cases are separated explicitly; the remaining natural-number index bounds follow by linear arithmetic; the defined maps and incidence equations are substituted.
\lean{Copying/ParametricityApplyJoin}{31}{65}{Parametricity.ApplyFirst.Join.inputs}
\par\penalty0

The outputs assertion computes the indicated ordered boundary coordinate in the assembled graph.
For the proof, constructor and boundary cases are separated explicitly; the remaining natural-number index bounds follow by linear arithmetic; the defined maps and incidence equations are substituted.
\lean{Copying/ParametricityApplyJoin}{68}{91}{Parametricity.ApplyFirst.Join.outputs}
\par\penalty0

The disjoint assertion excludes a shared edge or coordinate between the indicated independently inserted components.
For the proof, constructor and boundary cases are separated explicitly; the defined maps and incidence equations are substituted.
\lean{Copying/ParametricityApplyJoin}{95}{116}{Parametricity.ApplyFirst.Join.disjoint}
\par\penalty0

Every original ordered incidence is realized by the supplied complete occurrences. Their proven disjointness supplies all necessary freshness.
\lean{Copying/ParametricityApplyJoin}{120}{123}{Parametricity.ApplyFirst.Join.occurrence}
\par\penalty0

The input zero assertion computes the indicated ordered boundary coordinate in the assembled graph.
Both sides are definitionally equal.
\lean{Copying/ParametricityApplyJoin}{125}{126}{Parametricity.ApplyFirst.Join.input_zero}
\par\penalty0

The input succ assertion computes the indicated ordered boundary coordinate in the assembled graph.
For the proof, the remaining natural-number index bounds follow by linear arithmetic; the defined maps and incidence equations are substituted.
\lean{Copying/ParametricityApplyJoin}{128}{132}{Parametricity.ApplyFirst.Join.input_succ}
\par\penalty0

The output assertion computes the indicated ordered boundary coordinate in the assembled graph.
Both sides are definitionally equal.
\lean{Copying/ParametricityApplyJoin}{134}{135}{Parametricity.ApplyFirst.Join.output}
\par\penalty0

\subsection{Parametricity Apply Occurrence}
A selected member of a complete family gives an actual complete application occurrence; all remaining family members and private wires stay in its constructed complement.

The selected map is a complete occurrence, preserving original labels and every ordered source and target attachment.
\lean{Copying/ParametricityApplyOccurrence}{15}{15}{Parametricity.ApplyFirst.selected}
\par\penalty0

The selected argument disjoint assertion excludes a shared edge or coordinate between the indicated independently inserted components.
For the proof, constructor and boundary cases are separated explicitly.
\lean{Copying/ParametricityApplyOccurrence}{17}{18}{Parametricity.ApplyFirst.selected_argument_disjoint}
\par\penalty0

The binary Arguments map is a complete occurrence, preserving original labels and every ordered source and target attachment.
\lean{Copying/ParametricityApplyOccurrence}{20}{21}{Parametricity.ApplyFirst.binaryArguments}
\par\penalty0

The binary link assertion is the compatibility law for the indicated component of this construction. It relates the actual finite coordinates or certificates, with the hypotheses stated in the linked declaration.
For the proof, constructor and boundary cases are separated explicitly; the remaining natural-number index bounds follow by linear arithmetic; the defined maps and incidence equations are substituted.
\lean{Copying/ParametricityApplyOccurrence}{24}{39}{Parametricity.ApplyFirst.binary_link}
\par\penalty0

The head binary disjoint assertion excludes a shared edge or coordinate between the indicated independently inserted components.
For the proof, constructor and boundary cases are separated explicitly.
\lean{Copying/ParametricityApplyOccurrence}{41}{43}{Parametricity.ApplyFirst.head_binary_disjoint}
\par\penalty0

The unchanged original app, selected family member and arbitrary argument form a complete ordered local source, with no boundary or freshness premise.
\lean{Copying/ParametricityApplyOccurrence}{47}{49}{Parametricity.ApplyFirst.binaryOccurrence}
\par\penalty0

The binary input assertion computes the indicated ordered boundary coordinate in the assembled graph.
For the proof, the defined maps and incidence equations are substituted.
\lean{Copying/ParametricityApplyOccurrence}{51}{53}{Parametricity.ApplyFirst.binary_input}
\par\penalty0

\subsection{Parametricity Boundary Barrier}
Exact entry and no-exit boundaries of the private Figure 20 families, in both directions and in arbitrary original recursive filling contexts.

Every component input is sourced by a lambda value at the surface.
\lean{Copying/ParametricityBoundaryBarrier}{11}{12}{PreTemplate.LambdaEntries}
\par\penalty0

A non-lambda edge sourcing a host vertex cannot belong to a fresh shallow family. This retains the exact original incidence and label.
For the proof, constructor and boundary cases are separated explicitly; the defined maps and incidence equations are substituted.
\lean{Copying/ParametricityBoundaryBarrier}{21}{32}{LambdaEntryOverlap.source_retained}
\par\penalty0

Any sourced edge of a shallow family is its boundary lambda edge. The hypothesis records the actual original hole, not a chosen presentation.
For the proof, the defined maps and incidence equations are substituted.
\lean{Copying/ParametricityBoundaryBarrier}{36}{46}{LambdaEntryOverlap.shallow_source_label}
\par\penalty0

Hole freedom excludes retained unfilled holes at non-lambda sources.
For the proof, the preceding component lemmas are composed or specialized at the stated parameters.
\lean{Copying/ParametricityBoundaryBarrier}{49}{55}{LambdaEntryOverlap.source_not_hole}
\par\penalty0

A whole mapped operational ancestry with no lambda surface remains in the original context, while arbitrary deep box bodies stay unrestricted.
For the proof, induction follows the original derivation or finite structure.
\lean{Copying/ParametricityBoundaryBarrier}{59}{76}{LambdaEntryOverlap.ancestor_retained}
\par\penalty0

Both independent representations preserve the actual lambda entry.
For the proof, constructor and boundary cases are separated explicitly.
\lean{Copying/ParametricityBoundaryBarrier}{83}{95}{Parametricity.Pair.input_heads}
\par\penalty0

The template lambda Entries assertion is the compatibility law for the indicated component of this construction. It relates the actual finite coordinates or certificates, with the hypotheses stated in the linked declaration.
For the proof, the preceding component lemmas are composed or specialized at the stated parameters.
\lean{Copying/ParametricityBoundaryBarrier}{97}{97}{Parametricity.template_lambdaEntries}
\par\penalty0

The converse lambda Entries assertion is the compatibility law for the indicated component of this construction. It relates the actual finite coordinates or certificates, with the hypotheses stated in the linked declaration.
For the proof, the preceding component lemmas are composed or specialized at the stated parameters.
\lean{Copying/ParametricityBoundaryBarrier}{98}{98}{Parametricity.converse_lambdaEntries}
\par\penalty0

The complete private cell and all sharing resources have no escaping wire, regardless of family multiplicities or recursive graph carriers.
For the proof, the preceding component lemmas are composed or specialized at the stated parameters.
\lean{Copying/ParametricityBoundaryBarrier}{102}{102}{Parametricity.template_noOutputs}
\par\penalty0

The converse no Outputs assertion computes the indicated ordered boundary coordinate in the assembled graph.
For the proof, the preceding component lemmas are composed or specialized at the stated parameters.
\lean{Copying/ParametricityBoundaryBarrier}{103}{103}{Parametricity.converse_noOutputs}
\par\penalty0

\subsection{Parametricity Boundary Witness}
Reordering complete Figure 20 families is justified by an actual common input-permutation context. The private contraction tree is reindexed with the recorded output action; no closure axiom is added to parametricity.

Every bijection of a homogeneous term interface is a typed permutation.
\lean{Copying/ParametricityBoundaryWitness}{13}{28}{Parametricity.BoundaryWitness.termOrder}
\par\penalty0

Reindex every input of the original tree with the exact private-output action of the family isomorphism, retaining all primitive tree edges.
\lean{Copying/ParametricityBoundaryWitness}{32}{32}{Parametricity.BoundaryWitness.tree}
\par\penalty0

The tree boundary action agrees with every shared private wire.
\lean{Copying/ParametricityBoundaryWitness}{35}{37}{Parametricity.BoundaryWitness.treeIso}
\par\penalty0

The compatible assertion is the compatibility law for the indicated component of this construction. It relates the actual finite coordinates or certificates, with the hypotheses stated in the linked declaration.
For the proof, the defined maps and incidence equations are substituted.
\lean{Copying/ParametricityBoundaryWitness}{39}{41}{Parametricity.BoundaryWitness.compatible}
\par\penalty0

Whole private-cell assembly carries precisely the external input action of the complete family, and fixes its empty output interface.
\lean{Copying/ParametricityBoundaryWitness}{45}{49}{Parametricity.BoundaryWitness.leftIso}
\par\penalty0

The left Iso input assertion computes the indicated ordered boundary coordinate in the assembled graph.
Both sides are definitionally equal.
\lean{Copying/ParametricityBoundaryWitness}{51}{52}{Parametricity.BoundaryWitness.leftIso_input}
\par\penalty0

The original literal parametricity pair, placed in a genuine common permutation context, certifies arbitrary independently represented families whose external input actions agree.
\lean{Copying/ParametricityBoundaryWitness}{57}{66}{Parametricity.BoundaryWitness.filling}
\par\penalty0

A syntactically different presentation of the same homogeneous input list needs only its proved interface equality, not a new semantic premise.
\lean{Copying/ParametricityBoundaryWitness}{70}{81}{Parametricity.BoundaryWitness.filling_reboundary}
\par\penalty0

Natural-number reassociation of the tree arity transports the complete original assembly without changing its graph or its boundary coordinates.
\lean{Copying/ParametricityBoundaryWitness}{85}{96}{Parametricity.BoundaryWitness.filling_reboundary_tree}
\par\penalty0

\subsection{Parametricity Copy Assembly}
The exact residual-input context commutes with arbitrary serial continuations. The explicit input action is preserved, not quotiented away.

The serial map is an ordered recursive graph isomorphism between the stated original constructions.
\lean{Copying/ParametricityCopyAssembly}{12}{32}{CopyingInputPermutation.serial}
\par\penalty0

The congr map is an ordered recursive graph isomorphism between the stated original constructions.
\lean{Copying/ParametricityCopyAssembly}{41}{44}{Parametricity.CopyContext.congr}
\par\penalty0

The serial map is an ordered recursive graph isomorphism between the stated original constructions.
\lean{Copying/ParametricityCopyAssembly}{46}{50}{Parametricity.CopyContext.serial}
\par\penalty0

\subsection{Parametricity Copy Components}
The actual active and retained abstraction occurrences after one original copy, with an arbitrary retained input graph (including zero inputs).

This defines the ordered target list for the construction just described. Boundary positions are kept distinct from newly inserted internal vertices.
\lean{Copying/ParametricityCopyComponents}{15}{16}{Parametricity.CopyComponents.target}
\par\penalty0

The branch map is a complete occurrence, preserving original labels and every ordered source and target attachment.
\lean{Copying/ParametricityCopyComponents}{18}{20}{Parametricity.CopyComponents.branch}
\par\penalty0

The active map is a complete occurrence, preserving original labels and every ordered source and target attachment.
\lean{Copying/ParametricityCopyComponents}{21}{21}{Parametricity.CopyComponents.active}
\par\penalty0

The retained map is a complete occurrence, preserving original labels and every ordered source and target attachment.
\lean{Copying/ParametricityCopyComponents}{22}{23}{Parametricity.CopyComponents.retained}
\par\penalty0

The upstream map is a complete occurrence, preserving original labels and every ordered source and target attachment.
\lean{Copying/ParametricityCopyComponents}{24}{25}{Parametricity.CopyComponents.upstream}
\par\penalty0

The distributor map is a complete occurrence, preserving original labels and every ordered source and target attachment.
\lean{Copying/ParametricityCopyComponents}{26}{27}{Parametricity.CopyComponents.distributor}
\par\penalty0

The payload map is a complete occurrence, preserving original labels and every ordered source and target attachment.
\lean{Copying/ParametricityCopyComponents}{29}{32}{Parametricity.CopyComponents.payload}
\par\penalty0

The payload disjoint assertion excludes a shared edge or coordinate between the indicated independently inserted components.
For the proof, constructor and boundary cases are separated explicitly; the defined maps and incidence equations are substituted.
\lean{Copying/ParametricityCopyComponents}{34}{38}{Parametricity.CopyComponents.payload_disjoint}
\par\penalty0

The complete two-payload graph occurs before the original complete resource distributor. Neither copy's recursively boxed body is modified.
\lean{Copying/ParametricityCopyComponents}{42}{45}{Parametricity.CopyComponents.pair}
\par\penalty0

The pair component assertion is the compatibility law for the indicated component of this construction. It relates the actual finite coordinates or certificates, with the hypotheses stated in the linked declaration.
For the proof, constructor and boundary cases are separated explicitly; the defined maps and incidence equations are substituted.
\lean{Copying/ParametricityCopyComponents}{47}{67}{Parametricity.CopyComponents.pair_component}
\par\penalty0

The pair output assertion computes the indicated ordered boundary coordinate in the assembled graph.
For the proof, constructor and boundary cases are separated explicitly; the defined maps and incidence equations are substituted.
\lean{Copying/ParametricityCopyComponents}{85}{93}{Parametricity.CopyComponents.pair_output}
\par\penalty0

The pair distributor disjoint assertion excludes a shared edge or coordinate between the indicated independently inserted components.
For the proof, constructor and boundary cases are separated explicitly.
\lean{Copying/ParametricityCopyComponents}{95}{96}{Parametricity.CopyComponents.pair_distributor_disjoint}
\par\penalty0

The literal merged two-copy graph, with every canonical output row, occurs inside the original target; only its retained input tree lies outside.
\lean{Copying/ParametricityCopyComponents}{100}{102}{Parametricity.CopyComponents.duplicate}
\par\penalty0

The duplicate input assertion computes the indicated ordered boundary coordinate in the assembled graph.
For the proof, constructor and boundary cases are separated explicitly; the defined maps and incidence equations are substituted.
\lean{Copying/ParametricityCopyComponents}{104}{111}{Parametricity.CopyComponents.duplicate_input}
\par\penalty0

The upstream link assertion is the compatibility law for the indicated component of this construction. It relates the actual finite coordinates or certificates, with the hypotheses stated in the linked declaration.
For the proof, the defined maps and incidence equations are substituted.
\lean{Copying/ParametricityCopyComponents}{113}{119}{Parametricity.CopyComponents.upstream_link}
\par\penalty0

The upstream duplicate disjoint assertion excludes a shared edge or coordinate between the indicated independently inserted components.
For the proof, constructor and boundary cases are separated explicitly.
\lean{Copying/ParametricityCopyComponents}{121}{127}{Parametricity.CopyComponents.upstream_duplicate_disjoint}
\par\penalty0

The upstream input assertion computes the indicated ordered boundary coordinate in the assembled graph.
For the proof, the remaining natural-number index bounds follow by linear arithmetic; the defined maps and incidence equations are substituted.
\lean{Copying/ParametricityCopyComponents}{129}{135}{Parametricity.CopyComponents.upstream_input}
\par\penalty0

The active input assertion computes the indicated ordered boundary coordinate in the assembled graph.
For the proof, the defined maps and incidence equations are substituted.
\lean{Copying/ParametricityCopyComponents}{137}{140}{Parametricity.CopyComponents.active_input}
\par\penalty0

The duplicate output assertion computes the indicated ordered boundary coordinate in the assembled graph.
For the proof, the defined maps and incidence equations are substituted.
\lean{Copying/ParametricityCopyComponents}{142}{148}{Parametricity.CopyComponents.duplicate_output}
\par\penalty0

\subsection{Parametricity Copy Context}
The residual-input context for growing a Figure 20 abstraction family. The new active abstraction remains input zero. Every remaining upstream input keeps its original order, before the untouched family inputs.

This defines the forward Index data for the construction just described. Boundary positions are kept distinct from newly inserted internal vertices.
\lean{Copying/ParametricityCopyContext}{13}{15}{Parametricity.CopyContext.forwardIndex}
\par\penalty0

This defines the backward Index data for the construction just described. Boundary positions are kept distinct from newly inserted internal vertices.
\lean{Copying/ParametricityCopyContext}{16}{18}{Parametricity.CopyContext.backwardIndex}
\par\penalty0

This defines the rotation data for the construction just described. Boundary positions are kept distinct from newly inserted internal vertices.
\lean{Copying/ParametricityCopyContext}{20}{38}{Parametricity.CopyContext.rotation}
\par\penalty0

This defines the finite typed net for the construction just described. Boundary positions are kept distinct from newly inserted internal vertices.
\lean{Copying/ParametricityCopyContext}{40}{41}{Parametricity.CopyContext.net}
\par\penalty0

The upstream map is a complete occurrence, preserving original labels and every ordered source and target attachment.
\lean{Copying/ParametricityCopyContext}{43}{45}{Parametricity.CopyContext.upstream}
\par\penalty0

The family map is a complete occurrence, preserving original labels and every ordered source and target attachment.
\lean{Copying/ParametricityCopyContext}{46}{48}{Parametricity.CopyContext.family}
\par\penalty0

The upstream input assertion computes the indicated ordered boundary coordinate in the assembled graph.
For the proof, the remaining natural-number index bounds follow by linear arithmetic; the defined maps and incidence equations are substituted.
\lean{Copying/ParametricityCopyContext}{50}{60}{Parametricity.CopyContext.upstream_input}
\par\penalty0

The active input assertion computes the indicated ordered boundary coordinate in the assembled graph.
For the proof, the defined maps and incidence equations are substituted.
\lean{Copying/ParametricityCopyContext}{62}{69}{Parametricity.CopyContext.active_input}
\par\penalty0

The retained input assertion computes the indicated ordered boundary coordinate in the assembled graph.
For the proof, the defined maps and incidence equations are substituted.
\lean{Copying/ParametricityCopyContext}{71}{80}{Parametricity.CopyContext.retained_input}
\par\penalty0

The remaining input assertion computes the indicated ordered boundary coordinate in the assembled graph.
For the proof, the remaining natural-number index bounds follow by linear arithmetic; the defined maps and incidence equations are substituted.
\lean{Copying/ParametricityCopyContext}{82}{98}{Parametricity.CopyContext.remaining_input}
\par\penalty0

The family output assertion computes the indicated ordered boundary coordinate in the assembled graph.
For the proof, the defined maps and incidence equations are substituted.
\lean{Copying/ParametricityCopyContext}{100}{106}{Parametricity.CopyContext.family_output}
\par\penalty0

One original enlarged-family specimen is inserted once. The residual sharing tree is retained identically, without an assumed composition law.
\lean{Copying/ParametricityCopyContext}{110}{113}{Parametricity.CopyContext.witness}
\par\penalty0

\subsection{Parametricity Copy Distributors}
Identify the actual zero/one-private-resource copying distributions. Only empty representation carriers disappear: the nonempty private row remains the full original canonical D1,2, including its weakening.

The private Row map is an ordered recursive graph isomorphism between the stated original constructions.
\lean{Copying/ParametricityCopyDistributors}{11}{30}{Parametricity.CopyDistributors.privateRow}
\par\penalty0

The private Pair map is an ordered recursive graph isomorphism between the stated original constructions.
\lean{Copying/ParametricityCopyDistributors}{32}{47}{Parametricity.CopyDistributors.privatePair}
\par\penalty0

The private lambda's copy merge is exactly two original lambdas followed by the complete canonical two-input term tree.
\lean{Copying/ParametricityCopyDistributors}{51}{54}{Parametricity.CopyDistributors.privateMerge}
\par\penalty0

A closed lambda has no resource rows: the literal empty distributor contributes no edge and the two ordered payloads remain unchanged.
\lean{Copying/ParametricityCopyDistributors}{58}{60}{Parametricity.CopyDistributors.closedMerge}
\par\penalty0

\subsection{Parametricity Copy Execution}
One original lambda-copying step with an arbitrary retained family and arbitrary serial continuation. Exact target normalization exposes the enlarged family, retaining the full upstream residual tree in a common input context.

Lift an actual rewrite in the first tensor component without altering any edge or ordered port of the second component.
For the proof, the defined maps and incidence equations are substituted.
\lean{Copying/ParametricityCopyExecution}{15}{23}{Combined.Rewrite.tensor_left}
\par\penalty0

Serial continuation preserves an actual rewrite and its two input coordinates, including the resumed active family member.
For the proof, constructor and boundary cases are separated explicitly; the defined maps and incidence equations are substituted.
\lean{Copying/ParametricityCopyExecution}{27}{36}{Combined.Rewrite.serial_left}
\par\penalty0

This defines the ordered source list for the construction just described. Boundary positions are kept distinct from newly inserted internal vertices.
\lean{Copying/ParametricityCopyExecution}{46}{48}{Parametricity.CopyExecution.source}
\par\penalty0

This defines the typed family of inserted components for the construction just described. Boundary positions are kept distinct from newly inserted internal vertices.
\lean{Copying/ParametricityCopyExecution}{49}{50}{Parametricity.CopyExecution.family}
\par\penalty0

This defines the ordered target list for the construction just described. Boundary positions are kept distinct from newly inserted internal vertices.
\lean{Copying/ParametricityCopyExecution}{51}{51}{Parametricity.CopyExecution.target}
\par\penalty0

The whole original selected copying rule fires exactly once. Its target contains the literal enlarged two-payload family and every original other family member, with the original active input still zero.
For the proof, the remaining natural-number index bounds follow by linear arithmetic; the defined maps and incidence equations are substituted.
\lean{Copying/ParametricityCopyExecution}{56}{65}{Parametricity.CopyExecution.rewrite}
\par\penalty0

The same actual compute step retains an arbitrary private sharing/store continuation and exposes it inside the common residual-input context.
For the proof, the remaining natural-number index bounds follow by linear arithmetic; the defined maps and incidence equations are substituted.
\lean{Copying/ParametricityCopyExecution}{69}{75}{Parametricity.CopyExecution.continued}
\par\penalty0

\subsection{Parametricity Copy Extension}
Pull the residual copying-input context outside an arbitrary untouched family. The external order remains active, remaining upstream inputs, family.

This defines the ordered target list for the construction just described. Boundary positions are kept distinct from newly inserted internal vertices.
\lean{Copying/ParametricityCopyExtension}{16}{18}{Parametricity.CopyExtension.target}
\par\penalty0

This defines the whole data for the construction just described. Boundary positions are kept distinct from newly inserted internal vertices.
\lean{Copying/ParametricityCopyExtension}{19}{20}{Parametricity.CopyExtension.whole}
\par\penalty0

The left map is a complete occurrence, preserving original labels and every ordered source and target attachment.
\lean{Copying/ParametricityCopyExtension}{22}{23}{Parametricity.CopyExtension.left}
\par\penalty0

The right map is a complete occurrence, preserving original labels and every ordered source and target attachment.
\lean{Copying/ParametricityCopyExtension}{24}{25}{Parametricity.CopyExtension.right}
\par\penalty0

The upstream map is a complete occurrence, preserving original labels and every ordered source and target attachment.
\lean{Copying/ParametricityCopyExtension}{26}{26}{Parametricity.CopyExtension.upstream}
\par\penalty0

The selected map is a complete occurrence, preserving original labels and every ordered source and target attachment.
\lean{Copying/ParametricityCopyExtension}{27}{27}{Parametricity.CopyExtension.selected}
\par\penalty0

The selected disjoint assertion excludes a shared edge or coordinate between the indicated independently inserted components.
For the proof, constructor and boundary cases are separated explicitly.
\lean{Copying/ParametricityCopyExtension}{29}{30}{Parametricity.CopyExtension.selected_disjoint}
\par\penalty0

The family map is a complete occurrence, preserving original labels and every ordered source and target attachment.
\lean{Copying/ParametricityCopyExtension}{32}{34}{Parametricity.CopyExtension.family}
\par\penalty0

The family first assertion is the compatibility law for the indicated component of this construction. It relates the actual finite coordinates or certificates, with the hypotheses stated in the linked declaration.
For the proof, the defined maps and incidence equations are substituted.
\lean{Copying/ParametricityCopyExtension}{36}{40}{Parametricity.CopyExtension.family_first}
\par\penalty0

The family rest assertion is the compatibility law for the indicated component of this construction. It relates the actual finite coordinates or certificates, with the hypotheses stated in the linked declaration.
For the proof, the remaining natural-number index bounds follow by linear arithmetic; the defined maps and incidence equations are substituted.
\lean{Copying/ParametricityCopyExtension}{42}{47}{Parametricity.CopyExtension.family_rest}
\par\penalty0

The link assertion is the compatibility law for the indicated component of this construction. It relates the actual finite coordinates or certificates, with the hypotheses stated in the linked declaration.
For the proof, the defined maps and incidence equations are substituted.
\lean{Copying/ParametricityCopyExtension}{49}{52}{Parametricity.CopyExtension.link}
\par\penalty0

The disjoint assertion excludes a shared edge or coordinate between the indicated independently inserted components.
For the proof, constructor and boundary cases are separated explicitly.
\lean{Copying/ParametricityCopyExtension}{54}{60}{Parametricity.CopyExtension.disjoint}
\par\penalty0

The raw map is a complete occurrence, preserving original labels and every ordered source and target attachment.
\lean{Copying/ParametricityCopyExtension}{62}{65}{Parametricity.CopyExtension.raw}
\par\penalty0

The upstream input assertion computes the indicated ordered boundary coordinate in the assembled graph.
For the proof, the remaining natural-number index bounds follow by linear arithmetic; the defined maps and incidence equations are substituted.
\lean{Copying/ParametricityCopyExtension}{67}{72}{Parametricity.CopyExtension.upstream_input}
\par\penalty0

The family active assertion is the compatibility law for the indicated component of this construction. It relates the actual finite coordinates or certificates, with the hypotheses stated in the linked declaration.
For the proof, the defined maps and incidence equations are substituted.
\lean{Copying/ParametricityCopyExtension}{74}{78}{Parametricity.CopyExtension.family_active}
\par\penalty0

The family other assertion is the compatibility law for the indicated component of this construction. It relates the actual finite coordinates or certificates, with the hypotheses stated in the linked declaration.
For the proof, the remaining natural-number index bounds follow by linear arithmetic; the defined maps and incidence equations are substituted.
\lean{Copying/ParametricityCopyExtension}{80}{88}{Parametricity.CopyExtension.family_other}
\par\penalty0

The family output assertion computes the indicated ordered boundary coordinate in the assembled graph.
For the proof, constructor and boundary cases are separated explicitly; the defined maps and incidence equations are substituted.
\lean{Copying/ParametricityCopyExtension}{90}{106}{Parametricity.CopyExtension.family_output}
\par\penalty0

The raw input assertion computes the indicated ordered boundary coordinate in the assembled graph.
For the proof, constructor and boundary cases are separated explicitly; the remaining natural-number index bounds follow by linear arithmetic; the defined maps and incidence equations are substituted.
\lean{Copying/ParametricityCopyExtension}{108}{148}{Parametricity.CopyExtension.raw_input}
\par\penalty0

The raw output assertion computes the indicated ordered boundary coordinate in the assembled graph.
For the proof, the defined maps and incidence equations are substituted.
\lean{Copying/ParametricityCopyExtension}{150}{153}{Parametricity.CopyExtension.raw_output}
\par\penalty0

The occurrence map is a complete occurrence, preserving original labels and every ordered source and target attachment.
\lean{Copying/ParametricityCopyExtension}{155}{156}{Parametricity.CopyExtension.occurrence}
\par\penalty0

The input assertion computes the indicated ordered boundary coordinate in the assembled graph.
For the proof, the defined maps and incidence equations are substituted.
\lean{Copying/ParametricityCopyExtension}{158}{161}{Parametricity.CopyExtension.input}
\par\penalty0

The output assertion computes the indicated ordered boundary coordinate in the assembled graph.
For the proof, the preceding component lemmas are composed or specialized at the stated parameters.
\lean{Copying/ParametricityCopyExtension}{163}{164}{Parametricity.CopyExtension.output}
\par\penalty0

The cover assertion states that every target coordinate has an original preimage under the indicated map.
For the proof, constructor and boundary cases are separated explicitly; the defined maps and incidence equations are substituted.
\lean{Copying/ParametricityCopyExtension}{166}{173}{Parametricity.CopyExtension.cover}
\par\penalty0

Exact tensor extension of the residual-input context. No permutation of external inputs is hidden in this graph isomorphism.
\lean{Copying/ParametricityCopyExtension}{177}{178}{Parametricity.CopyExtension.iso}
\par\penalty0

\subsection{Parametricity Copy Families}
Genuine original Figure20 residual specimens after copying either kind of abstraction, for arbitrary finite large and small families.

Copying a large abstraction increases the large family by exactly one. The old private tree is extended by the entire canonical two-input term row.
\lean{Copying/ParametricityCopyFamilies}{27}{44}{Parametricity.CopyFamilies.large}
\par\penalty0

Copying a small abstraction increases the small family by exactly one. The two selected members move together past the large block; their private outputs use the same proved permutation as both external input interfaces.
\lean{Copying/ParametricityCopyFamilies}{49}{66}{Parametricity.CopyFamilies.small}
\par\penalty0

Full original copy/copy diagram for a selected large abstraction and arbitrarily many remaining large and small members. The only parameters are actual finite trees and family sizes; the residual is a constructed specimen.
For the proof, the remaining natural-number index bounds follow by linear arithmetic; the defined maps and incidence equations are substituted.
\lean{Copying/ParametricityCopyFamilies}{71}{85}{Parametricity.CopyFamilies.large_copy}
\par\penalty0

The corresponding complete original diagram for any selected small abstraction. The private-cell pointer is duplicated through the original canonical term row; the right closed abstraction has no private pointer.
For the proof, the remaining natural-number index bounds follow by linear arithmetic; the defined maps and incidence equations are substituted.
\lean{Copying/ParametricityCopyFamilies}{90}{104}{Parametricity.CopyFamilies.small_copy}
\par\penalty0

\subsection{Parametricity Copy Graft Execution}
Actual one-step copying through the complete original upstream graft at the selected family member. Only arithmetic equality of interface lengths is cast; every selected and untouched family input keeps its ordered position.

This defines the graft Source data for the construction just described. Boundary positions are kept distinct from newly inserted internal vertices.
\lean{Copying/ParametricityCopyGraftExecution}{15}{17}{Parametricity.CopyExecution.graftSource}
\par\penalty0

The source Iso map is an ordered recursive graph isomorphism between the stated original constructions.
\lean{Copying/ParametricityCopyGraftExecution}{19}{23}{Parametricity.CopyExecution.sourceIso}
\par\penalty0

The entire original graft takes one actual original copying transition, not a conditional step inferred from a future contextual-equivalence theorem.
For the proof, the preceding component lemmas are composed or specialized at the stated parameters.
\lean{Copying/ParametricityCopyGraftExecution}{27}{30}{Parametricity.CopyExecution.graft_rewrite}
\par\penalty0

This defines the continued Source data for the construction just described. Boundary positions are kept distinct from newly inserted internal vertices.
\lean{Copying/ParametricityCopyGraftExecution}{37}{40}{Parametricity.CopyExecution.continuedSource}
\par\penalty0

The continued Source Iso map is an ordered recursive graph isomorphism between the stated original constructions.
\lean{Copying/ParametricityCopyGraftExecution}{42}{48}{Parametricity.CopyExecution.continuedSourceIso}
\par\penalty0

Private sharing and store remain inside the original family graft and inside its computed successor; no closed-world or single-use condition occurs.
For the proof, the preceding component lemmas are composed or specialized at the stated parameters.
\lean{Copying/ParametricityCopyGraftExecution}{52}{57}{Parametricity.CopyExecution.continued_graft_rewrite}
\par\penalty0

\subsection{Parametricity Copy Growth}
Duplication increases the actual Figure20 abstraction family and grafts the complete canonical term D1,2 into its selected private-sharing input. The private cell and its canonical atom distributor remain literally present.

Two copies replace exactly the selected first private-variable use. The original old tree is retained as the whole downstream component.
\lean{Copying/ParametricityCopyGrowth}{21}{21}{Parametricity.CopyGrowth.tree}
\par\penalty0

The old private continuation has only its equal list presentation exposed.
\lean{Copying/ParametricityCopyGrowth}{24}{25}{Parametricity.CopyGrowth.tail}
\par\penalty0

Both selected abstractions precede the untouched original family.
\lean{Copying/ParametricityCopyGrowth}{28}{29}{Parametricity.CopyGrowth.doubled}
\par\penalty0

This defines the before data for the construction just described. Boundary positions are kept distinct from newly inserted internal vertices.
\lean{Copying/ParametricityCopyGrowth}{31}{32}{Parametricity.CopyGrowth.before}
\par\penalty0

This defines the after data for the construction just described. Boundary positions are kept distinct from newly inserted internal vertices.
\lean{Copying/ParametricityCopyGrowth}{34}{35}{Parametricity.CopyGrowth.after}
\par\penalty0

The tail Iso map is an ordered recursive graph isomorphism between the stated original constructions.
\lean{Copying/ParametricityCopyGrowth}{37}{40}{Parametricity.CopyGrowth.tailIso}
\par\penalty0

Exact graph normalization, not merely a sharing law: the full actual copying distributor becomes part of the original private contraction tree.
\lean{Copying/ParametricityCopyGrowth}{44}{59}{Parametricity.CopyGrowth.iso}
\par\penalty0

The right-hand family has no private outputs. Removing only the empty resource carriers leaves the two unchanged selected abstractions and rest.
\lean{Copying/ParametricityCopyGrowth}{63}{66}{Parametricity.CopyGrowth.closedIso}
\par\penalty0

\subsection{Parametricity Copy Normalization}
The actual original copying target is exactly an enlarged two-abstraction family inside the common residual-input context. This retains every canonical resource edge and allows arbitrary residual input arity, including zero.

The raw Occurrence map is a complete occurrence, preserving original labels and every ordered source and target attachment.
\lean{Copying/ParametricityCopyNormalization}{15}{20}{Parametricity.CopyNormalization.rawOccurrence}
\par\penalty0

The raw input assertion computes the indicated ordered boundary coordinate in the assembled graph.
For the proof, constructor and boundary cases are separated explicitly; the remaining natural-number index bounds follow by linear arithmetic; the defined maps and incidence equations are substituted.
\lean{Copying/ParametricityCopyNormalization}{22}{50}{Parametricity.CopyNormalization.raw_input}
\par\penalty0

The raw output assertion computes the indicated ordered boundary coordinate in the assembled graph.
For the proof, the defined maps and incidence equations are substituted.
\lean{Copying/ParametricityCopyNormalization}{52}{58}{Parametricity.CopyNormalization.raw_output}
\par\penalty0

The occurrence map is a complete occurrence, preserving original labels and every ordered source and target attachment.
\lean{Copying/ParametricityCopyNormalization}{60}{63}{Parametricity.CopyNormalization.occurrence}
\par\penalty0

The input assertion computes the indicated ordered boundary coordinate in the assembled graph.
For the proof, the defined maps and incidence equations are substituted.
\lean{Copying/ParametricityCopyNormalization}{65}{68}{Parametricity.CopyNormalization.input}
\par\penalty0

The output assertion computes the indicated ordered boundary coordinate in the assembled graph.
For the proof, the preceding component lemmas are composed or specialized at the stated parameters.
\lean{Copying/ParametricityCopyNormalization}{70}{71}{Parametricity.CopyNormalization.output}
\par\penalty0

Every old residual-tree, both actual payload, and every distributor edge is covered, so this normalization deletes or duplicates no graph component.
For the proof, constructor and boundary cases are separated explicitly; the defined maps and incidence equations are substituted.
\lean{Copying/ParametricityCopyNormalization}{75}{87}{Parametricity.CopyNormalization.cover}
\par\penalty0

The iso map is an ordered recursive graph isomorphism between the stated original constructions.
\lean{Copying/ParametricityCopyNormalization}{89}{92}{Parametricity.CopyNormalization.iso}
\par\penalty0

The normalization identifies the literal unchanged Figure 17 target.
\lean{Copying/ParametricityCopyNormalization}{95}{99}{Parametricity.CopyNormalization.afterIso}
\par\penalty0

One original compute step reaches the normalized enlarged pair. No translation or operational rule is changed by this exact representation.
For the proof, the defined maps and incidence equations are substituted.
\lean{Copying/ParametricityCopyNormalization}{103}{108}{Parametricity.CopyNormalization.rewrite}
\par\penalty0

\subsection{Parametricity Copy Residual}
The selected-copy residual keeps a single enlarged Figure20 specimen inside the actual common upstream residual context. Both compute steps are original Figure17 rules; the private cell is never replaced by an abstract store.

This defines the grown Left data for the construction just described. Boundary positions are kept distinct from newly inserted internal vertices.
\lean{Copying/ParametricityCopyResidual}{21}{22}{Parametricity.CopyGrowth.grownLeft}
\par\penalty0

Express the exact private-resource normalization at the common enlarged family interface expected by the residual-input context.
\lean{Copying/ParametricityCopyResidual}{26}{32}{Parametricity.CopyGrowth.normalizedIso}
\par\penalty0

This defines the grown Right data for the construction just described. Boundary positions are kept distinct from newly inserted internal vertices.
\lean{Copying/ParametricityCopyResidual}{34}{36}{Parametricity.CopyGrowth.grownRight}
\par\penalty0

The normalized Closed Iso map is an ordered recursive graph isomorphism between the stated original constructions.
\lean{Copying/ParametricityCopyResidual}{38}{41}{Parametricity.CopyGrowth.normalizedClosedIso}
\par\penalty0

An independently proved enlarged-family certificate lifts through exactly one common residual-input context, retaining the active and all other inputs.
\lean{Copying/ParametricityCopyResidual}{45}{54}{Parametricity.CopyGrowth.witness}
\par\penalty0

Actual original copying on both complete grafted families, together with its original enlarged-family specimen. The premises on copied bodies are the original copyability judgments; no machine or template law is assumed.
For the proof, the remaining natural-number index bounds follow by linear arithmetic; the defined maps and incidence equations are substituted.
\lean{Copying/ParametricityCopyResidual}{59}{77}{Parametricity.CopyGrowth.diagram}
\par\penalty0

\subsection{Parametricity Copy Source}
Exact isolation of an upstream sharing graph feeding the selected first abstraction of an arbitrary ordered family. No unused input is discarded.

This defines the ordered target list for the construction just described. Boundary positions are kept distinct from newly inserted internal vertices.
\lean{Copying/ParametricityCopySource}{15}{16}{Parametricity.CopySource.target}
\par\penalty0

This defines the whole data for the construction just described. Boundary positions are kept distinct from newly inserted internal vertices.
\lean{Copying/ParametricityCopySource}{17}{18}{Parametricity.CopySource.whole}
\par\penalty0

The left map is a complete occurrence, preserving original labels and every ordered source and target attachment.
\lean{Copying/ParametricityCopySource}{19}{20}{Parametricity.CopySource.left}
\par\penalty0

The right map is a complete occurrence, preserving original labels and every ordered source and target attachment.
\lean{Copying/ParametricityCopySource}{21}{22}{Parametricity.CopySource.right}
\par\penalty0

The upstream map is a complete occurrence, preserving original labels and every ordered source and target attachment.
\lean{Copying/ParametricityCopySource}{23}{24}{Parametricity.CopySource.upstream}
\par\penalty0

The selected map is a complete occurrence, preserving original labels and every ordered source and target attachment.
\lean{Copying/ParametricityCopySource}{25}{26}{Parametricity.CopySource.selected}
\par\penalty0

The selected disjoint assertion excludes a shared edge or coordinate between the indicated independently inserted components.
For the proof, constructor and boundary cases are separated explicitly.
\lean{Copying/ParametricityCopySource}{27}{28}{Parametricity.CopySource.selected_disjoint}
\par\penalty0

The family map is a complete occurrence, preserving original labels and every ordered source and target attachment.
\lean{Copying/ParametricityCopySource}{29}{31}{Parametricity.CopySource.family}
\par\penalty0

The family first assertion is the compatibility law for the indicated component of this construction. It relates the actual finite coordinates or certificates, with the hypotheses stated in the linked declaration.
For the proof, the defined maps and incidence equations are substituted.
\lean{Copying/ParametricityCopySource}{33}{36}{Parametricity.CopySource.family_first}
\par\penalty0

The family rest assertion is the compatibility law for the indicated component of this construction. It relates the actual finite coordinates or certificates, with the hypotheses stated in the linked declaration.
For the proof, the remaining natural-number index bounds follow by linear arithmetic; the defined maps and incidence equations are substituted.
\lean{Copying/ParametricityCopySource}{37}{45}{Parametricity.CopySource.family_rest}
\par\penalty0

The link assertion is the compatibility law for the indicated component of this construction. It relates the actual finite coordinates or certificates, with the hypotheses stated in the linked declaration.
For the proof, the defined maps and incidence equations are substituted.
\lean{Copying/ParametricityCopySource}{46}{48}{Parametricity.CopySource.link}
\par\penalty0

The disjoint assertion excludes a shared edge or coordinate between the indicated independently inserted components.
For the proof, constructor and boundary cases are separated explicitly.
\lean{Copying/ParametricityCopySource}{49}{54}{Parametricity.CopySource.disjoint}
\par\penalty0

The occurrence map is a complete occurrence, preserving original labels and every ordered source and target attachment.
\lean{Copying/ParametricityCopySource}{56}{58}{Parametricity.CopySource.occurrence}
\par\penalty0

The input assertion computes the indicated ordered boundary coordinate in the assembled graph.
For the proof, constructor and boundary cases are separated explicitly; the remaining natural-number index bounds follow by linear arithmetic; the defined maps and incidence equations are substituted.
\lean{Copying/ParametricityCopySource}{60}{85}{Parametricity.CopySource.input}
\par\penalty0

The family output assertion computes the indicated ordered boundary coordinate in the assembled graph.
For the proof, constructor and boundary cases are separated explicitly; the defined maps and incidence equations are substituted.
\lean{Copying/ParametricityCopySource}{87}{101}{Parametricity.CopySource.family_output}
\par\penalty0

The output assertion computes the indicated ordered boundary coordinate in the assembled graph.
For the proof, the defined maps and incidence equations are substituted.
\lean{Copying/ParametricityCopySource}{103}{106}{Parametricity.CopySource.output}
\par\penalty0

The cover assertion states that every target coordinate has an original preimage under the indicated map.
For the proof, constructor and boundary cases are separated explicitly; the defined maps and incidence equations are substituted.
\lean{Copying/ParametricityCopySource}{108}{115}{Parametricity.CopySource.cover}
\par\penalty0

The iso map is an ordered recursive graph isomorphism between the stated original constructions.
\lean{Copying/ParametricityCopySource}{117}{118}{Parametricity.CopySource.iso}
\par\penalty0

\subsection{Parametricity Double Family}
Exact enlarged Figure20 families after duplicating either abstraction. Both actual copies precede the remaining original family. The simultaneous input/private-output action is proved, so the private tree can be reindexed without changing the observational interface.

Tensor reassociation has the same boundary action on the two families and preserves the left private-output/input correspondence.
For the proof, the defined maps and incidence equations are substituted.
\lean{Copying/ParametricityDoubleFamily}{15}{32}{Parametricity.FamilyLiteral.association_synchronous}
\par\penalty0

Two actual left large abstractions, in their copy order, are followed by all original remaining abstractions. No box body is changed.
\lean{Copying/ParametricityDoubleFamily}{36}{40}{Parametricity.FamilyLiteral.leftLargeDouble}
\par\penalty0

This defines the right Large Double data for the construction just described. Boundary positions are kept distinct from newly inserted internal vertices.
\lean{Copying/ParametricityDoubleFamily}{42}{46}{Parametricity.FamilyLiteral.rightLargeDouble}
\par\penalty0

The two selected small abstractions are moved together past the original large block, with the same move on their two private outputs.
\lean{Copying/ParametricityDoubleFamily}{50}{54}{Parametricity.FamilyLiteral.leftSmallDouble}
\par\penalty0

This defines the right Small Double data for the construction just described. Boundary positions are kept distinct from newly inserted internal vertices.
\lean{Copying/ParametricityDoubleFamily}{56}{60}{Parametricity.FamilyLiteral.rightSmallDouble}
\par\penalty0

The enlarged large families have a single common external input action, which is also the left private-output action.
For the proof, the preceding component lemmas are composed or specialized at the stated parameters.
\lean{Copying/ParametricityDoubleFamily}{64}{68}{Parametricity.FamilyLiteral.largeDouble_synchronous}
\par\penalty0

The enlarged small families use the same complete two-member rotation on both sides and on the original private-tree attachments.
For the proof, the preceding component lemmas are composed or specialized at the stated parameters.
\lean{Copying/ParametricityDoubleFamily}{72}{76}{Parametricity.FamilyLiteral.smallDouble_synchronous}
\par\penalty0

\subsection{Parametricity Family Blocks}
The two ordered Figure20 families are separated as their actual complete large and small blocks, with exact input and private-output offsets.

The original tensor of large and small copies, exposing its arithmetic interface equality uniformly for the one-output and zero-output families.
\lean{Copying/ParametricityFamilyBlocks}{13}{16}{Parametricity.FamilyComponents.family}
\par\penalty0

The large Block map is a complete occurrence, preserving original labels and every ordered source and target attachment.
\lean{Copying/ParametricityFamilyBlocks}{18}{19}{Parametricity.FamilyComponents.largeBlock}
\par\penalty0

The small Block map is a complete occurrence, preserving original labels and every ordered source and target attachment.
\lean{Copying/ParametricityFamilyBlocks}{20}{21}{Parametricity.FamilyComponents.smallBlock}
\par\penalty0

Whole blocks retain their literal ordered coordinates.
Both sides are definitionally equal.
\lean{Copying/ParametricityFamilyBlocks}{24}{25}{Parametricity.FamilyComponents.largeBlock_input}
\par\penalty0

The small Block input assertion computes the indicated ordered boundary coordinate in the assembled graph.
For the proof, the remaining natural-number index bounds follow by linear arithmetic; the defined maps and incidence equations are substituted.
\lean{Copying/ParametricityFamilyBlocks}{27}{29}{Parametricity.FamilyComponents.smallBlock_input}
\par\penalty0

The large Block output assertion computes the indicated ordered boundary coordinate in the assembled graph.
Both sides are definitionally equal.
\lean{Copying/ParametricityFamilyBlocks}{31}{33}{Parametricity.FamilyComponents.largeBlock_output}
\par\penalty0

The small Block output assertion computes the indicated ordered boundary coordinate in the assembled graph.
For the proof, the remaining natural-number index bounds follow by linear arithmetic; the defined maps and incidence equations are substituted.
\lean{Copying/ParametricityFamilyBlocks}{35}{38}{Parametricity.FamilyComponents.smallBlock_output}
\par\penalty0

The large map is a complete occurrence, preserving original labels and every ordered source and target attachment.
\lean{Copying/ParametricityFamilyBlocks}{40}{41}{Parametricity.FamilyComponents.large}
\par\penalty0

The small map is a complete occurrence, preserving original labels and every ordered source and target attachment.
\lean{Copying/ParametricityFamilyBlocks}{42}{43}{Parametricity.FamilyComponents.small}
\par\penalty0

The large input assertion computes the indicated ordered boundary coordinate in the assembled graph.
For the proof, the defined maps and incidence equations are substituted.
\lean{Copying/ParametricityFamilyBlocks}{45}{48}{Parametricity.FamilyComponents.large_input}
\par\penalty0

The small input assertion computes the indicated ordered boundary coordinate in the assembled graph.
For the proof, the defined maps and incidence equations are substituted.
\lean{Copying/ParametricityFamilyBlocks}{50}{53}{Parametricity.FamilyComponents.small_input}
\par\penalty0

The private-output row follows exactly the same large-block ordering.
For the proof, the remaining natural-number index bounds follow by linear arithmetic; the defined maps and incidence equations are substituted.
\lean{Copying/ParametricityFamilyBlocks}{56}{63}{Parametricity.FamilyComponents.large_output}
\par\penalty0

Small private outputs follow all large outputs and retain their own within-copy ordering. Empty-output families satisfy the same statement.
For the proof, the remaining natural-number index bounds follow by linear arithmetic; the defined maps and incidence equations are substituted.
\lean{Copying/ParametricityFamilyBlocks}{67}{74}{Parametricity.FamilyComponents.small_output}
\par\penalty0

Large/small component ownership never aliases, independently of their bodies or of any isomorphism between two lambda representations.
For the proof, constructor and boundary cases are separated explicitly.
\lean{Copying/ParametricityFamilyBlocks}{78}{82}{Parametricity.FamilyComponents.blocks_disjoint}
\par\penalty0

Every family edge belongs to one actual original indexed abstraction.
For the proof, constructor and boundary cases are separated explicitly.
\lean{Copying/ParametricityFamilyBlocks}{85}{93}{Parametricity.FamilyComponents.family_cover}
\par\penalty0

\subsection{Parametricity Family Boundary}
Coincidence of the left private-output permutation and the two external input permutations is proved from the literal component coordinates.

The first-small cyclic move, as an explicit numerical coordinate map.
\lean{Copying/ParametricityFamilyBoundary}{10}{10}{Parametricity.FamilyLiteral.smallIndex}
\par\penalty0

The whole ordered left input action of first-large isolation.
For the proof, constructor and boundary cases are separated explicitly; the remaining natural-number index bounds follow by linear arithmetic; the defined maps and incidence equations are substituted.
\lean{Copying/ParametricityFamilyBoundary}{13}{32}{Parametricity.FamilyLiteral.leftLarge_inputs_val}
\par\penalty0

The whole ordered right input action of first-large isolation.
For the proof, constructor and boundary cases are separated explicitly; the remaining natural-number index bounds follow by linear arithmetic; the defined maps and incidence equations are substituted.
\lean{Copying/ParametricityFamilyBoundary}{35}{54}{Parametricity.FamilyLiteral.rightLarge_inputs_val}
\par\penalty0

The whole ordered left output action of first-large isolation.
For the proof, constructor and boundary cases are separated explicitly; the remaining natural-number index bounds follow by linear arithmetic; the defined maps and incidence equations are substituted.
\lean{Copying/ParametricityFamilyBoundary}{57}{76}{Parametricity.FamilyLiteral.leftLarge_outputs_val}
\par\penalty0

The whole ordered left input action of first-small isolation.
For the proof, constructor and boundary cases are separated explicitly; the remaining natural-number index bounds follow by linear arithmetic; the defined maps and incidence equations are substituted.
\lean{Copying/ParametricityFamilyBoundary}{79}{98}{Parametricity.FamilyLiteral.leftSmall_inputs_val}
\par\penalty0

The whole ordered right input action of first-small isolation.
For the proof, constructor and boundary cases are separated explicitly; the remaining natural-number index bounds follow by linear arithmetic; the defined maps and incidence equations are substituted.
\lean{Copying/ParametricityFamilyBoundary}{101}{120}{Parametricity.FamilyLiteral.rightSmall_inputs_val}
\par\penalty0

The whole ordered left output action of first-small isolation.
For the proof, constructor and boundary cases are separated explicitly; the remaining natural-number index bounds follow by linear arithmetic; the defined maps and incidence equations are substituted.
\lean{Copying/ParametricityFamilyBoundary}{123}{142}{Parametricity.FamilyLiteral.leftSmall_outputs_val}
\par\penalty0

Both sides of Figure20 perform the same external input permutation.
For the proof, the defined maps and incidence equations are substituted.
\lean{Copying/ParametricityFamilyBoundary}{145}{148}{Parametricity.FamilyLiteral.large_inputs}
\par\penalty0

The private sharing-tree attachment follows the input permutation exactly.
For the proof, the defined maps and incidence equations are substituted.
\lean{Copying/ParametricityFamilyBoundary}{151}{154}{Parametricity.FamilyLiteral.large_private}
\par\penalty0

Both sides of Figure20 perform the same external input permutation.
For the proof, the defined maps and incidence equations are substituted.
\lean{Copying/ParametricityFamilyBoundary}{157}{160}{Parametricity.FamilyLiteral.small_inputs}
\par\penalty0

The private sharing-tree attachment follows the input permutation exactly.
For the proof, the defined maps and incidence equations are substituted.
\lean{Copying/ParametricityFamilyBoundary}{163}{166}{Parametricity.FamilyLiteral.small_private}
\par\penalty0

\subsection{Parametricity Family Components}
Indexed complete occurrences of the original ordered abstraction copies. The recursion is the literal Figure20 family construction, including empty families; no alternate translation or family-order convention is introduced.

The selected original member of a finite ordered family.
\lean{Copying/ParametricityFamilyComponents}{14}{23}{Parametricity.FamilyComponents.copy}
\par\penalty0

Every copy consumes precisely its original ordered external input.
For the proof, induction follows the original derivation or finite structure; the remaining natural-number index bounds follow by linear arithmetic; the defined maps and incidence equations are substituted.
\lean{Copying/ParametricityFamilyComponents}{26}{44}{Parametricity.FamilyComponents.copy_input}
\par\penalty0

All outputs of a copy keep their within-copy order and the copy blocks remain ordered; for the paper's families k is respectively one or zero.
For the proof, induction follows the original derivation or finite structure; the remaining natural-number index bounds follow by linear arithmetic; the defined maps and incidence equations are substituted.
\lean{Copying/ParametricityFamilyComponents}{48}{79}{Parametricity.FamilyComponents.copy_output}
\par\penalty0

Different original members have disjoint complete edge sets, including recursive lambda box edges.
For the proof, induction follows the original derivation or finite structure; constructor and boundary cases are separated explicitly; the remaining natural-number index bounds follow by linear arithmetic.
\lean{Copying/ParametricityFamilyComponents}{83}{113}{Parametricity.FamilyComponents.copy_disjoint}
\par\penalty0

Every actual family edge belongs to one of its original indexed members; zero families contribute no hidden identity or administrative edge.
For the proof, induction follows the original derivation or finite structure; constructor and boundary cases are separated explicitly; the remaining natural-number index bounds follow by linear arithmetic.
\lean{Copying/ParametricityFamilyComponents}{117}{130}{Parametricity.FamilyComponents.copy_cover}
\par\penalty0

\subsection{Parametricity Family Coordinates}
Exact boundary coordinates of the original ordered abstraction families. Small-member isolation rotates inputs and their private outputs together.

Removing the first large member shifts every remaining input by one.
For the proof, constructor and boundary cases are separated explicitly; the remaining natural-number index bounds follow by linear arithmetic; the defined maps and incidence equations are substituted.
\lean{Copying/ParametricityFamilyCoordinates}{11}{28}{Parametricity.FamilyComponents.remainingLarge_input}
\par\penalty0

Removing the first small member leaves the large inputs fixed and shifts only the remaining small inputs past the removed member.
For the proof, constructor and boundary cases are separated explicitly; the remaining natural-number index bounds follow by linear arithmetic; the defined maps and incidence equations are substituted.
\lean{Copying/ParametricityFamilyCoordinates}{32}{51}{Parametricity.FamilyComponents.remainingSmall_input}
\par\penalty0

All remaining private-output blocks shift by the selected large block's output count, including the zero-output family.
For the proof, constructor and boundary cases are separated explicitly; the remaining natural-number index bounds follow by linear arithmetic; the defined maps and incidence equations are substituted.
\lean{Copying/ParametricityFamilyCoordinates}{55}{73}{Parametricity.FamilyComponents.remainingLarge_output}
\par\penalty0

Only the small private-output block moves past the selected small member. For one-output lambdas this is precisely the input permutation above.
For the proof, constructor and boundary cases are separated explicitly; the remaining natural-number index bounds follow by linear arithmetic; the defined maps and incidence equations are substituted.
\lean{Copying/ParametricityFamilyCoordinates}{77}{97}{Parametricity.FamilyComponents.remainingSmall_output}
\par\penalty0

The selected large abstraction occupies the original first input.
For the proof, the defined maps and incidence equations are substituted.
\lean{Copying/ParametricityFamilyCoordinates}{100}{104}{Parametricity.FamilyComponents.isolateLarge_first_input}
\par\penalty0

The selected small abstraction is the original input after all large ones.
For the proof, the defined maps and incidence equations are substituted.
\lean{Copying/ParametricityFamilyCoordinates}{107}{112}{Parametricity.FamilyComponents.isolateSmall_first_input}
\par\penalty0

Every other isolated-large input retains its original ordered index.
For the proof, the remaining natural-number index bounds follow by linear arithmetic; the defined maps and incidence equations are substituted.
\lean{Copying/ParametricityFamilyCoordinates}{115}{124}{Parametricity.FamilyComponents.isolateLarge_rest_input}
\par\penalty0

The selected-small rotation shifts large positions back one and leaves later small positions unchanged.
For the proof, the remaining natural-number index bounds follow by linear arithmetic; the defined maps and incidence equations are substituted.
\lean{Copying/ParametricityFamilyCoordinates}{128}{137}{Parametricity.FamilyComponents.isolateSmall_rest_input}
\par\penalty0

Outputs of the first large member remain at the first original output block.
For the proof, the remaining natural-number index bounds follow by linear arithmetic; the defined maps and incidence equations are substituted.
\lean{Copying/ParametricityFamilyCoordinates}{140}{145}{Parametricity.FamilyComponents.isolateLarge_first_output}
\par\penalty0

The selected small output block follows precisely the original large block.
For the proof, the remaining natural-number index bounds follow by linear arithmetic; the defined maps and incidence equations are substituted.
\lean{Copying/ParametricityFamilyCoordinates}{148}{153}{Parametricity.FamilyComponents.isolateSmall_first_output}
\par\penalty0

The remaining outputs under large isolation keep their original indices.
For the proof, the remaining natural-number index bounds follow by linear arithmetic; the defined maps and incidence equations are substituted.
\lean{Copying/ParametricityFamilyCoordinates}{156}{165}{Parametricity.FamilyComponents.isolateLarge_rest_output}
\par\penalty0

The remaining small-isolation outputs follow the same block rotation as inputs, with arbitrary output multiplicity k and its zero-output case.
For the proof, the remaining natural-number index bounds follow by linear arithmetic; the defined maps and incidence equations are substituted.
\lean{Copying/ParametricityFamilyCoordinates}{169}{180}{Parametricity.FamilyComponents.isolateSmall_rest_output}
\par\penalty0

\subsection{Parametricity Family Isolation}
Exact isolation of the first large or first small original abstraction. The residual family keeps every other original abstraction in large-then-small order. Its actual boundary action records the selected-input cyclic move.

The literal tail of a nonempty ordered copy block.
\lean{Copying/ParametricityFamilyIsolation}{13}{16}{Parametricity.FamilyComponents.tail}
\par\penalty0

The tail input assertion computes the indicated ordered boundary coordinate in the assembled graph.
For the proof, the remaining natural-number index bounds follow by linear arithmetic; the defined maps and incidence equations are substituted.
\lean{Copying/ParametricityFamilyIsolation}{18}{20}{Parametricity.FamilyComponents.tail_input}
\par\penalty0

The tail output assertion computes the indicated ordered boundary coordinate in the assembled graph.
For the proof, the remaining natural-number index bounds follow by linear arithmetic; the defined maps and incidence equations are substituted.
\lean{Copying/ParametricityFamilyIsolation}{22}{25}{Parametricity.FamilyComponents.tail_output}
\par\penalty0

The large Tail map is a complete occurrence, preserving original labels and every ordered source and target attachment.
\lean{Copying/ParametricityFamilyIsolation}{27}{28}{Parametricity.FamilyComponents.largeTail}
\par\penalty0

The small Tail map is a complete occurrence, preserving original labels and every ordered source and target attachment.
\lean{Copying/ParametricityFamilyIsolation}{29}{30}{Parametricity.FamilyComponents.smallTail}
\par\penalty0

The large Tail disjoint assertion excludes a shared edge or coordinate between the indicated independently inserted components.
For the proof, constructor and boundary cases are separated explicitly.
\lean{Copying/ParametricityFamilyIsolation}{32}{36}{Parametricity.FamilyComponents.largeTail_disjoint}
\par\penalty0

The small Tail disjoint assertion excludes a shared edge or coordinate between the indicated independently inserted components.
For the proof, constructor and boundary cases are separated explicitly.
\lean{Copying/ParametricityFamilyIsolation}{38}{42}{Parametricity.FamilyComponents.smallTail_disjoint}
\par\penalty0

Delete the first large member only; the residual blocks keep their original large-then-small ordering and every complete recursive label.
\lean{Copying/ParametricityFamilyIsolation}{46}{48}{Parametricity.FamilyComponents.remainingLarge}
\par\penalty0

Delete the first small member only, retaining the complete original large block before the remaining original small block.
\lean{Copying/ParametricityFamilyIsolation}{52}{54}{Parametricity.FamilyComponents.remainingSmall}
\par\penalty0

The remaining Large inputs assertion computes the indicated ordered boundary coordinate in the assembled graph.
For the proof, the defined maps and incidence equations are substituted.
\lean{Copying/ParametricityFamilyIsolation}{56}{66}{Parametricity.FamilyComponents.remainingLarge_inputs}
\par\penalty0

The remaining Small inputs assertion computes the indicated ordered boundary coordinate in the assembled graph.
For the proof, the defined maps and incidence equations are substituted.
\lean{Copying/ParametricityFamilyIsolation}{68}{78}{Parametricity.FamilyComponents.remainingSmall_inputs}
\par\penalty0

The remaining Large outputs assertion computes the indicated ordered boundary coordinate in the assembled graph.
For the proof, the defined maps and incidence equations are substituted.
\lean{Copying/ParametricityFamilyIsolation}{80}{90}{Parametricity.FamilyComponents.remainingLarge_outputs}
\par\penalty0

The remaining Small outputs assertion computes the indicated ordered boundary coordinate in the assembled graph.
For the proof, the defined maps and incidence equations are substituted.
\lean{Copying/ParametricityFamilyIsolation}{92}{102}{Parametricity.FamilyComponents.remainingSmall_outputs}
\par\penalty0

The first Large disjoint assertion excludes a shared edge or coordinate between the indicated independently inserted components.
For the proof, constructor and boundary cases are separated explicitly.
\lean{Copying/ParametricityFamilyIsolation}{104}{107}{Parametricity.FamilyComponents.firstLarge_disjoint}
\par\penalty0

The first Small disjoint assertion excludes a shared edge or coordinate between the indicated independently inserted components.
For the proof, constructor and boundary cases are separated explicitly.
\lean{Copying/ParametricityFamilyIsolation}{109}{112}{Parametricity.FamilyComponents.firstSmall_disjoint}
\par\penalty0

The selected large abstraction followed by the remaining ordered family is an actual complete occurrence covering the whole original family.
\lean{Copying/ParametricityFamilyIsolation}{116}{119}{Parametricity.FamilyComponents.isolateLarge}
\par\penalty0

The selected small abstraction moves to the front explicitly. The boundary isomorphism below records that move rather than pretending it fixes all original external coordinates.
\lean{Copying/ParametricityFamilyIsolation}{124}{127}{Parametricity.FamilyComponents.isolateSmall}
\par\penalty0

The isolate Large cover assertion states that every target coordinate has an original preimage under the indicated map.
For the proof, constructor and boundary cases are separated explicitly; the defined maps and incidence equations are substituted.
\lean{Copying/ParametricityFamilyIsolation}{129}{136}{Parametricity.FamilyComponents.isolateLarge_cover}
\par\penalty0

The isolate Small cover assertion states that every target coordinate has an original preimage under the indicated map.
For the proof, constructor and boundary cases are separated explicitly; the defined maps and incidence equations are substituted.
\lean{Copying/ParametricityFamilyIsolation}{138}{145}{Parametricity.FamilyComponents.isolateSmall_cover}
\par\penalty0

The literal first-large isolation is a full recursive boundary isomorphism, valid for arbitrary family sizes including empty remainders.
\lean{Copying/ParametricityFamilyIsolation}{149}{156}{Parametricity.FamilyComponents.isolateLargeIso}
\par\penalty0

First-small isolation records the simultaneous ordered input/private output rotation. It also applies to the actual zero-output right family.
\lean{Copying/ParametricityFamilyIsolation}{160}{167}{Parametricity.FamilyComponents.isolateSmallIso}
\par\penalty0

\subsection{Parametricity Family Literal}
The generic family decomposition is identified with the original Figure20 families, preserving all actual edges, boxes, and boundary coordinates.

The generic one-output presentation differs only in arithmetic index proofs from the literal left family.
\lean{Copying/ParametricityFamilyLiteral}{12}{15}{Parametricity.FamilyLiteral.leftOccurrence}
\par\penalty0

The zero-output counterpart retains the original right family exactly.
\lean{Copying/ParametricityFamilyLiteral}{18}{21}{Parametricity.FamilyLiteral.rightOccurrence}
\par\penalty0

This defines the left Iso data for the construction just described. Boundary positions are kept distinct from newly inserted internal vertices.
\lean{Copying/ParametricityFamilyLiteral}{23}{26}{Parametricity.FamilyLiteral.leftIso}
\par\penalty0

This defines the right Iso data for the construction just described. Boundary positions are kept distinct from newly inserted internal vertices.
\lean{Copying/ParametricityFamilyLiteral}{28}{31}{Parametricity.FamilyLiteral.rightIso}
\par\penalty0

The left input assertion computes the indicated ordered boundary coordinate in the assembled graph.
Both sides are definitionally equal.
\lean{Copying/ParametricityFamilyLiteral}{33}{34}{Parametricity.FamilyLiteral.left_input}
\par\penalty0

The right input assertion computes the indicated ordered boundary coordinate in the assembled graph.
Both sides are definitionally equal.
\lean{Copying/ParametricityFamilyLiteral}{36}{37}{Parametricity.FamilyLiteral.right_input}
\par\penalty0

The left output assertion computes the indicated ordered boundary coordinate in the assembled graph.
Both sides are definitionally equal.
\lean{Copying/ParametricityFamilyLiteral}{39}{40}{Parametricity.FamilyLiteral.left_output}
\par\penalty0

The left inv input assertion computes the indicated ordered boundary coordinate in the assembled graph.
For the proof, the defined maps and incidence equations are substituted.
\lean{Copying/ParametricityFamilyLiteral}{42}{44}{Parametricity.FamilyLiteral.left_inv_input}
\par\penalty0

The right inv input assertion computes the indicated ordered boundary coordinate in the assembled graph.
For the proof, the defined maps and incidence equations are substituted.
\lean{Copying/ParametricityFamilyLiteral}{46}{48}{Parametricity.FamilyLiteral.right_inv_input}
\par\penalty0

The left inv output assertion computes the indicated ordered boundary coordinate in the assembled graph.
For the proof, the defined maps and incidence equations are substituted.
\lean{Copying/ParametricityFamilyLiteral}{50}{53}{Parametricity.FamilyLiteral.left_inv_output}
\par\penalty0

The first large abstraction and the literal remaining left family.
\lean{Copying/ParametricityFamilyLiteral}{56}{59}{Parametricity.FamilyLiteral.leftLarge}
\par\penalty0

The first small abstraction and the literal remaining left family; the boundary permutation rotates its private output with its input.
\lean{Copying/ParametricityFamilyLiteral}{63}{66}{Parametricity.FamilyLiteral.leftSmall}
\par\penalty0

Literal large-right isolation with exactly the same input ordering.
\lean{Copying/ParametricityFamilyLiteral}{69}{72}{Parametricity.FamilyLiteral.rightLarge}
\par\penalty0

Literal small-right isolation with the same selected-input rotation.
\lean{Copying/ParametricityFamilyLiteral}{75}{78}{Parametricity.FamilyLiteral.rightSmall}
\par\penalty0

\subsection{Parametricity Family Literal Coordinates}
Literal Figure20 isolation keeps the exact selected and remaining external coordinates; left private outputs undergo the same permutation.

First-large isolation fixes the selected original input.
For the proof, the defined maps and incidence equations are substituted.
\lean{Copying/ParametricityFamilyLiteralCoordinates}{11}{18}{Parametricity.FamilyLiteral.leftLarge_first_input}
\par\penalty0

Right isolation fixes the very same selected first-large input.
For the proof, the defined maps and incidence equations are substituted.
\lean{Copying/ParametricityFamilyLiteralCoordinates}{21}{28}{Parametricity.FamilyLiteral.rightLarge_first_input}
\par\penalty0

First-small left isolation selects precisely the original index m.
For the proof, the defined maps and incidence equations are substituted.
\lean{Copying/ParametricityFamilyLiteralCoordinates}{31}{38}{Parametricity.FamilyLiteral.leftSmall_first_input}
\par\penalty0

First-small right isolation selects that same original index.
For the proof, the defined maps and incidence equations are substituted.
\lean{Copying/ParametricityFamilyLiteralCoordinates}{41}{48}{Parametricity.FamilyLiteral.rightSmall_first_input}
\par\penalty0

Remaining large-isolation inputs preserve their original order.
For the proof, the remaining natural-number index bounds follow by linear arithmetic; the defined maps and incidence equations are substituted.
\lean{Copying/ParametricityFamilyLiteralCoordinates}{51}{67}{Parametricity.FamilyLiteral.leftLarge_rest_input}
\par\penalty0

The right family's remaining input map agrees exactly with the left.
For the proof, the remaining natural-number index bounds follow by linear arithmetic; the defined maps and incidence equations are substituted.
\lean{Copying/ParametricityFamilyLiteralCoordinates}{70}{86}{Parametricity.FamilyLiteral.rightLarge_rest_input}
\par\penalty0

Remaining inputs under first-small isolation retain large-first order and shift only the original remaining small block.
For the proof, the remaining natural-number index bounds follow by linear arithmetic; the defined maps and incidence equations are substituted.
\lean{Copying/ParametricityFamilyLiteralCoordinates}{90}{106}{Parametricity.FamilyLiteral.leftSmall_rest_input}
\par\penalty0

Remaining inputs under first-small isolation retain large-first order and shift only the original remaining small block.
For the proof, the remaining natural-number index bounds follow by linear arithmetic; the defined maps and incidence equations are substituted.
\lean{Copying/ParametricityFamilyLiteralCoordinates}{110}{126}{Parametricity.FamilyLiteral.rightSmall_rest_input}
\par\penalty0

The selected private output follows the same original coordinate as its selected term input.
For the proof, the defined maps and incidence equations are substituted.
\lean{Copying/ParametricityFamilyLiteralCoordinates}{130}{137}{Parametricity.FamilyLiteral.leftLarge_first_output}
\par\penalty0

The selected private output follows the same original coordinate as its selected term input.
For the proof, the defined maps and incidence equations are substituted.
\lean{Copying/ParametricityFamilyLiteralCoordinates}{141}{149}{Parametricity.FamilyLiteral.leftSmall_first_output}
\par\penalty0

Every remaining private output has exactly the same ordered action as the corresponding remaining term input.
For the proof, the remaining natural-number index bounds follow by linear arithmetic; the defined maps and incidence equations are substituted.
\lean{Copying/ParametricityFamilyLiteralCoordinates}{153}{169}{Parametricity.FamilyLiteral.leftLarge_rest_output}
\par\penalty0

Every remaining private output has exactly the same ordered action as the corresponding remaining term input.
For the proof, the remaining natural-number index bounds follow by linear arithmetic; the defined maps and incidence equations are substituted.
\lean{Copying/ParametricityFamilyLiteralCoordinates}{173}{189}{Parametricity.FamilyLiteral.leftSmall_rest_output}
\par\penalty0

\subsection{Parametricity Isolated Family}
Genuine original parametricity specimens with one selected abstraction first, preserving the complete remaining ordered family and private tree.

Only equal-list presentations are exposed; the selected small lambda and every remaining original lambda are unchanged.
\lean{Copying/ParametricityIsolatedFamily}{13}{16}{Parametricity.IsolatedFamily.smallLeftFamily}
\par\penalty0

This defines the small Right Family data for the construction just described. Boundary positions are kept distinct from newly inserted internal vertices.
\lean{Copying/ParametricityIsolatedFamily}{18}{21}{Parametricity.IsolatedFamily.smallRightFamily}
\par\penalty0

This defines the small Left Iso data for the construction just described. Boundary positions are kept distinct from newly inserted internal vertices.
\lean{Copying/ParametricityIsolatedFamily}{23}{25}{Parametricity.IsolatedFamily.smallLeftIso}
\par\penalty0

This defines the small Right Iso data for the construction just described. Boundary positions are kept distinct from newly inserted internal vertices.
\lean{Copying/ParametricityIsolatedFamily}{27}{29}{Parametricity.IsolatedFamily.smallRightIso}
\par\penalty0

The small same inputs assertion computes the indicated ordered boundary coordinate in the assembled graph.
For the proof, the defined maps and incidence equations are substituted.
\lean{Copying/ParametricityIsolatedFamily}{31}{34}{Parametricity.IsolatedFamily.small_same_inputs}
\par\penalty0

An arbitrary original private sharing tree attaches to all selected and remaining free x outputs, in their explicitly isolated order.
\lean{Copying/ParametricityIsolatedFamily}{38}{40}{Parametricity.IsolatedFamily.smallClosed}
\par\penalty0

This is an actual original-template filling, obtained by placing the literal family in a common input permutation context and reindexing C.
\lean{Copying/ParametricityIsolatedFamily}{44}{47}{Parametricity.IsolatedFamily.smallWitness}
\par\penalty0

Only equal-list presentations are exposed; the selected large lambda and every remaining original lambda are unchanged.
\lean{Copying/ParametricityIsolatedFamily}{51}{54}{Parametricity.IsolatedFamily.largeLeftFamily}
\par\penalty0

This defines the large Right Family data for the construction just described. Boundary positions are kept distinct from newly inserted internal vertices.
\lean{Copying/ParametricityIsolatedFamily}{56}{59}{Parametricity.IsolatedFamily.largeRightFamily}
\par\penalty0

This defines the large Left Iso data for the construction just described. Boundary positions are kept distinct from newly inserted internal vertices.
\lean{Copying/ParametricityIsolatedFamily}{61}{63}{Parametricity.IsolatedFamily.largeLeftIso}
\par\penalty0

This defines the large Right Iso data for the construction just described. Boundary positions are kept distinct from newly inserted internal vertices.
\lean{Copying/ParametricityIsolatedFamily}{65}{67}{Parametricity.IsolatedFamily.largeRightIso}
\par\penalty0

The large same inputs assertion computes the indicated ordered boundary coordinate in the assembled graph.
For the proof, the defined maps and incidence equations are substituted.
\lean{Copying/ParametricityIsolatedFamily}{69}{72}{Parametricity.IsolatedFamily.large_same_inputs}
\par\penalty0

An arbitrary original private sharing tree attaches to all selected and remaining free x outputs, in their explicitly isolated order.
\lean{Copying/ParametricityIsolatedFamily}{76}{78}{Parametricity.IsolatedFamily.largeClosed}
\par\penalty0

This is an actual original-template filling, obtained by placing the literal family in a common input permutation context and reindexing C.
\lean{Copying/ParametricityIsolatedFamily}{82}{85}{Parametricity.IsolatedFamily.largeWitness}
\par\penalty0

\subsection{Parametricity Large Family Application}
The larger abstractions open in one original beta step each. Their unevaluated callbacks remain common context around a genuine complete parametricity specimen for the increased small-abstraction family.

The opened left family retains all its remaining original closures, its original arbitrary sharing tree and its literal private cell.
\lean{Copying/ParametricityLargeFamilyApplication}{19}{19}{Parametricity.LargeFamilyApplication.leftAfter}
\par\penalty0

This defines the right After data for the construction just described. Boundary positions are kept distinct from newly inserted internal vertices.
\lean{Copying/ParametricityLargeFamilyApplication}{20}{20}{Parametricity.LargeFamilyApplication.rightAfter}
\par\penalty0

The left Member map is a complete occurrence, preserving original labels and every ordered source and target attachment.
\lean{Copying/ParametricityLargeFamilyApplication}{22}{26}{Parametricity.LargeFamilyApplication.leftMember}
\par\penalty0

The right Member map is a complete occurrence, preserving original labels and every ordered source and target attachment.
\lean{Copying/ParametricityLargeFamilyApplication}{28}{29}{Parametricity.LargeFamilyApplication.rightMember}
\par\penalty0

The left member input assertion computes the indicated ordered boundary coordinate in the assembled graph.
For the proof, the defined maps and incidence equations are substituted.
\lean{Copying/ParametricityLargeFamilyApplication}{31}{35}{Parametricity.LargeFamilyApplication.left_member_input}
\par\penalty0

The right member input assertion computes the indicated ordered boundary coordinate in the assembled graph.
Both sides are definitionally equal.
\lean{Copying/ParametricityLargeFamilyApplication}{36}{36}{Parametricity.LargeFamilyApplication.right_member_input}
\par\penalty0

The complete original opened large-left body, including its unchanged callback and argument, occurs in the actual smaller-family continuation.
\lean{Copying/ParametricityLargeFamilyApplication}{40}{41}{Parametricity.LargeFamilyApplication.leftOpened}
\par\penalty0

The right Opened map is a complete occurrence, preserving original labels and every ordered source and target attachment.
\lean{Copying/ParametricityLargeFamilyApplication}{43}{44}{Parametricity.LargeFamilyApplication.rightOpened}
\par\penalty0

The left opened input assertion computes the indicated ordered boundary coordinate in the assembled graph.
For the proof, the preceding component lemmas are composed or specialized at the stated parameters.
\lean{Copying/ParametricityLargeFamilyApplication}{46}{47}{Parametricity.LargeFamilyApplication.left_opened_input}
\par\penalty0

The right opened input assertion computes the indicated ordered boundary coordinate in the assembled graph.
For the proof, the defined maps and incidence equations are substituted.
\lean{Copying/ParametricityLargeFamilyApplication}{48}{49}{Parametricity.LargeFamilyApplication.right_opened_input}
\par\penalty0

Insert the literal original large beta source in its computed complement.
\lean{Copying/ParametricityLargeFamilyApplication}{52}{53}{Parametricity.LargeFamilyApplication.leftBefore}
\par\penalty0

This defines the right Before data for the construction just described. Boundary positions are kept distinct from newly inserted internal vertices.
\lean{Copying/ParametricityLargeFamilyApplication}{54}{55}{Parametricity.LargeFamilyApplication.rightBefore}
\par\penalty0

The actual left beta rewrite resumes at the original external input.
For the proof, the defined maps and incidence equations are substituted.
\lean{Copying/ParametricityLargeFamilyApplication}{58}{65}{Parametricity.LargeFamilyApplication.left_beta}
\par\penalty0

The right beta rewrite has the same one-step cost; the callback need not terminate and is not evaluated by either proof.
For the proof, the defined maps and incidence equations are substituted.
\lean{Copying/ParametricityLargeFamilyApplication}{69}{76}{Parametricity.LargeFamilyApplication.right_beta}
\par\penalty0

A fully constructed original parametricity filling after the two beta steps: the smaller abstraction count increases, the larger count decreases, and the unevaluated callback plus V remains common context.
\lean{Copying/ParametricityLargeFamilyApplication}{81}{84}{Parametricity.LargeFamilyApplication.residual}
\par\penalty0

This defines the pointed data for the construction just described. Boundary positions are kept distinct from newly inserted internal vertices.
\lean{Copying/ParametricityLargeFamilyApplication}{86}{89}{Parametricity.LargeFamilyApplication.pointed}
\par\penalty0

The complete selected-family large-beta comparison has one genuine original rewrite on each side and an actual common residual specimen.
For the proof, the defined maps and incidence equations are substituted.
\lean{Copying/ParametricityLargeFamilyApplication}{93}{103}{Parametricity.LargeFamilyApplication.one_vs_one}
\par\penalty0

\subsection{Parametricity Large Family Diagram}
The complete one-versus-one larger-beta diagram starts at the actual original application contexts around both full Figure 20 families.

This defines the left Source data for the construction just described. Boundary positions are kept distinct from newly inserted internal vertices.
\lean{Copying/ParametricityLargeFamilyDiagram}{13}{13}{Parametricity.LargeFamilyDiagram.leftSource}
\par\penalty0

This defines the right Source data for the construction just described. Boundary positions are kept distinct from newly inserted internal vertices.
\lean{Copying/ParametricityLargeFamilyDiagram}{14}{14}{Parametricity.LargeFamilyDiagram.rightSource}
\par\penalty0

The actual initial whole-family sources already have a genuine original parametricity certificate, with the unchanged argument retained.
\lean{Copying/ParametricityLargeFamilyDiagram}{18}{20}{Parametricity.LargeFamilyDiagram.sourceWitness}
\par\penalty0

Opening a selected large abstraction reduces its family count by one and increases the small count by one. Both original machine steps preserve the entire unevaluated callback and every argument resource output.
For the proof, the defined maps and incidence equations are substituted.
\lean{Copying/ParametricityLargeFamilyDiagram}{25}{35}{Parametricity.LargeFamilyDiagram.one_vs_one}
\par\penalty0

\subsection{Parametricity Large Family Source}
The computed larger-beta predecessor contains the entire original isolated Figure20 large family, its arbitrary original private tree, and its literal allocated cell. The only external components are application and argument.

This defines the old-frame port map for the construction just described. Boundary positions are kept distinct from newly inserted internal vertices.
\lean{Copying/ParametricityLargeFamilySource}{16}{16}{Parametricity.LargeFamilySource.host}
\par\penalty0

This defines the local Source data for the construction just described. Boundary positions are kept distinct from newly inserted internal vertices.
\lean{Copying/ParametricityLargeFamilySource}{17}{17}{Parametricity.LargeFamilySource.localSource}
\par\penalty0

The inserted map is a complete occurrence, preserving original labels and every ordered source and target attachment.
\lean{Copying/ParametricityLargeFamilySource}{18}{19}{Parametricity.LargeFamilySource.inserted}
\par\penalty0

This defines the head data for the construction just described. Boundary positions are kept distinct from newly inserted internal vertices.
\lean{Copying/ParametricityLargeFamilySource}{20}{22}{Parametricity.LargeFamilySource.head}
\par\penalty0

This defines the arguments data for the construction just described. Boundary positions are kept distinct from newly inserted internal vertices.
\lean{Copying/ParametricityLargeFamilySource}{23}{25}{Parametricity.LargeFamilySource.arguments}
\par\penalty0

This defines the selected data for the construction just described. Boundary positions are kept distinct from newly inserted internal vertices.
\lean{Copying/ParametricityLargeFamilySource}{26}{27}{Parametricity.LargeFamilySource.selected}
\par\penalty0

This defines the argument data for the construction just described. Boundary positions are kept distinct from newly inserted internal vertices.
\lean{Copying/ParametricityLargeFamilySource}{28}{29}{Parametricity.LargeFamilySource.argument}
\par\penalty0

This defines the private Family data for the construction just described. Boundary positions are kept distinct from newly inserted internal vertices.
\lean{Copying/ParametricityLargeFamilySource}{31}{31}{Parametricity.LargeFamilySource.privateFamily}
\par\penalty0

The old Family map is a complete occurrence, preserving original labels and every ordered source and target attachment.
\lean{Copying/ParametricityLargeFamilySource}{32}{37}{Parametricity.LargeFamilySource.oldFamily}
\par\penalty0

The old Tree map is a complete occurrence, preserving original labels and every ordered source and target attachment.
\lean{Copying/ParametricityLargeFamilySource}{38}{42}{Parametricity.LargeFamilySource.oldTree}
\par\penalty0

The old Store map is a complete occurrence, preserving original labels and every ordered source and target attachment.
\lean{Copying/ParametricityLargeFamilySource}{43}{46}{Parametricity.LargeFamilySource.oldStore}
\par\penalty0

The tree store link assertion is the compatibility law for the indicated component of this construction. It relates the actual finite coordinates or certificates, with the hypotheses stated in the linked declaration.
For the proof, the defined maps and incidence equations are substituted.
\lean{Copying/ParametricityLargeFamilySource}{48}{53}{Parametricity.LargeFamilySource.tree_store_link}
\par\penalty0

The tree store disjoint assertion excludes a shared edge or coordinate between the indicated independently inserted components.
For the proof, constructor and boundary cases are separated explicitly.
\lean{Copying/ParametricityLargeFamilySource}{55}{57}{Parametricity.LargeFamilySource.tree_store_disjoint}
\par\penalty0

The old Cell map is a complete occurrence, preserving original labels and every ordered source and target attachment.
\lean{Copying/ParametricityLargeFamilySource}{59}{62}{Parametricity.LargeFamilySource.oldCell}
\par\penalty0

The family avoids assertion is the compatibility law for the indicated component of this construction. It relates the actual finite coordinates or certificates, with the hypotheses stated in the linked declaration.
For the proof, constructor and boundary cases are separated explicitly.
\lean{Copying/ParametricityLargeFamilySource}{64}{70}{Parametricity.LargeFamilySource.family_avoids}
\par\penalty0

The cell avoids assertion is the compatibility law for the indicated component of this construction. It relates the actual finite coordinates or certificates, with the hypotheses stated in the linked declaration.
For the proof, constructor and boundary cases are separated explicitly.
\lean{Copying/ParametricityLargeFamilySource}{72}{78}{Parametricity.LargeFamilySource.cell_avoids}
\par\penalty0

The remaining map is a complete occurrence, preserving original labels and every ordered source and target attachment.
\lean{Copying/ParametricityLargeFamilySource}{80}{82}{Parametricity.LargeFamilySource.remaining}
\par\penalty0

The cell map is a complete occurrence, preserving original labels and every ordered source and target attachment.
\lean{Copying/ParametricityLargeFamilySource}{83}{85}{Parametricity.LargeFamilySource.cell}
\par\penalty0

The selected remaining disjoint assertion excludes a shared edge or coordinate between the indicated independently inserted components.
For the proof, constructor and boundary cases are separated explicitly.
\lean{Copying/ParametricityLargeFamilySource}{87}{89}{Parametricity.LargeFamilySource.selected_remaining_disjoint}
\par\penalty0

The family map is a complete occurrence, preserving original labels and every ordered source and target attachment.
\lean{Copying/ParametricityLargeFamilySource}{91}{93}{Parametricity.LargeFamilySource.family}
\par\penalty0

The selected cell assertion is the compatibility law for the indicated component of this construction. It relates the actual finite coordinates or certificates, with the hypotheses stated in the linked declaration.
For the proof, the defined maps and incidence equations are substituted.
\lean{Copying/ParametricityLargeFamilySource}{95}{116}{Parametricity.LargeFamilySource.selected_cell}
\par\penalty0

The remaining cell assertion is the compatibility law for the indicated component of this construction. It relates the actual finite coordinates or certificates, with the hypotheses stated in the linked declaration.
For the proof, the remaining natural-number index bounds follow by linear arithmetic; the defined maps and incidence equations are substituted.
\lean{Copying/ParametricityLargeFamilySource}{118}{136}{Parametricity.LargeFamilySource.remaining_cell}
\par\penalty0

Before beta the selected large lambda retains exactly its original private variable attachment, although its callback is not part of the private tree.
For the proof, constructor and boundary cases are separated explicitly; the remaining natural-number index bounds follow by linear arithmetic; the defined maps and incidence equations are substituted.
\lean{Copying/ParametricityLargeFamilySource}{140}{155}{Parametricity.LargeFamilySource.family_cell}
\par\penalty0

The family cell disjoint assertion excludes a shared edge or coordinate between the indicated independently inserted components.
For the proof, constructor and boundary cases are separated explicitly; the defined maps and incidence equations are substituted.
\lean{Copying/ParametricityLargeFamilySource}{157}{164}{Parametricity.LargeFamilySource.family_cell_disjoint}
\par\penalty0

This defines the closed Graph data for the construction just described. Boundary positions are kept distinct from newly inserted internal vertices.
\lean{Copying/ParametricityLargeFamilySource}{166}{167}{Parametricity.LargeFamilySource.closedGraph}
\par\penalty0

Every original family abstraction, the entire arbitrary C, and the unchanged allocated cell constitute one complete actual source occurrence.
\lean{Copying/ParametricityLargeFamilySource}{171}{173}{Parametricity.LargeFamilySource.closedFamily}
\par\penalty0

The closed zero assertion is the compatibility law for the indicated component of this construction. It relates the actual finite coordinates or certificates, with the hypotheses stated in the linked declaration.
For the proof, the defined maps and incidence equations are substituted.
\lean{Copying/ParametricityLargeFamilySource}{175}{180}{Parametricity.LargeFamilySource.closed_zero}
\par\penalty0

The head function assertion is the compatibility law for the indicated component of this construction. It relates the actual finite coordinates or certificates, with the hypotheses stated in the linked declaration.
For the proof, the defined maps and incidence equations are substituted.
\lean{Copying/ParametricityLargeFamilySource}{182}{185}{Parametricity.LargeFamilySource.head_function}
\par\penalty0

The head argument assertion is the compatibility law for the indicated component of this construction. It relates the actual finite coordinates or certificates, with the hypotheses stated in the linked declaration.
For the proof, the defined maps and incidence equations are substituted.
\lean{Copying/ParametricityLargeFamilySource}{187}{188}{Parametricity.LargeFamilySource.head_argument}
\par\penalty0

The head family disjoint assertion excludes a shared edge or coordinate between the indicated independently inserted components.
For the proof, constructor and boundary cases are separated explicitly.
\lean{Copying/ParametricityLargeFamilySource}{190}{195}{Parametricity.LargeFamilySource.head_family_disjoint}
\par\penalty0

The head argument disjoint assertion excludes a shared edge or coordinate between the indicated independently inserted components.
For the proof, constructor and boundary cases are separated explicitly.
\lean{Copying/ParametricityLargeFamilySource}{197}{198}{Parametricity.LargeFamilySource.head_argument_disjoint}
\par\penalty0

The family argument disjoint assertion excludes a shared edge or coordinate between the indicated independently inserted components.
For the proof, constructor and boundary cases are separated explicitly.
\lean{Copying/ParametricityLargeFamilySource}{200}{205}{Parametricity.LargeFamilySource.family_argument_disjoint}
\par\penalty0

Che isolated source family is certified by the actual original parametricity pre-template, after the proved explicit input permutation.
\lean{Copying/ParametricityLargeFamilySource}{209}{213}{Parametricity.LargeFamilySource.familyWitness}
\par\penalty0

\subsection{Parametricity Large Family Source Iso}
Complete reconstruction identifies the computed larger-beta source with the original application context around its whole Figure20 family.

The head input assertion computes the indicated ordered boundary coordinate in the assembled graph.
For the proof, the defined maps and incidence equations are substituted.
\lean{Copying/ParametricityLargeFamilySourceIso}{14}{17}{Parametricity.LargeFamilySource.head_input}
\par\penalty0

The remaining input assertion computes the indicated ordered boundary coordinate in the assembled graph.
For the proof, the remaining natural-number index bounds follow by linear arithmetic; the defined maps and incidence equations are substituted.
\lean{Copying/ParametricityLargeFamilySourceIso}{19}{34}{Parametricity.LargeFamilySource.remaining_input}
\par\penalty0

The closed input succ assertion computes the indicated ordered boundary coordinate in the assembled graph.
For the proof, the defined maps and incidence equations are substituted.
\lean{Copying/ParametricityLargeFamilySourceIso}{36}{43}{Parametricity.LargeFamilySource.closed_input_succ}
\par\penalty0

The argument output assertion computes the indicated ordered boundary coordinate in the assembled graph.
For the proof, the remaining natural-number index bounds follow by linear arithmetic; the defined maps and incidence equations are substituted.
\lean{Copying/ParametricityLargeFamilySourceIso}{45}{61}{Parametricity.LargeFamilySource.argument_output}
\par\penalty0

The generic original application context is realized by its actual head, whole original private family, and unchanged arbitrary stable argument.
\lean{Copying/ParametricityLargeFamilySourceIso}{65}{69}{Parametricity.LargeFamilySource.sourceOccurrence}
\par\penalty0

The input assertion computes the indicated ordered boundary coordinate in the assembled graph.
For the proof, constructor and boundary cases are separated explicitly; the remaining natural-number index bounds follow by linear arithmetic; the defined maps and incidence equations are substituted.
\lean{Copying/ParametricityLargeFamilySourceIso}{71}{82}{Parametricity.LargeFamilySource.input}
\par\penalty0

The output assertion computes the indicated ordered boundary coordinate in the assembled graph.
For the proof, the defined maps and incidence equations are substituted.
\lean{Copying/ParametricityLargeFamilySourceIso}{84}{86}{Parametricity.LargeFamilySource.output}
\par\penalty0

The covers head assertion is the compatibility law for the indicated component of this construction. It relates the actual finite coordinates or certificates, with the hypotheses stated in the linked declaration.
For the proof, the defined maps and incidence equations are substituted.
\lean{Copying/ParametricityLargeFamilySourceIso}{89}{90}{Parametricity.LargeFamilySource.covers_head}
\par\penalty0

The covers argument assertion is the compatibility law for the indicated component of this construction. It relates the actual finite coordinates or certificates, with the hypotheses stated in the linked declaration.
For the proof, the defined maps and incidence equations are substituted.
\lean{Copying/ParametricityLargeFamilySourceIso}{91}{92}{Parametricity.LargeFamilySource.covers_argument}
\par\penalty0

The covers selected assertion is the compatibility law for the indicated component of this construction. It relates the actual finite coordinates or certificates, with the hypotheses stated in the linked declaration.
For the proof, the defined maps and incidence equations are substituted.
\lean{Copying/ParametricityLargeFamilySourceIso}{93}{95}{Parametricity.LargeFamilySource.covers_selected}
\par\penalty0

The covers remaining assertion is the compatibility law for the indicated component of this construction. It relates the actual finite coordinates or certificates, with the hypotheses stated in the linked declaration.
For the proof, the defined maps and incidence equations are substituted.
\lean{Copying/ParametricityLargeFamilySourceIso}{96}{98}{Parametricity.LargeFamilySource.covers_remaining}
\par\penalty0

The covers cell assertion is the compatibility law for the indicated component of this construction. It relates the actual finite coordinates or certificates, with the hypotheses stated in the linked declaration.
For the proof, the defined maps and incidence equations are substituted.
\lean{Copying/ParametricityLargeFamilySourceIso}{99}{101}{Parametricity.LargeFamilySource.covers_cell}
\par\penalty0

Before the inverse beta cut, the removed component consists of the application, selected smaller abstraction and complete unevaluated callback invocation. Every other edge belongs to the unchanged remaining family or the original private tree and allocated cell.
For the proof, constructor and boundary cases are separated explicitly; the defined maps and incidence equations are substituted.
\lean{Copying/ParametricityLargeFamilySourceIso}{107}{126}{Parametricity.LargeFamilySource.opened_cover}
\par\penalty0

Every actual source edge is covered by the three original components; this includes all unused sharing siblings and every argument resource edge.
For the proof, constructor and boundary cases are separated explicitly; the defined maps and incidence equations are substituted.
\lean{Copying/ParametricityLargeFamilySourceIso}{130}{158}{Parametricity.LargeFamilySource.cover}
\par\penalty0

The complete source is the original application context around the entire original Figure20 family, with only equal-interface presentation casts. The map fixes every external input and argument output.
\lean{Copying/ParametricityLargeFamilySourceIso}{163}{166}{Parametricity.LargeFamilySource.iso}
\par\penalty0

\subsection{Parametricity Large Right Source}
The larger constant-family predecessor is exactly application to the original larger closed lambda, not merely an unspecified inverse beta cut.

The remaining map is a complete occurrence, preserving original labels and every ordered source and target attachment.
\lean{Copying/ParametricityLargeRightSource}{11}{15}{Parametricity.LargeRightSource.remaining}
\par\penalty0

The apart assertion is the compatibility law for the indicated component of this construction. It relates the actual finite coordinates or certificates, with the hypotheses stated in the linked declaration.
For the proof, constructor and boundary cases are separated explicitly.
\lean{Copying/ParametricityLargeRightSource}{17}{23}{Parametricity.LargeRightSource.apart}
\par\penalty0

The remaining input assertion computes the indicated ordered boundary coordinate in the assembled graph.
For the proof, the remaining natural-number index bounds follow by linear arithmetic; the defined maps and incidence equations are substituted.
\lean{Copying/ParametricityLargeRightSource}{25}{32}{Parametricity.LargeRightSource.remaining_input}
\par\penalty0

The body output assertion computes the indicated ordered boundary coordinate in the assembled graph.
For the proof, the defined maps and incidence equations are substituted.
\lean{Copying/ParametricityLargeRightSource}{34}{47}{Parametricity.LargeRightSource.body_output}
\par\penalty0

The full opened larger body and unchanged remainder cover exactly the application, selected smaller closure, callback, argument, and old family.
For the proof, constructor and boundary cases are separated explicitly; the defined maps and incidence equations are substituted.
\lean{Copying/ParametricityLargeRightSource}{51}{63}{Parametricity.LargeRightSource.cover}
\par\penalty0

Exact original larger-family application source, fixing every external input and every output of the arbitrary callback argument V.
\lean{Copying/ParametricityLargeRightSource}{67}{75}{Parametricity.LargeRightSource.iso}
\par\penalty0

\subsection{Parametricity Pointed Residual}
The private-read correction followed by a genuine parametricity specimen satisfies the paper's directed quasi-specimen clauses. Its intermediate search root is proved, not supplied as an additional assumption.

The parametricity hole Free assertion excludes unfilled holes throughout the indicated graph and every nested box.
For the proof, the preceding component lemmas are composed or specialized at the stated parameters.
\lean{Copying/ParametricityPointedResidual}{11}{14}{BoxPullWeakening.PointedWitness.parametricity_holeFree}
\par\penalty0

Lambda entries return immediately and no private output escapes, so the actual pointed parametricity filling preserves original search rootedness.
For the proof, the defined maps and incidence equations are substituted.
\lean{Copying/ParametricityPointedResidual}{18}{27}{BoxPullWeakening.PointedWitness.parametricity_rooted_iff}
\par\penalty0

Every structural correction retains the actual focus wire's type.
For the proof, the preceding component lemmas are composed or specialized at the stated parameters.
\lean{Copying/ParametricityPointedResidual}{36}{39}{Parametricity.ReadResidual.labels}
\par\penalty0

The actual four-link correction retains closedness of the running host.
For the proof, the preceding component lemmas are composed or specialized at the stated parameters.
\lean{Copying/ParametricityPointedResidual}{42}{44}{Parametricity.ReadResidual.holeFree}
\par\penalty0

A private-read correction followed by an independently constructed remaining-family specimen has the forward improvement quasi-clause. The corrected state's root follows from the opposite endpoint by original search.
For the proof, the defined maps and incidence equations are substituted.
\lean{Copying/ParametricityPointedResidual}{49}{64}{Parametricity.ReadResidual.parametricity_quasi}
\par\penalty0

The symmetrical one-versus-nine branch retains the converse cost inequality, rather than silently replacing it by exact-count equivalence.
For the proof, the defined maps and incidence equations are substituted.
\lean{Copying/ParametricityPointedResidual}{68}{83}{Parametricity.ReadResidual.parametricity_quasi_converse}
\par\penalty0

\subsection{Parametricity Private Atom Residual}
The two untouched atom weakenings in the private-cell complement attach literally to the reconstructed read spine. Closing those inputs produces a real one-input atom tree; the remaining private instance is not discarded.

This defines the spine data for the construction just described. Boundary positions are kept distinct from newly inserted internal vertices.
\lean{Copying/ParametricityPrivateAtomResidual}{12}{12}{Parametricity.PrivateRead.AtomResidual.spine}
\par\penalty0

This defines the local Tree data for the construction just described. Boundary positions are kept distinct from newly inserted internal vertices.
\lean{Copying/ParametricityPrivateAtomResidual}{13}{13}{Parametricity.PrivateRead.AtomResidual.localTree}
\par\penalty0

The inserted map is a complete occurrence, preserving original labels and every ordered source and target attachment.
\lean{Copying/ParametricityPrivateAtomResidual}{15}{16}{Parametricity.PrivateRead.AtomResidual.inserted}
\par\penalty0

The local Cell map is a complete occurrence, preserving original labels and every ordered source and target attachment.
\lean{Copying/ParametricityPrivateAtomResidual}{18}{23}{Parametricity.PrivateRead.AtomResidual.localCell}
\par\penalty0

The tree map is a complete occurrence, preserving original labels and every ordered source and target attachment.
\lean{Copying/ParametricityPrivateAtomResidual}{25}{26}{Parametricity.PrivateRead.AtomResidual.tree}
\par\penalty0

The old Weak map is a complete occurrence, preserving original labels and every ordered source and target attachment.
\lean{Copying/ParametricityPrivateAtomResidual}{28}{33}{Parametricity.PrivateRead.AtomResidual.oldWeak}
\par\penalty0

The weak disjoint assertion excludes a shared edge or coordinate between the indicated independently inserted components.
For the proof, constructor and boundary cases are separated explicitly; the defined maps and incidence equations are substituted.
\lean{Copying/ParametricityPrivateAtomResidual}{35}{45}{Parametricity.PrivateRead.AtomResidual.weak_disjoint}
\par\penalty0

The weak map is a complete occurrence, preserving original labels and every ordered source and target attachment.
\lean{Copying/ParametricityPrivateAtomResidual}{47}{49}{Parametricity.PrivateRead.AtomResidual.weak}
\par\penalty0

The old Weak link assertion is the compatibility law for the indicated component of this construction. It relates the actual finite coordinates or certificates, with the hypotheses stated in the linked declaration.
For the proof, constructor and boundary cases are separated explicitly; the defined maps and incidence equations are substituted.
\lean{Copying/ParametricityPrivateAtomResidual}{51}{54}{Parametricity.PrivateRead.AtomResidual.oldWeak_link}
\par\penalty0

The tree input assertion computes the indicated ordered boundary coordinate in the assembled graph.
For the proof, constructor and boundary cases are separated explicitly; the remaining natural-number index bounds follow by linear arithmetic; the defined maps and incidence equations are substituted.
\lean{Copying/ParametricityPrivateAtomResidual}{56}{59}{Parametricity.PrivateRead.AtomResidual.tree_input}
\par\penalty0

The weak link assertion is the compatibility law for the indicated component of this construction. It relates the actual finite coordinates or certificates, with the hypotheses stated in the linked declaration.
For the proof, constructor and boundary cases are separated explicitly; the defined maps and incidence equations are substituted.
\lean{Copying/ParametricityPrivateAtomResidual}{61}{67}{Parametricity.PrivateRead.AtomResidual.weak_link}
\par\penalty0

The weak tree disjoint assertion excludes a shared edge or coordinate between the indicated independently inserted components.
For the proof, constructor and boundary cases are separated explicitly.
\lean{Copying/ParametricityPrivateAtomResidual}{69}{70}{Parametricity.PrivateRead.AtomResidual.weak_tree_disjoint}
\par\penalty0

The weak pair disjoint assertion excludes a shared edge or coordinate between the indicated independently inserted components.
For the proof, constructor and boundary cases are separated explicitly.
\lean{Copying/ParametricityPrivateAtomResidual}{72}{78}{Parametricity.PrivateRead.AtomResidual.weak_pair_disjoint}
\par\penalty0

Close the old cell's weakening attachment first, keeping both of the other original atom inputs in order.
\lean{Copying/ParametricityPrivateAtomResidual}{82}{83}{Parametricity.PrivateRead.AtomResidual.firstTree}
\par\penalty0

The first map is a complete occurrence, preserving original labels and every ordered source and target attachment.
\lean{Copying/ParametricityPrivateAtomResidual}{85}{88}{Parametricity.PrivateRead.AtomResidual.first}
\par\penalty0

The first input assertion computes the indicated ordered boundary coordinate in the assembled graph.
For the proof, constructor and boundary cases are separated explicitly; the remaining natural-number index bounds follow by linear arithmetic; the defined maps and incidence equations are substituted.
\lean{Copying/ParametricityPrivateAtomResidual}{90}{98}{Parametricity.PrivateRead.AtomResidual.first_input}
\par\penalty0

The weak first disjoint assertion excludes a shared edge or coordinate between the indicated independently inserted components.
For the proof, constructor and boundary cases are separated explicitly; the defined maps and incidence equations are substituted.
\lean{Copying/ParametricityPrivateAtomResidual}{100}{108}{Parametricity.PrivateRead.AtomResidual.weak_first_disjoint}
\par\penalty0

The complete residual private atom network is genuinely a one-input contraction tree. It includes all three consumed-path contractions and all three actual weakening leaves.
\lean{Copying/ParametricityPrivateAtomResidual}{113}{114}{Parametricity.PrivateRead.AtomResidual.fullTree}
\par\penalty0

The full map is a complete occurrence, preserving original labels and every ordered source and target attachment.
\lean{Copying/ParametricityPrivateAtomResidual}{116}{119}{Parametricity.PrivateRead.AtomResidual.full}
\par\penalty0

The full input assertion computes the indicated ordered boundary coordinate in the assembled graph.
For the proof, constructor and boundary cases are separated explicitly; the defined maps and incidence equations are substituted.
\lean{Copying/ParametricityPrivateAtomResidual}{121}{128}{Parametricity.PrivateRead.AtomResidual.full_input}
\par\penalty0

The first output assertion computes the indicated ordered boundary coordinate in the assembled graph.
For the proof, the defined maps and incidence equations are substituted.
\lean{Copying/ParametricityPrivateAtomResidual}{130}{134}{Parametricity.PrivateRead.AtomResidual.first_output}
\par\penalty0

The full output assertion computes the indicated ordered boundary coordinate in the assembled graph.
For the proof, the defined maps and incidence equations are substituted.
\lean{Copying/ParametricityPrivateAtomResidual}{136}{140}{Parametricity.PrivateRead.AtomResidual.full_output}
\par\penalty0

Replace the complete residual atom network by the literal original single-use distributor, leaving its private instance and store attached.
\lean{Copying/ParametricityPrivateAtomResidual}{144}{146}{Parametricity.PrivateRead.AtomResidual.normalized}
\par\penalty0

The witness data give an actual recursive paired filling, with its common context and endpoint representation maps.
\lean{Copying/ParametricityPrivateAtomResidual}{148}{152}{Parametricity.PrivateRead.AtomResidual.witness}
\par\penalty0

This defines the pointed data for the construction just described. Boundary positions are kept distinct from newly inserted internal vertices.
\lean{Copying/ParametricityPrivateAtomResidual}{154}{157}{Parametricity.PrivateRead.AtomResidual.pointed}
\par\penalty0

\subsection{Parametricity Private Residual}
The private read's family-ready residual preserves its literal original cell and complete weakened term tree. Only the proved sharing-tree equation normalizes the three-contraction private atom network.

One returned numeral is separated from the whole remaining private cell. The selected term input is closed by the original residual-tree weakening.
\lean{Copying/ParametricityPrivateResidual}{14}{17}{Parametricity.PrivateRead.result}
\par\penalty0

The cell association map is an ordered recursive graph isomorphism between the stated original constructions.
\lean{Copying/ParametricityPrivateResidual}{19}{26}{Parametricity.PrivateRead.cell_association}
\par\penalty0

The stored association map is an ordered recursive graph isomorphism between the stated original constructions.
\lean{Copying/ParametricityPrivateResidual}{28}{34}{Parametricity.PrivateRead.stored_association}
\par\penalty0

The only structural correction is a literal pair of finite atom trees. Complete reconstruction and serial association recover the unchanged store encoding, including its canonical D1,1 and numeral.
\lean{Copying/ParametricityPrivateResidual}{39}{53}{Parametricity.PrivateRead.sharing_residual}
\par\penalty0

This defines the sharing pointed data for the construction just described. Boundary positions are kept distinct from newly inserted internal vertices.
\lean{Copying/ParametricityPrivateResidual}{55}{58}{Parametricity.PrivateRead.sharing_pointed}
\par\penalty0

The unchanged eight-step execution has a fully constructed residual chain ending in an ordered returned value and the original private cell.
For the proof, the preceding component lemmas are composed or specialized at the stated parameters.
\lean{Copying/ParametricityPrivateResidual}{62}{68}{Parametricity.PrivateRead.residual_factorization}
\par\penalty0

\subsection{Parametricity Private Residual Cover}
Every edge left by the private read is explicitly accounted for. The normal form preserves the complete ordered residual term tree and cell.

The old read rule's primitive-frame presentation and original rule have bijective edge carriers, even when the inverse is not definitionally simple.
For the proof, the defined maps and incidence equations are substituted.
\lean{Copying/ParametricityPrivateResidualCover}{12}{22}{Parametricity.PrivateRead.ResidualNormal.read_cover}
\par\penalty0

The tail covers reader assertion is the compatibility law for the indicated component of this construction. It relates the actual finite coordinates or certificates, with the hypotheses stated in the linked declaration.
For the proof, constructor and boundary cases are separated explicitly.
\lean{Copying/ParametricityPrivateResidualCover}{24}{26}{Parametricity.PrivateRead.ResidualNormal.tail_covers_reader}
\par\penalty0

The tail covers instance assertion is the compatibility law for the indicated component of this construction. It relates the actual finite coordinates or certificates, with the hypotheses stated in the linked declaration.
For the proof, constructor and boundary cases are separated explicitly.
\lean{Copying/ParametricityPrivateResidualCover}{28}{30}{Parametricity.PrivateRead.ResidualNormal.tail_covers_instance}
\par\penalty0

The tail covers atom assertion is the compatibility law for the indicated component of this construction. It relates the actual finite coordinates or certificates, with the hypotheses stated in the linked declaration.
For the proof, constructor and boundary cases are separated explicitly.
\lean{Copying/ParametricityPrivateResidualCover}{32}{34}{Parametricity.PrivateRead.ResidualNormal.tail_covers_atom}
\par\penalty0

The tail covers store assertion is the compatibility law for the indicated component of this construction. It relates the actual finite coordinates or certificates, with the hypotheses stated in the linked declaration.
For the proof, constructor and boundary cases are separated explicitly.
\lean{Copying/ParametricityPrivateResidualCover}{36}{38}{Parametricity.PrivateRead.ResidualNormal.tail_covers_store}
\par\penalty0

The ready private-cell graph consists precisely of the matched read, the retained reference branch, and the two untouched weakening leaves.
For the proof, constructor and boundary cases are separated explicitly; the remaining natural-number index bounds follow by linear arithmetic; the defined maps and incidence equations are substituted.
\lean{Copying/ParametricityPrivateResidualCover}{42}{94}{Parametricity.PrivateRead.ResidualNormal.original_cover}
\par\penalty0

The actual computed cut keeps every unmatched old edge and exactly the new returned numeral, residual atom tree, store and stored numeral.
For the proof, constructor and boundary cases are separated explicitly; the defined maps and incidence equations are substituted.
\lean{Copying/ParametricityPrivateResidualCover}{98}{124}{Parametricity.PrivateRead.ResidualNormal.corrected_cover}
\par\penalty0

The full covers weak assertion is the compatibility law for the indicated component of this construction. It relates the actual finite coordinates or certificates, with the hypotheses stated in the linked declaration.
For the proof, constructor and boundary cases are separated explicitly; the defined maps and incidence equations are substituted.
\lean{Copying/ParametricityPrivateResidualCover}{126}{152}{Parametricity.PrivateRead.ResidualNormal.full_covers_weak}
\par\penalty0

The full covers tree assertion is the compatibility law for the indicated component of this construction. It relates the actual finite coordinates or certificates, with the hypotheses stated in the linked declaration.
For the proof, the defined maps and incidence equations are substituted.
\lean{Copying/ParametricityPrivateResidualCover}{154}{171}{Parametricity.PrivateRead.ResidualNormal.full_covers_tree}
\par\penalty0

The covers reference assertion is the compatibility law for the indicated component of this construction. It relates the actual finite coordinates or certificates, with the hypotheses stated in the linked declaration.
For the proof, constructor and boundary cases are separated explicitly; the defined maps and incidence equations are substituted.
\lean{Copying/ParametricityPrivateResidualCover}{173}{175}{Parametricity.PrivateRead.ResidualNormal.covers_reference}
\par\penalty0

The covers atom assertion is the compatibility law for the indicated component of this construction. It relates the actual finite coordinates or certificates, with the hypotheses stated in the linked declaration.
For the proof, constructor and boundary cases are separated explicitly; the defined maps and incidence equations are substituted.
\lean{Copying/ParametricityPrivateResidualCover}{177}{179}{Parametricity.PrivateRead.ResidualNormal.covers_atom}
\par\penalty0

The covers numeral assertion is the compatibility law for the indicated component of this construction. It relates the actual finite coordinates or certificates, with the hypotheses stated in the linked declaration.
For the proof, constructor and boundary cases are separated explicitly; the defined maps and incidence equations are substituted.
\lean{Copying/ParametricityPrivateResidualCover}{181}{183}{Parametricity.PrivateRead.ResidualNormal.covers_numeral}
\par\penalty0

The covers cell assertion is the compatibility law for the indicated component of this construction. It relates the actual finite coordinates or certificates, with the hypotheses stated in the linked declaration.
For the proof, constructor and boundary cases are separated explicitly; the defined maps and incidence equations are substituted.
\lean{Copying/ParametricityPrivateResidualCover}{185}{187}{Parametricity.PrivateRead.ResidualNormal.covers_cell}
\par\penalty0

The cover assertion is the compatibility law for the indicated component of this construction. It relates the actual finite coordinates or certificates, with the hypotheses stated in the linked declaration.
For the proof, constructor and boundary cases are separated explicitly; the defined maps and incidence equations are substituted.
\lean{Copying/ParametricityPrivateResidualCover}{189}{199}{Parametricity.PrivateRead.ResidualNormal.cover}
\par\penalty0

The original input assertion computes the indicated ordered boundary coordinate in the assembled graph.
For the proof, constructor and boundary cases are separated explicitly; the remaining natural-number index bounds follow by linear arithmetic; the defined maps and incidence equations are substituted.
\lean{Copying/ParametricityPrivateResidualCover}{201}{222}{Parametricity.PrivateRead.ResidualNormal.original_input}
\par\penalty0

The reference input assertion computes the indicated ordered boundary coordinate in the assembled graph.
For the proof, the remaining natural-number index bounds follow by linear arithmetic.
\lean{Copying/ParametricityPrivateResidualCover}{224}{227}{Parametricity.PrivateRead.ResidualNormal.reference_input}
\par\penalty0

The numeral input assertion computes the indicated ordered boundary coordinate in the assembled graph.
For the proof, the defined maps and incidence equations are substituted.
\lean{Copying/ParametricityPrivateResidualCover}{229}{231}{Parametricity.PrivateRead.ResidualNormal.numeral_input}
\par\penalty0

The input assertion computes the indicated ordered boundary coordinate in the assembled graph.
For the proof, constructor and boundary cases are separated explicitly; the remaining natural-number index bounds follow by linear arithmetic; the defined maps and incidence equations are substituted.
\lean{Copying/ParametricityPrivateResidualCover}{233}{262}{Parametricity.PrivateRead.ResidualNormal.input}
\par\penalty0

Exact graph reconstruction: the only active component is the returned numeral; every other original term input still owns the retained private cell.
\lean{Copying/ParametricityPrivateResidualCover}{266}{267}{Parametricity.PrivateRead.ResidualNormal.iso}
\par\penalty0

\subsection{Parametricity Private Residual Normal}
Reassemble the private read's residual from its actual complementary reference branch, returned numeral, and whole residual atom tree.

This defines the reference data for the construction just described. Boundary positions are kept distinct from newly inserted internal vertices.
\lean{Copying/ParametricityPrivateResidualNormal}{11}{11}{Parametricity.PrivateRead.ResidualNormal.reference}
\par\penalty0

The reader Copy map is a complete occurrence, preserving original labels and every ordered source and target attachment.
\lean{Copying/ParametricityPrivateResidualNormal}{13}{16}{Parametricity.PrivateRead.ResidualNormal.readerCopy}
\par\penalty0

The original map is a complete occurrence, preserving original labels and every ordered source and target attachment.
\lean{Copying/ParametricityPrivateResidualNormal}{18}{21}{Parametricity.PrivateRead.ResidualNormal.original}
\par\penalty0

The read disjoint assertion excludes a shared edge or coordinate between the indicated independently inserted components.
For the proof, constructor and boundary cases are separated explicitly; the remaining natural-number index bounds follow by linear arithmetic; the defined maps and incidence equations are substituted.
\lean{Copying/ParametricityPrivateResidualNormal}{23}{48}{Parametricity.PrivateRead.ResidualNormal.read_disjoint}
\par\penalty0

The disjoint assertion excludes a shared edge or coordinate between the indicated independently inserted components.
For the proof, constructor and boundary cases are separated explicitly.
\lean{Copying/ParametricityPrivateResidualNormal}{50}{54}{Parametricity.PrivateRead.ResidualNormal.disjoint}
\par\penalty0

The reference Occurrence map is a complete occurrence, preserving original labels and every ordered source and target attachment.
\lean{Copying/ParametricityPrivateResidualNormal}{56}{57}{Parametricity.PrivateRead.ResidualNormal.referenceOccurrence}
\par\penalty0

The numeral Occurrence map is a complete occurrence, preserving original labels and every ordered source and target attachment.
\lean{Copying/ParametricityPrivateResidualNormal}{59}{64}{Parametricity.PrivateRead.ResidualNormal.numeralOccurrence}
\par\penalty0

The cell Occurrence map is a complete occurrence, preserving original labels and every ordered source and target attachment.
\lean{Copying/ParametricityPrivateResidualNormal}{66}{67}{Parametricity.PrivateRead.ResidualNormal.cellOccurrence}
\par\penalty0

The reference atom assertion is the compatibility law for the indicated component of this construction. It relates the actual finite coordinates or certificates, with the hypotheses stated in the linked declaration.
For the proof, constructor and boundary cases are separated explicitly; the defined maps and incidence equations are substituted.
\lean{Copying/ParametricityPrivateResidualNormal}{69}{73}{Parametricity.PrivateRead.ResidualNormal.reference_atom}
\par\penalty0

The atom cell assertion is the compatibility law for the indicated component of this construction. It relates the actual finite coordinates or certificates, with the hypotheses stated in the linked declaration.
For the proof, the defined maps and incidence equations are substituted.
\lean{Copying/ParametricityPrivateResidualNormal}{75}{79}{Parametricity.PrivateRead.ResidualNormal.atom_cell}
\par\penalty0

The reference labels assertion verifies the indicated label compatibility or label exclusion at the original coordinate.
For the proof, constructor and boundary cases are separated explicitly; the defined maps and incidence equations are substituted.
\lean{Copying/ParametricityPrivateResidualNormal}{81}{87}{Parametricity.PrivateRead.ResidualNormal.reference_labels}
\par\penalty0

The reference atom disjoint assertion excludes a shared edge or coordinate between the indicated independently inserted components.
For the proof, constructor and boundary cases are separated explicitly; the defined maps and incidence equations are substituted.
\lean{Copying/ParametricityPrivateResidualNormal}{89}{99}{ResidualNormal.reference_atom_disjoint}
\par\penalty0

The atom cell disjoint assertion excludes a shared edge or coordinate between the indicated independently inserted components.
For the proof, constructor and boundary cases are separated explicitly; the defined maps and incidence equations are substituted.
\lean{Copying/ParametricityPrivateResidualNormal}{101}{115}{Parametricity.PrivateRead.ResidualNormal.atom_cell_disjoint}
\par\penalty0

The reference cell disjoint assertion excludes a shared edge or coordinate between the indicated independently inserted components.
For the proof, constructor and boundary cases are separated explicitly.
\lean{Copying/ParametricityPrivateResidualNormal}{117}{119}{ResidualNormal.reference_cell_disjoint}
\par\penalty0

This defines the atom Cell data for the construction just described. Boundary positions are kept distinct from newly inserted internal vertices.
\lean{Copying/ParametricityPrivateResidualNormal}{121}{121}{Parametricity.PrivateRead.ResidualNormal.atomCell}
\par\penalty0

The atom Cell Occurrence map is a complete occurrence, preserving original labels and every ordered source and target attachment.
\lean{Copying/ParametricityPrivateResidualNormal}{123}{126}{Parametricity.PrivateRead.ResidualNormal.atomCellOccurrence}
\par\penalty0

This defines the stored data for the construction just described. Boundary positions are kept distinct from newly inserted internal vertices.
\lean{Copying/ParametricityPrivateResidualNormal}{128}{128}{Parametricity.PrivateRead.ResidualNormal.stored}
\par\penalty0

The stored Occurrence map is a complete occurrence, preserving original labels and every ordered source and target attachment.
\lean{Copying/ParametricityPrivateResidualNormal}{130}{142}{Parametricity.PrivateRead.ResidualNormal.storedOccurrence}
\par\penalty0

The numeral reference disjoint assertion excludes a shared edge or coordinate between the indicated independently inserted components.
For the proof, constructor and boundary cases are separated explicitly.
\lean{Copying/ParametricityPrivateResidualNormal}{144}{146}{ResidualNormal.numeral_reference_disjoint}
\par\penalty0

The numeral atom disjoint assertion excludes a shared edge or coordinate between the indicated independently inserted components.
For the proof, constructor and boundary cases are separated explicitly; the defined maps and incidence equations are substituted.
\lean{Copying/ParametricityPrivateResidualNormal}{148}{157}{ResidualNormal.numeral_atom_disjoint}
\par\penalty0

The numeral cell disjoint assertion excludes a shared edge or coordinate between the indicated independently inserted components.
For the proof, constructor and boundary cases are separated explicitly.
\lean{Copying/ParametricityPrivateResidualNormal}{159}{165}{ResidualNormal.numeral_cell_disjoint}
\par\penalty0

The numeral stored disjoint assertion excludes a shared edge or coordinate between the indicated independently inserted components.
For the proof, constructor and boundary cases are separated explicitly.
\lean{Copying/ParametricityPrivateResidualNormal}{167}{174}{ResidualNormal.numeral_stored_disjoint}
\par\penalty0

This defines the normal data for the construction just described. Boundary positions are kept distinct from newly inserted internal vertices.
\lean{Copying/ParametricityPrivateResidualNormal}{176}{178}{Parametricity.PrivateRead.ResidualNormal.normal}
\par\penalty0

The occurrence map is a complete occurrence, preserving original labels and every ordered source and target attachment.
\lean{Copying/ParametricityPrivateResidualNormal}{180}{191}{Parametricity.PrivateRead.ResidualNormal.occurrence}
\par\penalty0

\subsection{Parametricity Private Source}
The entire original private term tree and allocated cell form one complete occurrence beneath the active reader. Every other term input and the literal canonical atom distributor remain present.

This defines the value data for the construction just described. Boundary positions are kept distinct from newly inserted internal vertices.
\lean{Copying/ParametricityPrivateSource}{11}{11}{Parametricity.PrivateSource.value}
\par\penalty0

This defines the payload data for the construction just described. Boundary positions are kept distinct from newly inserted internal vertices.
\lean{Copying/ParametricityPrivateSource}{12}{12}{Parametricity.PrivateSource.payload}
\par\penalty0

The prefix Occurrence map is a complete occurrence, preserving original labels and every ordered source and target attachment.
\lean{Copying/ParametricityPrivateSource}{14}{15}{Parametricity.PrivateSource.prefixOccurrence}
\par\penalty0

The tail map is a complete occurrence, preserving original labels and every ordered source and target attachment.
\lean{Copying/ParametricityPrivateSource}{16}{17}{Parametricity.PrivateSource.tail}
\par\penalty0

The contents map is a complete occurrence, preserving original labels and every ordered source and target attachment.
\lean{Copying/ParametricityPrivateSource}{18}{21}{Parametricity.PrivateSource.contents}
\par\penalty0

The linked assertion is the compatibility law for the indicated component of this construction. It relates the actual finite coordinates or certificates, with the hypotheses stated in the linked declaration.
For the proof, constructor and boundary cases are separated explicitly; the defined maps and incidence equations are substituted.
\lean{Copying/ParametricityPrivateSource}{23}{28}{Parametricity.PrivateSource.linked}
\par\penalty0

The apart assertion is the compatibility law for the indicated component of this construction. It relates the actual finite coordinates or certificates, with the hypotheses stated in the linked declaration.
For the proof, constructor and boundary cases are separated explicitly.
\lean{Copying/ParametricityPrivateSource}{30}{31}{Parametricity.PrivateSource.apart}
\par\penalty0

The combined map is a complete occurrence, preserving original labels and every ordered source and target attachment.
\lean{Copying/ParametricityPrivateSource}{33}{34}{Parametricity.PrivateSource.combined}
\par\penalty0

Recover the original cell's exact association of instance, atom D1,1, store, and numeral, retaining all four genuine components.
\lean{Copying/ParametricityPrivateSource}{38}{42}{Parametricity.PrivateSource.cell_prefix}
\par\penalty0

The source iso map is an ordered recursive graph isomorphism between the stated original constructions.
\lean{Copying/ParametricityPrivateSource}{44}{47}{Parametricity.PrivateSource.source_iso}
\par\penalty0

The occurrence map is a complete occurrence, preserving original labels and every ordered source and target attachment.
\lean{Copying/ParametricityPrivateSource}{49}{50}{Parametricity.PrivateSource.occurrence}
\par\penalty0

The complete private component attaches to the actual reader output.
For the proof, constructor and boundary cases are separated explicitly; the defined maps and incidence equations are substituted.
\lean{Copying/ParametricityPrivateSource}{53}{59}{Parametricity.PrivateSource.selected_input}
\par\penalty0

The remaining input assertion computes the indicated ordered boundary coordinate in the assembled graph.
For the proof, constructor and boundary cases are separated explicitly; the defined maps and incidence equations are substituted.
\lean{Copying/ParametricityPrivateSource}{61}{78}{Parametricity.PrivateSource.remaining_input}
\par\penalty0

This defines the reader Edge data for the construction just described. Boundary positions are kept distinct from newly inserted internal vertices.
\lean{Copying/ParametricityPrivateSource}{81}{82}{Parametricity.PrivateSource.readerEdge}
\par\penalty0

The covers contents assertion is the compatibility law for the indicated component of this construction. It relates the actual finite coordinates or certificates, with the hypotheses stated in the linked declaration.
For the proof, constructor and boundary cases are separated explicitly.
\lean{Copying/ParametricityPrivateSource}{84}{90}{Parametricity.PrivateSource.covers_contents}
\par\penalty0

The covers tail assertion is the compatibility law for the indicated component of this construction. It relates the actual finite coordinates or certificates, with the hypotheses stated in the linked declaration.
For the proof, constructor and boundary cases are separated explicitly.
\lean{Copying/ParametricityPrivateSource}{92}{98}{Parametricity.PrivateSource.covers_tail}
\par\penalty0

The complete original private component accounts for every edge except the one active dereference itself. No tree sibling or cell edge is omitted.
For the proof, constructor and boundary cases are separated explicitly; the defined maps and incidence equations are substituted.
\lean{Copying/ParametricityPrivateSource}{102}{115}{Parametricity.PrivateSource.cover}
\par\penalty0

\subsection{Parametricity Read Comparison}
Quantitative semantics of the fully constructed private-read corrections. Copying supplies a directed improvement; GC and atom-tree normalization have exact cost. Every intermediate root is derived from actual search reflection.

The completed copying template yields its directed comparison at the actual common nonrewrite focus, without an exact-cost assertion.
For the proof, constructor and boundary cases are separated explicitly; the defined maps and incidence equations are substituted.
\lean{Copying/ParametricityReadComparison}{13}{35}{BoxPullWeakening.PointedWitness.copying_comparison}
\par\penalty0

The hole Free assertion excludes unfilled holes throughout the indicated graph and every nested box.
For the proof, the preceding component lemmas are composed or specialized at the stated parameters.
\lean{Copying/ParametricityReadComparison}{44}{48}{Parametricity.ReadCorrection.holeFree}
\par\penalty0

The actual copying/GC/GC residual chain improves successful runs in the paper's forward direction; converse runs retain the opposite cost inequality.
For the proof, constructor and boundary cases are separated explicitly.
\lean{Copying/ParametricityReadComparison}{52}{71}{Parametricity.ReadCorrection.comparison}
\par\penalty0

Boundary lifting uses the same original interface input on every link, without changing the graph certificates or presuming transitivity of fillings.
\lean{Copying/ParametricityReadComparison}{75}{79}{Parametricity.ReadCorrection.atInput}
\par\penalty0

The actual read correction followed by exact normalization of its remaining atom network to the private cell's original distributor.
\lean{Copying/ParametricityReadComparison}{85}{90}{Parametricity.ReadResidual}
\par\penalty0

The congr map is an ordered recursive graph isomorphism between the stated original constructions.
\lean{Copying/ParametricityReadComparison}{93}{100}{Parametricity.ReadResidual.congr}
\par\penalty0

The cut Input map is a complete occurrence, preserving original labels and every ordered source and target attachment.
\lean{Copying/ParametricityReadComparison}{102}{112}{Parametricity.ReadResidual.cutInput}
\par\penalty0

Every correction at a closed actual focus has the intended quantitative comparison, with the intermediate original search root proved by reflection.
For the proof, constructor and boundary cases are separated explicitly.
\lean{Copying/ParametricityReadComparison}{116}{131}{Parametricity.ReadResidual.comparison}
\par\penalty0

\subsection{Parametricity Right Application Cut}
A closed selected abstraction and its complete untouched family are reconstructed from an actual local beta-source replacement.

This defines the local Source data for the construction just described. Boundary positions are kept distinct from newly inserted internal vertices.
\lean{Copying/ParametricityRightApplicationCut}{17}{17}{Parametricity.RightApplicationCut.localSource}
\par\penalty0

This defines the ordered target list for the construction just described. Boundary positions are kept distinct from newly inserted internal vertices.
\lean{Copying/ParametricityRightApplicationCut}{18}{18}{Parametricity.RightApplicationCut.target}
\par\penalty0

The inserted map is a complete occurrence, preserving original labels and every ordered source and target attachment.
\lean{Copying/ParametricityRightApplicationCut}{19}{19}{Parametricity.RightApplicationCut.inserted}
\par\penalty0

This defines the head data for the construction just described. Boundary positions are kept distinct from newly inserted internal vertices.
\lean{Copying/ParametricityRightApplicationCut}{20}{21}{Parametricity.RightApplicationCut.head}
\par\penalty0

This defines the arguments data for the construction just described. Boundary positions are kept distinct from newly inserted internal vertices.
\lean{Copying/ParametricityRightApplicationCut}{22}{23}{Parametricity.RightApplicationCut.arguments}
\par\penalty0

This defines the selected data for the construction just described. Boundary positions are kept distinct from newly inserted internal vertices.
\lean{Copying/ParametricityRightApplicationCut}{24}{24}{Parametricity.RightApplicationCut.selected}
\par\penalty0

This defines the argument data for the construction just described. Boundary positions are kept distinct from newly inserted internal vertices.
\lean{Copying/ParametricityRightApplicationCut}{25}{25}{Parametricity.RightApplicationCut.argument}
\par\penalty0

This defines the remaining data for the construction just described. Boundary positions are kept distinct from newly inserted internal vertices.
\lean{Copying/ParametricityRightApplicationCut}{26}{26}{Parametricity.RightApplicationCut.remaining}
\par\penalty0

This defines the typed family of inserted components for the construction just described. Boundary positions are kept distinct from newly inserted internal vertices.
\lean{Copying/ParametricityRightApplicationCut}{28}{29}{Parametricity.RightApplicationCut.family}
\par\penalty0

The selected remaining assertion is the compatibility law for the indicated component of this construction. It relates the actual finite coordinates or certificates, with the hypotheses stated in the linked declaration.
For the proof, constructor and boundary cases are separated explicitly.
\lean{Copying/ParametricityRightApplicationCut}{31}{32}{Parametricity.RightApplicationCut.selected_remaining}
\par\penalty0

The family Occurrence map is a complete occurrence, preserving original labels and every ordered source and target attachment.
\lean{Copying/ParametricityRightApplicationCut}{33}{35}{Parametricity.RightApplicationCut.familyOccurrence}
\par\penalty0

The family zero assertion is the compatibility law for the indicated component of this construction. It relates the actual finite coordinates or certificates, with the hypotheses stated in the linked declaration.
For the proof, the defined maps and incidence equations are substituted.
\lean{Copying/ParametricityRightApplicationCut}{37}{41}{Parametricity.RightApplicationCut.family_zero}
\par\penalty0

The head function assertion is the compatibility law for the indicated component of this construction. It relates the actual finite coordinates or certificates, with the hypotheses stated in the linked declaration.
For the proof, the defined maps and incidence equations are substituted.
\lean{Copying/ParametricityRightApplicationCut}{43}{44}{Parametricity.RightApplicationCut.head_function}
\par\penalty0

The head argument assertion is the compatibility law for the indicated component of this construction. It relates the actual finite coordinates or certificates, with the hypotheses stated in the linked declaration.
Both sides are definitionally equal.
\lean{Copying/ParametricityRightApplicationCut}{45}{46}{Parametricity.RightApplicationCut.head_argument}
\par\penalty0

The head family assertion is the compatibility law for the indicated component of this construction. It relates the actual finite coordinates or certificates, with the hypotheses stated in the linked declaration.
For the proof, constructor and boundary cases are separated explicitly.
\lean{Copying/ParametricityRightApplicationCut}{48}{49}{Parametricity.RightApplicationCut.head_family}
\par\penalty0

The head arg assertion is the compatibility law for the indicated component of this construction. It relates the actual finite coordinates or certificates, with the hypotheses stated in the linked declaration.
For the proof, constructor and boundary cases are separated explicitly.
\lean{Copying/ParametricityRightApplicationCut}{50}{51}{Parametricity.RightApplicationCut.head_arg}
\par\penalty0

The family arg assertion is the compatibility law for the indicated component of this construction. It relates the actual finite coordinates or certificates, with the hypotheses stated in the linked declaration.
For the proof, constructor and boundary cases are separated explicitly.
\lean{Copying/ParametricityRightApplicationCut}{52}{53}{Parametricity.RightApplicationCut.family_arg}
\par\penalty0

The occurrence map is a complete occurrence, preserving original labels and every ordered source and target attachment.
\lean{Copying/ParametricityRightApplicationCut}{55}{62}{Parametricity.RightApplicationCut.occurrence}
\par\penalty0

The head input assertion computes the indicated ordered boundary coordinate in the assembled graph.
For the proof, the defined maps and incidence equations are substituted.
\lean{Copying/ParametricityRightApplicationCut}{65}{67}{Parametricity.RightApplicationCut.head_input}
\par\penalty0

The family succ assertion is the compatibility law for the indicated component of this construction. It relates the actual finite coordinates or certificates, with the hypotheses stated in the linked declaration.
For the proof, the remaining natural-number index bounds follow by linear arithmetic; the defined maps and incidence equations are substituted.
\lean{Copying/ParametricityRightApplicationCut}{70}{78}{Parametricity.RightApplicationCut.family_succ}
\par\penalty0

The argument output assertion computes the indicated ordered boundary coordinate in the assembled graph.
For the proof, the defined maps and incidence equations are substituted.
\lean{Copying/ParametricityRightApplicationCut}{81}{94}{Parametricity.RightApplicationCut.argument_output}
\par\penalty0

The input assertion computes the indicated ordered boundary coordinate in the assembled graph.
For the proof, constructor and boundary cases are separated explicitly; the remaining natural-number index bounds follow by linear arithmetic; the defined maps and incidence equations are substituted.
\lean{Copying/ParametricityRightApplicationCut}{97}{106}{Parametricity.RightApplicationCut.input}
\par\penalty0

The output assertion computes the indicated ordered boundary coordinate in the assembled graph.
For the proof, the preceding component lemmas are composed or specialized at the stated parameters.
\lean{Copying/ParametricityRightApplicationCut}{109}{112}{Parametricity.RightApplicationCut.output}
\par\penalty0

The covers assertion states that every target coordinate has an original preimage under the indicated map.
For the proof, constructor and boundary cases are separated explicitly; the defined maps and incidence equations are substituted.
\lean{Copying/ParametricityRightApplicationCut}{114}{136}{Parametricity.RightApplicationCut.covers}
\par\penalty0

The actual computed predecessor is exactly the ordinary application to its original closed lambda and untouched family, not a conditional residual.
\lean{Copying/ParametricityRightApplicationCut}{140}{143}{Parametricity.RightApplicationCut.iso}
\par\penalty0

\subsection{Parametricity Serial Association}
Serial reassociation preserves every ordered boundary and recursive edge label. This is a graph isomorphism, not an observational equation.

This defines the first data for the construction just described. Boundary positions are kept distinct from newly inserted internal vertices.
\lean{Copying/ParametricitySerialAssociation}{11}{11}{Serial.Association.first}
\par\penalty0

This defines the second data for the construction just described. Boundary positions are kept distinct from newly inserted internal vertices.
\lean{Copying/ParametricitySerialAssociation}{12}{12}{Serial.Association.second}
\par\penalty0

This defines the third data for the construction just described. Boundary positions are kept distinct from newly inserted internal vertices.
\lean{Copying/ParametricitySerialAssociation}{13}{13}{Serial.Association.third}
\par\penalty0

The first boundary assertion is the compatibility law for the indicated component of this construction. It relates the actual finite coordinates or certificates, with the hypotheses stated in the linked declaration.
For the proof, the defined maps and incidence equations are substituted.
\lean{Copying/ParametricitySerialAssociation}{15}{20}{Serial.Association.first_boundary}
\par\penalty0

The first disjoint assertion excludes a shared edge or coordinate between the indicated independently inserted components.
For the proof, constructor and boundary cases are separated explicitly.
\lean{Copying/ParametricitySerialAssociation}{22}{23}{Serial.Association.first_disjoint}
\par\penalty0

The pair Occurrence map is a complete occurrence, preserving original labels and every ordered source and target attachment.
\lean{Copying/ParametricitySerialAssociation}{25}{26}{Serial.Association.pairOccurrence}
\par\penalty0

The second boundary assertion is the compatibility law for the indicated component of this construction. It relates the actual finite coordinates or certificates, with the hypotheses stated in the linked declaration.
For the proof, the defined maps and incidence equations are substituted.
\lean{Copying/ParametricitySerialAssociation}{28}{33}{Serial.Association.second_boundary}
\par\penalty0

The second disjoint assertion excludes a shared edge or coordinate between the indicated independently inserted components.
For the proof, constructor and boundary cases are separated explicitly.
\lean{Copying/ParametricitySerialAssociation}{35}{39}{Serial.Association.second_disjoint}
\par\penalty0

The occurrence map is a complete occurrence, preserving original labels and every ordered source and target attachment.
\lean{Copying/ParametricitySerialAssociation}{41}{43}{Serial.Association.occurrence}
\par\penalty0

The input assertion computes the indicated ordered boundary coordinate in the assembled graph.
For the proof, the defined maps and incidence equations are substituted.
\lean{Copying/ParametricitySerialAssociation}{45}{48}{Serial.Association.input}
\par\penalty0

The output assertion computes the indicated ordered boundary coordinate in the assembled graph.
For the proof, the defined maps and incidence equations are substituted.
\lean{Copying/ParametricitySerialAssociation}{50}{54}{Serial.Association.output}
\par\penalty0

The cover assertion is the compatibility law for the indicated component of this construction. It relates the actual finite coordinates or certificates, with the hypotheses stated in the linked declaration.
For the proof, constructor and boundary cases are separated explicitly; the defined maps and incidence equations are substituted.
\lean{Copying/ParametricitySerialAssociation}{56}{62}{Serial.Association.cover}
\par\penalty0

The iso map is an ordered recursive graph isomorphism between the stated original constructions.
\lean{Copying/ParametricitySerialAssociation}{64}{65}{Serial.Association.iso}
\par\penalty0

\subsection{Parametricity Small Family Application}
The actual smaller-abstraction family overlap: one original beta plus eight original private-store transitions on the left, and one original beta on the right. Every unused stable argument output remains in both graphs.

This defines the discarded data for the construction just described. Boundary positions are kept distinct from newly inserted internal vertices.
\lean{Copying/ParametricitySmallFamilyApplication}{14}{14}{Parametricity.SmallFamilyApplication.discarded}
\par\penalty0

This defines the opened data for the construction just described. Boundary positions are kept distinct from newly inserted internal vertices.
\lean{Copying/ParametricitySmallFamilyApplication}{16}{16}{Parametricity.SmallFamilyApplication.opened}
\par\penalty0

The private Occurrence map is a complete occurrence, preserving original labels and every ordered source and target attachment.
\lean{Copying/ParametricitySmallFamilyApplication}{17}{18}{Parametricity.SmallFamilyApplication.privateOccurrence}
\par\penalty0

The argument Occurrence map is a complete occurrence, preserving original labels and every ordered source and target attachment.
\lean{Copying/ParametricitySmallFamilyApplication}{19}{20}{Parametricity.SmallFamilyApplication.argumentOccurrence}
\par\penalty0

The raw Reader map is a complete occurrence, preserving original labels and every ordered source and target attachment.
\lean{Copying/ParametricitySmallFamilyApplication}{22}{24}{Parametricity.SmallFamilyApplication.rawReader}
\par\penalty0

The reader map is a complete occurrence, preserving original labels and every ordered source and target attachment.
\lean{Copying/ParametricitySmallFamilyApplication}{26}{28}{Parametricity.SmallFamilyApplication.reader}
\par\penalty0

The raw Reader input assertion computes the indicated ordered boundary coordinate in the assembled graph.
For the proof, the defined maps and incidence equations are substituted.
\lean{Copying/ParametricitySmallFamilyApplication}{30}{36}{Parametricity.SmallFamilyApplication.rawReader_input}
\par\penalty0

The reader input assertion computes the indicated ordered boundary coordinate in the assembled graph.
For the proof, the defined maps and incidence equations are substituted.
\lean{Copying/ParametricitySmallFamilyApplication}{38}{43}{Parametricity.SmallFamilyApplication.reader_input}
\par\penalty0

The body disjoint assertion excludes a shared edge or coordinate between the indicated independently inserted components.
For the proof, constructor and boundary cases are separated explicitly.
\lean{Copying/ParametricitySmallFamilyApplication}{45}{47}{Parametricity.SmallFamilyApplication.body_disjoint}
\par\penalty0

The body Occurrence map is a complete occurrence, preserving original labels and every ordered source and target attachment.
\lean{Copying/ParametricitySmallFamilyApplication}{49}{51}{Parametricity.SmallFamilyApplication.bodyOccurrence}
\par\penalty0

The body input assertion computes the indicated ordered boundary coordinate in the assembled graph.
For the proof, the defined maps and incidence equations are substituted.
\lean{Copying/ParametricitySmallFamilyApplication}{53}{57}{Parametricity.SmallFamilyApplication.body_input}
\par\penalty0

The predecessor is obtained by inserting the original complete beta redex into the actual read/family/private-cell complement.
\lean{Copying/ParametricitySmallFamilyApplication}{61}{62}{Parametricity.SmallFamilyApplication.before}
\par\penalty0

This defines the token data for the construction just described. Boundary positions are kept distinct from newly inserted internal vertices.
\lean{Copying/ParametricitySmallFamilyApplication}{63}{64}{Parametricity.SmallFamilyApplication.token}
\par\penalty0

The beta assertion is the compatibility law for the indicated component of this construction. It relates the actual finite coordinates or certificates, with the hypotheses stated in the linked declaration.
For the proof, the defined maps and incidence equations are substituted.
\lean{Copying/ParametricitySmallFamilyApplication}{66}{72}{Parametricity.SmallFamilyApplication.beta}
\par\penalty0

Literal paired residuals, including every ignored argument resource.
\lean{Copying/ParametricitySmallFamilyApplication}{75}{75}{Parametricity.SmallFamilyApplication.corrected}
\par\penalty0

This defines the counterpart data for the construction just described. Boundary positions are kept distinct from newly inserted internal vertices.
\lean{Copying/ParametricitySmallFamilyApplication}{76}{76}{Parametricity.SmallFamilyApplication.counterpart}
\par\penalty0

The paired data give an actual recursive paired filling, with its common context and endpoint representation maps.
\lean{Copying/ParametricitySmallFamilyApplication}{78}{81}{Parametricity.SmallFamilyApplication.paired}
\par\penalty0

This defines the pointed data for the construction just described. Boundary positions are kept distinct from newly inserted internal vertices.
\lean{Copying/ParametricitySmallFamilyApplication}{83}{85}{Parametricity.SmallFamilyApplication.pointed}
\par\penalty0

The opened focus Free assertion excludes focus labels throughout the indicated graph, including recursively inside boxes.
For the proof, constructor and boundary cases are separated explicitly.
\lean{Copying/ParametricitySmallFamilyApplication}{88}{92}{Parametricity.SmallFamilyApplication.opened_focusFree}
\par\penalty0

The read after input assertion computes the indicated ordered boundary coordinate in the assembled graph.
For the proof, the defined maps and incidence equations are substituted.
\lean{Copying/ParametricitySmallFamilyApplication}{94}{96}{Parametricity.SmallFamilyApplication.read_after_input}
\par\penalty0

The correction map is a complete occurrence, preserving original labels and every ordered source and target attachment.
\lean{Copying/ParametricitySmallFamilyApplication}{98}{117}{Parametricity.SmallFamilyApplication.correction}
\par\penalty0

All nine original transitions are preserved after reattaching every remaining closure and every output of the arbitrary stable argument.
For the proof, the preceding component lemmas are composed or specialized at the stated parameters.
\lean{Copying/ParametricitySmallFamilyApplication}{121}{129}{Parametricity.SmallFamilyApplication.run}
\par\penalty0

The numeral Occurrence map is a complete occurrence, preserving original labels and every ordered source and target attachment.
\lean{Copying/ParametricitySmallFamilyApplication}{131}{135}{Parametricity.SmallFamilyApplication.numeralOccurrence}
\par\penalty0

The right Argument map is a complete occurrence, preserving original labels and every ordered source and target attachment.
\lean{Copying/ParametricitySmallFamilyApplication}{136}{138}{Parametricity.SmallFamilyApplication.rightArgument}
\par\penalty0

The right disjoint assertion excludes a shared edge or coordinate between the indicated independently inserted components.
For the proof, constructor and boundary cases are separated explicitly.
\lean{Copying/ParametricitySmallFamilyApplication}{140}{143}{Parametricity.SmallFamilyApplication.right_disjoint}
\par\penalty0

The right Body map is a complete occurrence, preserving original labels and every ordered source and target attachment.
\lean{Copying/ParametricitySmallFamilyApplication}{145}{149}{Parametricity.SmallFamilyApplication.rightBody}
\par\penalty0

The right Body input assertion computes the indicated ordered boundary coordinate in the assembled graph.
For the proof, the defined maps and incidence equations are substituted.
\lean{Copying/ParametricitySmallFamilyApplication}{151}{156}{Parametricity.SmallFamilyApplication.rightBody_input}
\par\penalty0

The counterpart retains the same stable argument and the entire remaining right family, and uses the original constant-body beta redex.
\lean{Copying/ParametricitySmallFamilyApplication}{160}{162}{Parametricity.SmallFamilyApplication.rightBefore}
\par\penalty0

This defines the right Token data for the construction just described. Boundary positions are kept distinct from newly inserted internal vertices.
\lean{Copying/ParametricitySmallFamilyApplication}{163}{165}{Parametricity.SmallFamilyApplication.rightToken}
\par\penalty0

The right beta assertion is the compatibility law for the indicated component of this construction. It relates the actual finite coordinates or certificates, with the hypotheses stated in the linked declaration.
For the proof, the defined maps and incidence equations are substituted.
\lean{Copying/ParametricitySmallFamilyApplication}{167}{175}{Parametricity.SmallFamilyApplication.right_beta}
\par\penalty0

The actual 9-versus-1 family diagram. Its residual is constructed from the original copying template, two exact garbage collections, exact sharing, and an actual smaller Figure20 family. No termination of V is presumed.
For the proof, the defined maps and incidence equations are substituted.
\lean{Copying/ParametricitySmallFamilyApplication}{180}{197}{Parametricity.SmallFamilyApplication.nine_vs_one}
\par\penalty0

\subsection{Parametricity Small Family Diagram}
The arbitrary-family smaller-abstraction diagram starts from the actual original application context around a genuine Figure20 template filling. Its nine-versus-one prefixes and all residual certificates are constructed.

The original application head observes the selected small abstraction; the complete remaining left family and original private cell stay inside.
\lean{Copying/ParametricitySmallFamilyDiagram}{16}{17}{Parametricity.SmallFamilyDiagram.left}
\par\penalty0

The identical original application context contains the matching right family, with every output of the same stable argument preserved.
\lean{Copying/ParametricitySmallFamilyDiagram}{21}{22}{Parametricity.SmallFamilyDiagram.right}
\par\penalty0

The sources themselves are certified instances of the original template, not a new relation defined by the desired execution comparison.
\lean{Copying/ParametricitySmallFamilyDiagram}{26}{29}{Parametricity.SmallFamilyDiagram.sourceWitness}
\par\penalty0

The left token assertion is the compatibility law for the indicated component of this construction. It relates the actual finite coordinates or certificates, with the hypotheses stated in the linked declaration.
For the proof, the preceding component lemmas are composed or specialized at the stated parameters.
\lean{Copying/ParametricitySmallFamilyDiagram}{31}{32}{Parametricity.SmallFamilyDiagram.left_token}
\par\penalty0

The right token assertion is the compatibility law for the indicated component of this construction. It relates the actual finite coordinates or certificates, with the hypotheses stated in the linked declaration.
For the proof, the defined maps and incidence equations are substituted.
\lean{Copying/ParametricitySmallFamilyDiagram}{33}{34}{Parametricity.SmallFamilyDiagram.right_token}
\par\penalty0

The right prefix is the actual constant-body beta rule inside its unchanged application and remaining family context.
For the proof, the defined maps and incidence equations are substituted.
\lean{Copying/ParametricitySmallFamilyDiagram}{38}{44}{Parametricity.SmallFamilyDiagram.right_beta}
\par\penalty0

The complete original first-small family case: nine actual left transitions versus one actual right transition, followed by the independently proved copying/GC/GC/sharing corrections and a genuine remaining-family parametricity specimen. Stable arguments are retained and never evaluated.
For the proof, the defined maps and incidence equations are substituted.
\lean{Copying/ParametricitySmallFamilyDiagram}{50}{71}{Parametricity.SmallFamilyDiagram.nine_vs_one}
\par\penalty0

\subsection{Parametricity Small Family Read}
The complete private read is reattached to both actual remaining Figure20 families. Its terminal correction contains a genuine parametricity specimen with the selected term input weakened, not an assumed residual law.

This defines the before data for the construction just described. Boundary positions are kept distinct from newly inserted internal vertices.
\lean{Copying/ParametricitySmallFamilyRead}{15}{16}{Parametricity.SmallFamilyRead.before}
\par\penalty0

This defines the after data for the construction just described. Boundary positions are kept distinct from newly inserted internal vertices.
\lean{Copying/ParametricitySmallFamilyRead}{18}{19}{Parametricity.SmallFamilyRead.after}
\par\penalty0

All unchanged left closures are reattached in their original order.
\lean{Copying/ParametricitySmallFamilyRead}{22}{23}{Parametricity.SmallFamilyRead.corrected}
\par\penalty0

The same returned numeral precedes the unchanged ordered right family.
\lean{Copying/ParametricitySmallFamilyRead}{26}{28}{Parametricity.SmallFamilyRead.counterpart}
\par\penalty0

This defines the residual data for the construction just described. Boundary positions are kept distinct from newly inserted internal vertices.
\lean{Copying/ParametricitySmallFamilyRead}{30}{34}{Parametricity.SmallFamilyRead.residual}
\par\penalty0

The remaining left private family is literally the original Figure20 family for the actual residual tree. Serial association only changes graph representation, not the original sharing, store encoding, or translation.
\lean{Copying/ParametricitySmallFamilyRead}{39}{49}{Parametricity.SmallFamilyRead.residual_family_iso}
\par\penalty0

A genuine remaining-family specimen. Both multiplicities are arbitrary, including zero; its tree is the actual original tree with only the selected input closed by the copying rule's literal weakening.
\lean{Copying/ParametricitySmallFamilyRead}{54}{60}{Parametricity.SmallFamilyRead.paired}
\par\penalty0

This defines the pointed data for the construction just described. Boundary positions are kept distinct from newly inserted internal vertices.
\lean{Copying/ParametricitySmallFamilyRead}{62}{64}{Parametricity.SmallFamilyRead.pointed}
\par\penalty0

Exactly eight original transitions after beta, followed by the fully constructed directed/exact correction and original remaining-family specimen.
For the proof, the defined maps and incidence equations are substituted.
\lean{Copying/ParametricitySmallFamilyRead}{68}{77}{Parametricity.SmallFamilyRead.factorization}
\par\penalty0

\subsection{Parametricity Small Family Source}
The computed smaller-beta predecessor contains the entire original isolated Figure20 family, its arbitrary original private tree, and its literal allocated cell. The only external components are application and argument.

This defines the old-frame port map for the construction just described. Boundary positions are kept distinct from newly inserted internal vertices.
\lean{Copying/ParametricitySmallFamilySource}{14}{14}{Parametricity.SmallFamilySource.host}
\par\penalty0

This defines the local Source data for the construction just described. Boundary positions are kept distinct from newly inserted internal vertices.
\lean{Copying/ParametricitySmallFamilySource}{15}{15}{Parametricity.SmallFamilySource.localSource}
\par\penalty0

The inserted map is a complete occurrence, preserving original labels and every ordered source and target attachment.
\lean{Copying/ParametricitySmallFamilySource}{16}{17}{Parametricity.SmallFamilySource.inserted}
\par\penalty0

This defines the head data for the construction just described. Boundary positions are kept distinct from newly inserted internal vertices.
\lean{Copying/ParametricitySmallFamilySource}{18}{20}{Parametricity.SmallFamilySource.head}
\par\penalty0

This defines the arguments data for the construction just described. Boundary positions are kept distinct from newly inserted internal vertices.
\lean{Copying/ParametricitySmallFamilySource}{21}{23}{Parametricity.SmallFamilySource.arguments}
\par\penalty0

This defines the selected data for the construction just described. Boundary positions are kept distinct from newly inserted internal vertices.
\lean{Copying/ParametricitySmallFamilySource}{24}{25}{Parametricity.SmallFamilySource.selected}
\par\penalty0

This defines the argument data for the construction just described. Boundary positions are kept distinct from newly inserted internal vertices.
\lean{Copying/ParametricitySmallFamilySource}{26}{27}{Parametricity.SmallFamilySource.argument}
\par\penalty0

The old Family map is a complete occurrence, preserving original labels and every ordered source and target attachment.
\lean{Copying/ParametricitySmallFamilySource}{29}{31}{Parametricity.SmallFamilySource.oldFamily}
\par\penalty0

The old Cell map is a complete occurrence, preserving original labels and every ordered source and target attachment.
\lean{Copying/ParametricitySmallFamilySource}{32}{36}{Parametricity.SmallFamilySource.oldCell}
\par\penalty0

The family avoids assertion is the compatibility law for the indicated component of this construction. It relates the actual finite coordinates or certificates, with the hypotheses stated in the linked declaration.
For the proof, constructor and boundary cases are separated explicitly.
\lean{Copying/ParametricitySmallFamilySource}{38}{41}{Parametricity.SmallFamilySource.family_avoids}
\par\penalty0

The cell avoids assertion is the compatibility law for the indicated component of this construction. It relates the actual finite coordinates or certificates, with the hypotheses stated in the linked declaration.
For the proof, constructor and boundary cases are separated explicitly; the defined maps and incidence equations are substituted.
\lean{Copying/ParametricitySmallFamilySource}{43}{54}{Parametricity.SmallFamilySource.cell_avoids}
\par\penalty0

The remaining map is a complete occurrence, preserving original labels and every ordered source and target attachment.
\lean{Copying/ParametricitySmallFamilySource}{56}{58}{Parametricity.SmallFamilySource.remaining}
\par\penalty0

The cell map is a complete occurrence, preserving original labels and every ordered source and target attachment.
\lean{Copying/ParametricitySmallFamilySource}{59}{61}{Parametricity.SmallFamilySource.cell}
\par\penalty0

The selected remaining disjoint assertion excludes a shared edge or coordinate between the indicated independently inserted components.
For the proof, constructor and boundary cases are separated explicitly.
\lean{Copying/ParametricitySmallFamilySource}{63}{65}{Parametricity.SmallFamilySource.selected_remaining_disjoint}
\par\penalty0

The family map is a complete occurrence, preserving original labels and every ordered source and target attachment.
\lean{Copying/ParametricitySmallFamilySource}{67}{69}{Parametricity.SmallFamilySource.family}
\par\penalty0

Before beta, the selected lambda's free x is exactly the original private reader attachment reconstructed by the inverse beta cut.
For the proof, the defined maps and incidence equations are substituted.
\lean{Copying/ParametricitySmallFamilySource}{73}{91}{Parametricity.SmallFamilySource.selected_cell}
\par\penalty0

The remaining cell assertion is the compatibility law for the indicated component of this construction. It relates the actual finite coordinates or certificates, with the hypotheses stated in the linked declaration.
For the proof, the defined maps and incidence equations are substituted.
\lean{Copying/ParametricitySmallFamilySource}{93}{106}{Parametricity.SmallFamilySource.remaining_cell}
\par\penalty0

The family cell assertion is the compatibility law for the indicated component of this construction. It relates the actual finite coordinates or certificates, with the hypotheses stated in the linked declaration.
For the proof, constructor and boundary cases are separated explicitly; the remaining natural-number index bounds follow by linear arithmetic; the defined maps and incidence equations are substituted.
\lean{Copying/ParametricitySmallFamilySource}{109}{124}{Parametricity.SmallFamilySource.family_cell}
\par\penalty0

The family cell disjoint assertion excludes a shared edge or coordinate between the indicated independently inserted components.
For the proof, constructor and boundary cases are separated explicitly; the defined maps and incidence equations are substituted.
\lean{Copying/ParametricitySmallFamilySource}{126}{132}{Parametricity.SmallFamilySource.family_cell_disjoint}
\par\penalty0

This defines the closed Graph data for the construction just described. Boundary positions are kept distinct from newly inserted internal vertices.
\lean{Copying/ParametricitySmallFamilySource}{134}{135}{Parametricity.SmallFamilySource.closedGraph}
\par\penalty0

Every original family abstraction, the entire arbitrary T, and the unchanged allocated cell constitute one complete actual source occurrence.
\lean{Copying/ParametricitySmallFamilySource}{139}{141}{Parametricity.SmallFamilySource.closedFamily}
\par\penalty0

The closed zero assertion is the compatibility law for the indicated component of this construction. It relates the actual finite coordinates or certificates, with the hypotheses stated in the linked declaration.
For the proof, the defined maps and incidence equations are substituted.
\lean{Copying/ParametricitySmallFamilySource}{143}{148}{Parametricity.SmallFamilySource.closed_zero}
\par\penalty0

The head function assertion is the compatibility law for the indicated component of this construction. It relates the actual finite coordinates or certificates, with the hypotheses stated in the linked declaration.
For the proof, the defined maps and incidence equations are substituted.
\lean{Copying/ParametricitySmallFamilySource}{150}{153}{Parametricity.SmallFamilySource.head_function}
\par\penalty0

The head argument assertion is the compatibility law for the indicated component of this construction. It relates the actual finite coordinates or certificates, with the hypotheses stated in the linked declaration.
For the proof, the defined maps and incidence equations are substituted.
\lean{Copying/ParametricitySmallFamilySource}{155}{156}{Parametricity.SmallFamilySource.head_argument}
\par\penalty0

The head family disjoint assertion excludes a shared edge or coordinate between the indicated independently inserted components.
For the proof, constructor and boundary cases are separated explicitly.
\lean{Copying/ParametricitySmallFamilySource}{158}{163}{Parametricity.SmallFamilySource.head_family_disjoint}
\par\penalty0

The head argument disjoint assertion excludes a shared edge or coordinate between the indicated independently inserted components.
For the proof, constructor and boundary cases are separated explicitly.
\lean{Copying/ParametricitySmallFamilySource}{165}{166}{Parametricity.SmallFamilySource.head_argument_disjoint}
\par\penalty0

The family argument disjoint assertion excludes a shared edge or coordinate between the indicated independently inserted components.
For the proof, constructor and boundary cases are separated explicitly.
\lean{Copying/ParametricitySmallFamilySource}{168}{173}{Parametricity.SmallFamilySource.family_argument_disjoint}
\par\penalty0

The isolated source family is certified by the actual original parametricity pre-template, after the proved explicit input permutation.
\lean{Copying/ParametricitySmallFamilySource}{177}{181}{Parametricity.SmallFamilySource.familyWitness}
\par\penalty0

\subsection{Parametricity Small Family Source Iso}
Complete reconstruction identifies the computed smaller-beta source with the original application context around its whole Figure20 family.

The head input assertion computes the indicated ordered boundary coordinate in the assembled graph.
For the proof, the defined maps and incidence equations are substituted.
\lean{Copying/ParametricitySmallFamilySourceIso}{13}{16}{Parametricity.SmallFamilySource.head_input}
\par\penalty0

The remaining input assertion computes the indicated ordered boundary coordinate in the assembled graph.
For the proof, the remaining natural-number index bounds follow by linear arithmetic; the defined maps and incidence equations are substituted.
\lean{Copying/ParametricitySmallFamilySourceIso}{18}{24}{Parametricity.SmallFamilySource.remaining_input}
\par\penalty0

The closed input succ assertion computes the indicated ordered boundary coordinate in the assembled graph.
For the proof, the defined maps and incidence equations are substituted.
\lean{Copying/ParametricitySmallFamilySourceIso}{26}{33}{Parametricity.SmallFamilySource.closed_input_succ}
\par\penalty0

The argument output assertion computes the indicated ordered boundary coordinate in the assembled graph.
For the proof, the remaining natural-number index bounds follow by linear arithmetic; the defined maps and incidence equations are substituted.
\lean{Copying/ParametricitySmallFamilySourceIso}{35}{50}{Parametricity.SmallFamilySource.argument_output}
\par\penalty0

The generic original application context is realized by its actual head, whole original private family, and unchanged arbitrary stable argument.
\lean{Copying/ParametricitySmallFamilySourceIso}{54}{58}{Parametricity.SmallFamilySource.sourceOccurrence}
\par\penalty0

The input assertion computes the indicated ordered boundary coordinate in the assembled graph.
For the proof, constructor and boundary cases are separated explicitly; the remaining natural-number index bounds follow by linear arithmetic; the defined maps and incidence equations are substituted.
\lean{Copying/ParametricitySmallFamilySourceIso}{60}{71}{Parametricity.SmallFamilySource.input}
\par\penalty0

The output assertion computes the indicated ordered boundary coordinate in the assembled graph.
For the proof, the defined maps and incidence equations are substituted.
\lean{Copying/ParametricitySmallFamilySourceIso}{73}{75}{Parametricity.SmallFamilySource.output}
\par\penalty0

The covers head assertion is the compatibility law for the indicated component of this construction. It relates the actual finite coordinates or certificates, with the hypotheses stated in the linked declaration.
For the proof, the defined maps and incidence equations are substituted.
\lean{Copying/ParametricitySmallFamilySourceIso}{78}{79}{Parametricity.SmallFamilySource.covers_head}
\par\penalty0

The covers argument assertion is the compatibility law for the indicated component of this construction. It relates the actual finite coordinates or certificates, with the hypotheses stated in the linked declaration.
For the proof, the defined maps and incidence equations are substituted.
\lean{Copying/ParametricitySmallFamilySourceIso}{80}{81}{Parametricity.SmallFamilySource.covers_argument}
\par\penalty0

The covers selected assertion is the compatibility law for the indicated component of this construction. It relates the actual finite coordinates or certificates, with the hypotheses stated in the linked declaration.
For the proof, the defined maps and incidence equations are substituted.
\lean{Copying/ParametricitySmallFamilySourceIso}{82}{84}{Parametricity.SmallFamilySource.covers_selected}
\par\penalty0

The covers remaining assertion is the compatibility law for the indicated component of this construction. It relates the actual finite coordinates or certificates, with the hypotheses stated in the linked declaration.
For the proof, the defined maps and incidence equations are substituted.
\lean{Copying/ParametricitySmallFamilySourceIso}{85}{87}{Parametricity.SmallFamilySource.covers_remaining}
\par\penalty0

The covers cell assertion is the compatibility law for the indicated component of this construction. It relates the actual finite coordinates or certificates, with the hypotheses stated in the linked declaration.
For the proof, the defined maps and incidence equations are substituted.
\lean{Copying/ParametricitySmallFamilySourceIso}{88}{90}{Parametricity.SmallFamilySource.covers_cell}
\par\penalty0

Before the inverse beta cut, only the actual reader and discarded argument belong to the removed component; all other edges belong to the unchanged remaining family or the complete original private cell.
For the proof, constructor and boundary cases are separated explicitly; the defined maps and incidence equations are substituted.
\lean{Copying/ParametricitySmallFamilySourceIso}{95}{114}{Parametricity.SmallFamilySource.opened_cover}
\par\penalty0

Every actual source edge is covered by the three original components; this includes all unused sharing siblings and every argument resource edge.
For the proof, constructor and boundary cases are separated explicitly; the defined maps and incidence equations are substituted.
\lean{Copying/ParametricitySmallFamilySourceIso}{118}{146}{Parametricity.SmallFamilySource.cover}
\par\penalty0

The complete source is the original application context around the entire original Figure20 family, with only equal-interface presentation casts. The map fixes every external input and argument output.
\lean{Copying/ParametricitySmallFamilySourceIso}{151}{154}{Parametricity.SmallFamilySource.iso}
\par\penalty0

\subsection{Parametricity Small Read Residual}
The small-abstraction case retains the original nine transitions and ends, through proved directed/exact corrections, at the family-ready graph. The stable callback argument is not evaluated and none of its outputs vanish.

This defines the full Residual data for the construction just described. Boundary positions are kept distinct from newly inserted internal vertices.
\lean{Copying/ParametricitySmallReadResidual}{13}{15}{Parametricity.PrivateRead.fullResidual}
\par\penalty0

Insert the returned numeral and the unchanged original private cell in the original opened-beta context, retaining every ignored argument resource.
\lean{Copying/ParametricitySmallReadResidual}{24}{26}{Parametricity.PrivateApplication.result}
\par\penalty0

This defines the result Point data for the construction just described. Boundary positions are kept distinct from newly inserted internal vertices.
\lean{Copying/ParametricitySmallReadResidual}{28}{30}{Parametricity.PrivateApplication.resultPoint}
\par\penalty0

The full Residual map is a complete occurrence, preserving original labels and every ordered source and target attachment.
\lean{Copying/ParametricitySmallReadResidual}{32}{42}{Parametricity.PrivateApplication.fullResidual}
\par\penalty0

Exactly nine actual original transitions, followed by the constructed Copying/GC/GC/Sharing chain. This is the full local small-read calculation, including the original private store and arbitrary stable argument resources.
For the proof, the defined maps and incidence equations are substituted.
\lean{Copying/ParametricitySmallReadResidual}{47}{57}{Parametricity.PrivateApplication.nine_transition_residual}
\par\penalty0

\subsection{Parametricity Small Right Source}
Exact source correspondence for the smaller constant-returning family's original beta predecessor, including all remaining closures and V outputs.

The remaining map is a complete occurrence, preserving original labels and every ordered source and target attachment.
\lean{Copying/ParametricitySmallRightSource}{12}{17}{Parametricity.SmallRightSource.remaining}
\par\penalty0

The apart assertion is the compatibility law for the indicated component of this construction. It relates the actual finite coordinates or certificates, with the hypotheses stated in the linked declaration.
For the proof, constructor and boundary cases are separated explicitly.
\lean{Copying/ParametricitySmallRightSource}{19}{21}{Parametricity.SmallRightSource.apart}
\par\penalty0

The remaining input assertion computes the indicated ordered boundary coordinate in the assembled graph.
For the proof, the remaining natural-number index bounds follow by linear arithmetic; the defined maps and incidence equations are substituted.
\lean{Copying/ParametricitySmallRightSource}{23}{25}{Parametricity.SmallRightSource.remaining_input}
\par\penalty0

The body output assertion computes the indicated ordered boundary coordinate in the assembled graph.
For the proof, the defined maps and incidence equations are substituted.
\lean{Copying/ParametricitySmallRightSource}{27}{38}{Parametricity.SmallRightSource.body_output}
\par\penalty0

The local constant/discarded-argument reduct and unchanged remaining family partition every original edge.
For the proof, constructor and boundary cases are separated explicitly; the defined maps and incidence equations are substituted.
\lean{Copying/ParametricitySmallRightSource}{42}{51}{Parametricity.SmallRightSource.cover}
\par\penalty0

The computed predecessor is exactly application to the original isolated small-right family, with the original ordered argument resources.
\lean{Copying/ParametricitySmallRightSource}{55}{60}{Parametricity.SmallRightSource.iso}
\par\penalty0

\subsection{Parametricity Tail Precompose}
Reattach an arbitrary ordered remaining family to all but the active input of a private computation. The zero-member family is included literally; no identity edge, dummy closure, or boundary quotient is inserted.

This defines the Vertex data for the construction just described. Boundary positions are kept distinct from newly inserted internal vertices.
\lean{Copying/ParametricityTailPrecompose}{11}{11}{Parametricity.TailPrecompose.Vertex}
\par\penalty0

This defines the family Input data for the construction just described. Boundary positions are kept distinct from newly inserted internal vertices.
\lean{Copying/ParametricityTailPrecompose}{13}{13}{Parametricity.TailPrecompose.familyInput}
\par\penalty0

This defines the active data for the construction just described. Boundary positions are kept distinct from newly inserted internal vertices.
\lean{Copying/ParametricityTailPrecompose}{14}{14}{Parametricity.TailPrecompose.active}
\par\penalty0

This defines the hole data for the construction just described. Boundary positions are kept distinct from newly inserted internal vertices.
\lean{Copying/ParametricityTailPrecompose}{15}{17}{Parametricity.TailPrecompose.hole}
\par\penalty0

This defines the sources data for the construction just described. Boundary positions are kept distinct from newly inserted internal vertices.
\lean{Copying/ParametricityTailPrecompose}{18}{20}{Parametricity.TailPrecompose.sources}
\par\penalty0

This defines the targets data for the construction just described. Boundary positions are kept distinct from newly inserted internal vertices.
\lean{Copying/ParametricityTailPrecompose}{21}{23}{Parametricity.TailPrecompose.targets}
\par\penalty0

This defines the common graph frame for the construction just described. Boundary positions are kept distinct from newly inserted internal vertices.
\lean{Copying/ParametricityTailPrecompose}{25}{94}{Parametricity.TailPrecompose.frame}
\par\penalty0

This defines the pattern data for the construction just described. Boundary positions are kept distinct from newly inserted internal vertices.
\lean{Copying/ParametricityTailPrecompose}{96}{99}{Parametricity.TailPrecompose.pattern}
\par\penalty0

This defines the filling data for the construction just described. Boundary positions are kept distinct from newly inserted internal vertices.
\lean{Copying/ParametricityTailPrecompose}{101}{106}{Parametricity.TailPrecompose.filling}
\par\penalty0

This defines the finite typed net for the construction just described. Boundary positions are kept distinct from newly inserted internal vertices.
\lean{Copying/ParametricityTailPrecompose}{108}{114}{Parametricity.TailPrecompose.net}
\par\penalty0

This defines the left Insert data for the construction just described. Boundary positions are kept distinct from newly inserted internal vertices.
\lean{Copying/ParametricityTailPrecompose}{116}{117}{Parametricity.TailPrecompose.leftInsert}
\par\penalty0

This defines the right Insert data for the construction just described. Boundary positions are kept distinct from newly inserted internal vertices.
\lean{Copying/ParametricityTailPrecompose}{119}{120}{Parametricity.TailPrecompose.rightInsert}
\par\penalty0

The left Occurrence map is a complete occurrence, preserving original labels and every ordered source and target attachment.
\lean{Copying/ParametricityTailPrecompose}{122}{123}{Parametricity.TailPrecompose.leftOccurrence}
\par\penalty0

The right Occurrence map is a complete occurrence, preserving original labels and every ordered source and target attachment.
\lean{Copying/ParametricityTailPrecompose}{124}{125}{Parametricity.TailPrecompose.rightOccurrence}
\par\penalty0

The left input assertion computes the indicated ordered boundary coordinate in the assembled graph.
For the proof, constructor and boundary cases are separated explicitly; the remaining natural-number index bounds follow by linear arithmetic; the defined maps and incidence equations are substituted.
\lean{Copying/ParametricityTailPrecompose}{127}{132}{Parametricity.TailPrecompose.left_input}
\par\penalty0

The right zero assertion is the compatibility law for the indicated component of this construction. It relates the actual finite coordinates or certificates, with the hypotheses stated in the linked declaration.
For the proof, constructor and boundary cases are separated explicitly; the defined maps and incidence equations are substituted.
\lean{Copying/ParametricityTailPrecompose}{134}{135}{Parametricity.TailPrecompose.right_zero}
\par\penalty0

The right succ assertion is the compatibility law for the indicated component of this construction. It relates the actual finite coordinates or certificates, with the hypotheses stated in the linked declaration.
For the proof, constructor and boundary cases are separated explicitly; the defined maps and incidence equations are substituted.
\lean{Copying/ParametricityTailPrecompose}{137}{144}{Parametricity.TailPrecompose.right_succ}
\par\penalty0

The right output assertion computes the indicated ordered boundary coordinate in the assembled graph.
For the proof, constructor and boundary cases are separated explicitly; the defined maps and incidence equations are substituted.
\lean{Copying/ParametricityTailPrecompose}{146}{150}{Parametricity.TailPrecompose.right_output}
\par\penalty0

The focus Free assertion excludes focus labels throughout the indicated graph, including recursively inside boxes.
For the proof, constructor and boundary cases are separated explicitly.
\lean{Copying/ParametricityTailPrecompose}{152}{155}{Parametricity.TailPrecompose.focusFree}
\par\penalty0

The hole Free assertion excludes unfilled holes throughout the indicated graph and every nested box.
For the proof, constructor and boundary cases are separated explicitly.
\lean{Copying/ParametricityTailPrecompose}{157}{160}{Parametricity.TailPrecompose.holeFree}
\par\penalty0

Every original transition remains the original transition under the literal ordered family frame.
For the proof, constructor and boundary cases are separated explicitly; the defined maps and incidence equations are substituted.
\lean{Copying/ParametricityTailPrecompose}{164}{173}{Parametricity.TailPrecompose.rewrite_right}
\par\penalty0

\subsection{Parametricity Tail Tensor}
After the private read returns its numeral, the remaining family attaches only to the weakened original tree and cell. The numeral is a genuinely separate component, not an observational simplification.

This defines the selected data for the construction just described. Boundary positions are kept distinct from newly inserted internal vertices.
\lean{Copying/ParametricityTailTensor}{15}{16}{Parametricity.TailPrecompose.TensorResult.selected}
\par\penalty0

This defines the ordered source list for the construction just described. Boundary positions are kept distinct from newly inserted internal vertices.
\lean{Copying/ParametricityTailTensor}{18}{18}{Parametricity.TailPrecompose.TensorResult.source}
\par\penalty0

This defines the ordered target list for the construction just described. Boundary positions are kept distinct from newly inserted internal vertices.
\lean{Copying/ParametricityTailTensor}{19}{20}{Parametricity.TailPrecompose.TensorResult.target}
\par\penalty0

The family Occurrence map is a complete occurrence, preserving original labels and every ordered source and target attachment.
\lean{Copying/ParametricityTailTensor}{22}{22}{Parametricity.TailPrecompose.TensorResult.familyOccurrence}
\par\penalty0

The active Occurrence map is a complete occurrence, preserving original labels and every ordered source and target attachment.
\lean{Copying/ParametricityTailTensor}{23}{25}{Parametricity.TailPrecompose.TensorResult.activeOccurrence}
\par\penalty0

The remainder Occurrence map is a complete occurrence, preserving original labels and every ordered source and target attachment.
\lean{Copying/ParametricityTailTensor}{26}{28}{Parametricity.TailPrecompose.TensorResult.remainderOccurrence}
\par\penalty0

The link assertion is the compatibility law for the indicated component of this construction. It relates the actual finite coordinates or certificates, with the hypotheses stated in the linked declaration.
For the proof, constructor and boundary cases are separated explicitly; the remaining natural-number index bounds follow by linear arithmetic; the defined maps and incidence equations are substituted.
\lean{Copying/ParametricityTailTensor}{30}{39}{Parametricity.TailPrecompose.TensorResult.link}
\par\penalty0

The family remainder disjoint assertion excludes a shared edge or coordinate between the indicated independently inserted components.
For the proof, constructor and boundary cases are separated explicitly.
\lean{Copying/ParametricityTailTensor}{41}{43}{TensorResult.family_remainder_disjoint}
\par\penalty0

The continuation map is a complete occurrence, preserving original labels and every ordered source and target attachment.
\lean{Copying/ParametricityTailTensor}{45}{47}{Parametricity.TailPrecompose.TensorResult.continuation}
\par\penalty0

The active disjoint assertion excludes a shared edge or coordinate between the indicated independently inserted components.
For the proof, constructor and boundary cases are separated explicitly.
\lean{Copying/ParametricityTailTensor}{49}{54}{Parametricity.TailPrecompose.TensorResult.active_disjoint}
\par\penalty0

The occurrence map is a complete occurrence, preserving original labels and every ordered source and target attachment.
\lean{Copying/ParametricityTailTensor}{56}{58}{Parametricity.TailPrecompose.TensorResult.occurrence}
\par\penalty0

The active input assertion computes the indicated ordered boundary coordinate in the assembled graph.
For the proof, constructor and boundary cases are separated explicitly.
\lean{Copying/ParametricityTailTensor}{60}{61}{Parametricity.TailPrecompose.TensorResult.active_input}
\par\penalty0

The continuation input assertion computes the indicated ordered boundary coordinate in the assembled graph.
For the proof, constructor and boundary cases are separated explicitly; the remaining natural-number index bounds follow by linear arithmetic; the defined maps and incidence equations are substituted.
\lean{Copying/ParametricityTailTensor}{63}{66}{Parametricity.TailPrecompose.TensorResult.continuation_input}
\par\penalty0

The remainder output assertion computes the indicated ordered boundary coordinate in the assembled graph.
For the proof, constructor and boundary cases are separated explicitly; the remaining natural-number index bounds follow by linear arithmetic; the defined maps and incidence equations are substituted.
\lean{Copying/ParametricityTailTensor}{68}{74}{Parametricity.TailPrecompose.TensorResult.remainder_output}
\par\penalty0

The input assertion computes the indicated ordered boundary coordinate in the assembled graph.
For the proof, constructor and boundary cases are separated explicitly; the remaining natural-number index bounds follow by linear arithmetic; the defined maps and incidence equations are substituted.
\lean{Copying/ParametricityTailTensor}{76}{94}{Parametricity.TailPrecompose.TensorResult.input}
\par\penalty0

The output assertion computes the indicated ordered boundary coordinate in the assembled graph.
For the proof, the defined maps and incidence equations are substituted.
\lean{Copying/ParametricityTailTensor}{96}{101}{Parametricity.TailPrecompose.TensorResult.output}
\par\penalty0

The cover assertion states that every target coordinate has an original preimage under the indicated map.
For the proof, constructor and boundary cases are separated explicitly; the defined maps and incidence equations are substituted.
\lean{Copying/ParametricityTailTensor}{103}{112}{Parametricity.TailPrecompose.TensorResult.cover}
\par\penalty0

Literal reassembly into the returned numeral and the remaining family followed by the original weakened tree/private cell.
\lean{Copying/ParametricityTailTensor}{116}{117}{Parametricity.TailPrecompose.TensorResult.iso}
\par\penalty0

\subsection{Parametricity Tail Transfer}
Exact original executions and genuine directed read corrections lift through the remaining ordered Figure20 family. No restriction on the size or on the enclosing hypernet's observations is introduced by this construction.

The original mode and complete local focus are inserted at their actual family attachment. This permits intermediate internal foci, not only inputs.
\lean{Copying/ParametricityTailTransfer}{15}{17}{Parametricity.TailPrecompose.state}
\par\penalty0

The step assertion supplies the indicated actual machine transition or its transport through the common construction.
For the proof, constructor and boundary cases are separated explicitly.
\lean{Copying/ParametricityTailTransfer}{19}{25}{Parametricity.TailPrecompose.step}
\par\penalty0

Every original transition in a finite execution is preserved, hence its full numerical cost is unchanged by reattaching the remaining family.
For the proof, induction follows the original derivation or finite structure.
\lean{Copying/ParametricityTailTransfer}{29}{34}{Parametricity.TailPrecompose.run}
\par\penalty0

Assemble the local pointed witness and the unchanged complete family in the same original finite frame. The token is its actual component insertion.
\lean{Copying/ParametricityTailTransfer}{38}{48}{Parametricity.TailPrecompose.pointed}
\par\penalty0

This defines the correction data for the construction just described. Boundary positions are kept distinct from newly inserted internal vertices.
\lean{Copying/ParametricityTailTransfer}{50}{55}{Parametricity.TailPrecompose.correction}
\par\penalty0

All four actual correction links lift without assuming that filling witnesses compose transitively. Intermediate graphs and foci remain explicit.
\lean{Copying/ParametricityTailTransfer}{59}{63}{Parametricity.TailPrecompose.residual}
\par\penalty0

\subsection{Parametricity Tensor Cut}
Replacing the first complete tensor component retains the literal second component and all ordered external ports. This identifies the computed small read complement with the actual returned-value/argument parallel graph.

The replace output assertion computes the indicated ordered boundary coordinate in the assembled graph.
For the proof, the defined maps and incidence equations are substituted.
\lean{Copying/ParametricityTensorCut}{13}{25}{CompleteMatchCut.replace_output}
\par\penalty0

The selected map is a complete occurrence, preserving original labels and every ordered source and target attachment.
\lean{Copying/ParametricityTensorCut}{33}{34}{CompleteMatchCut.TensorLeft.selected}
\par\penalty0

The retained map is a complete occurrence, preserving original labels and every ordered source and target attachment.
\lean{Copying/ParametricityTensorCut}{35}{36}{CompleteMatchCut.TensorLeft.retained}
\par\penalty0

The disjoint assertion excludes a shared edge or coordinate between the indicated independently inserted components.
For the proof, constructor and boundary cases are separated explicitly.
\lean{Copying/ParametricityTensorCut}{38}{40}{CompleteMatchCut.TensorLeft.disjoint}
\par\penalty0

This defines the ordered target list for the construction just described. Boundary positions are kept distinct from newly inserted internal vertices.
\lean{Copying/ParametricityTensorCut}{42}{42}{CompleteMatchCut.TensorLeft.target}
\par\penalty0

The first map is a complete occurrence, preserving original labels and every ordered source and target attachment.
\lean{Copying/ParametricityTensorCut}{43}{43}{CompleteMatchCut.TensorLeft.first}
\par\penalty0

The second map is a complete occurrence, preserving original labels and every ordered source and target attachment.
\lean{Copying/ParametricityTensorCut}{44}{45}{CompleteMatchCut.TensorLeft.second}
\par\penalty0

The pieces disjoint assertion excludes a shared edge or coordinate between the indicated independently inserted components.
For the proof, constructor and boundary cases are separated explicitly.
\lean{Copying/ParametricityTensorCut}{47}{49}{CompleteMatchCut.TensorLeft.pieces_disjoint}
\par\penalty0

The occurrence map is a complete occurrence, preserving original labels and every ordered source and target attachment.
\lean{Copying/ParametricityTensorCut}{51}{53}{CompleteMatchCut.TensorLeft.occurrence}
\par\penalty0

The first input assertion computes the indicated ordered boundary coordinate in the assembled graph.
For the proof, the defined maps and incidence equations are substituted.
\lean{Copying/ParametricityTensorCut}{55}{57}{CompleteMatchCut.TensorLeft.first_input}
\par\penalty0

The first output assertion computes the indicated ordered boundary coordinate in the assembled graph.
For the proof, the defined maps and incidence equations are substituted.
\lean{Copying/ParametricityTensorCut}{59}{61}{CompleteMatchCut.TensorLeft.first_output}
\par\penalty0

The second input assertion computes the indicated ordered boundary coordinate in the assembled graph.
For the proof, the defined maps and incidence equations are substituted.
\lean{Copying/ParametricityTensorCut}{63}{65}{CompleteMatchCut.TensorLeft.second_input}
\par\penalty0

The second output assertion computes the indicated ordered boundary coordinate in the assembled graph.
For the proof, the defined maps and incidence equations are substituted.
\lean{Copying/ParametricityTensorCut}{67}{69}{CompleteMatchCut.TensorLeft.second_output}
\par\penalty0

The input assertion computes the indicated ordered boundary coordinate in the assembled graph.
For the proof, constructor and boundary cases are separated explicitly; the defined maps and incidence equations are substituted.
\lean{Copying/ParametricityTensorCut}{71}{82}{CompleteMatchCut.TensorLeft.input}
\par\penalty0

The output assertion computes the indicated ordered boundary coordinate in the assembled graph.
For the proof, constructor and boundary cases are separated explicitly; the defined maps and incidence equations are substituted.
\lean{Copying/ParametricityTensorCut}{84}{95}{CompleteMatchCut.TensorLeft.output}
\par\penalty0

The cover assertion states that every target coordinate has an original preimage under the indicated map.
For the proof, constructor and boundary cases are separated explicitly; the defined maps and incidence equations are substituted.
\lean{Copying/ParametricityTensorCut}{97}{106}{CompleteMatchCut.TensorLeft.cover}
\par\penalty0

The computed complement replacement is the literal tensor, with ordered boundaries fixed and every original second-component edge retained.
\lean{Copying/ParametricityTensorCut}{110}{111}{CompleteMatchCut.TensorLeft.iso}
\par\penalty0

\subsection{Parametricity Tensor Graft}
Move an actual complete two-input resource row from a selected family branch into the private sharing component. No row edge is simplified.

This defines the ordered source list for the construction just described. Boundary positions are kept distinct from newly inserted internal vertices.
\lean{Copying/ParametricityTensorGraft}{16}{16}{Parametricity.TensorGraft.source}
\par\penalty0

This defines the tail data for the construction just described. Boundary positions are kept distinct from newly inserted internal vertices.
\lean{Copying/ParametricityTensorGraft}{18}{19}{Parametricity.TensorGraft.tail}
\par\penalty0

This defines the unmerged data for the construction just described. Boundary positions are kept distinct from newly inserted internal vertices.
\lean{Copying/ParametricityTensorGraft}{21}{22}{Parametricity.TensorGraft.unmerged}
\par\penalty0

This defines the ordered target list for the construction just described. Boundary positions are kept distinct from newly inserted internal vertices.
\lean{Copying/ParametricityTensorGraft}{24}{25}{Parametricity.TensorGraft.target}
\par\penalty0

This defines the block data for the construction just described. Boundary positions are kept distinct from newly inserted internal vertices.
\lean{Copying/ParametricityTensorGraft}{27}{28}{Parametricity.TensorGraft.block}
\par\penalty0

This defines the first Block data for the construction just described. Boundary positions are kept distinct from newly inserted internal vertices.
\lean{Copying/ParametricityTensorGraft}{29}{29}{Parametricity.TensorGraft.firstBlock}
\par\penalty0

This defines the x data for the construction just described. Boundary positions are kept distinct from newly inserted internal vertices.
\lean{Copying/ParametricityTensorGraft}{30}{30}{Parametricity.TensorGraft.x}
\par\penalty0

This defines the row data for the construction just described. Boundary positions are kept distinct from newly inserted internal vertices.
\lean{Copying/ParametricityTensorGraft}{31}{31}{Parametricity.TensorGraft.row}
\par\penalty0

This defines the rest data for the construction just described. Boundary positions are kept distinct from newly inserted internal vertices.
\lean{Copying/ParametricityTensorGraft}{32}{32}{Parametricity.TensorGraft.rest}
\par\penalty0

The continuation map is a complete occurrence, preserving original labels and every ordered source and target attachment.
\lean{Copying/ParametricityTensorGraft}{33}{34}{Parametricity.TensorGraft.continuation}
\par\penalty0

The selected link assertion is the compatibility law for the indicated component of this construction. It relates the actual finite coordinates or certificates, with the hypotheses stated in the linked declaration.
For the proof, constructor and boundary cases are separated explicitly; the defined maps and incidence equations are substituted.
\lean{Copying/ParametricityTensorGraft}{36}{38}{Parametricity.TensorGraft.selected_link}
\par\penalty0

The row continuation disjoint assertion excludes a shared edge or coordinate between the indicated independently inserted components.
For the proof, constructor and boundary cases are separated explicitly.
\lean{Copying/ParametricityTensorGraft}{40}{42}{Parametricity.TensorGraft.row_continuation_disjoint}
\par\penalty0

This defines the grafted data for the construction just described. Boundary positions are kept distinct from newly inserted internal vertices.
\lean{Copying/ParametricityTensorGraft}{44}{46}{Parametricity.TensorGraft.grafted}
\par\penalty0

The x rest disjoint assertion excludes a shared edge or coordinate between the indicated independently inserted components.
For the proof, constructor and boundary cases are separated explicitly.
\lean{Copying/ParametricityTensorGraft}{48}{49}{Parametricity.TensorGraft.x_rest_disjoint}
\par\penalty0

This defines the pair data for the construction just described. Boundary positions are kept distinct from newly inserted internal vertices.
\lean{Copying/ParametricityTensorGraft}{51}{52}{Parametricity.TensorGraft.pair}
\par\penalty0

The unmerged Occurrence map is a complete occurrence, preserving original labels and every ordered source and target attachment.
\lean{Copying/ParametricityTensorGraft}{53}{54}{Parametricity.TensorGraft.unmergedOccurrence}
\par\penalty0

The x row link assertion is the compatibility law for the indicated component of this construction. It relates the actual finite coordinates or certificates, with the hypotheses stated in the linked declaration.
For the proof, constructor and boundary cases are separated explicitly; the defined maps and incidence equations are substituted.
\lean{Copying/ParametricityTensorGraft}{56}{58}{Parametricity.TensorGraft.x_row_link}
\par\penalty0

The rest continuation link assertion is the compatibility law for the indicated component of this construction. It relates the actual finite coordinates or certificates, with the hypotheses stated in the linked declaration.
For the proof, constructor and boundary cases are separated explicitly; the remaining natural-number index bounds follow by linear arithmetic; the defined maps and incidence equations are substituted.
\lean{Copying/ParametricityTensorGraft}{60}{79}{Parametricity.TensorGraft.rest_continuation_link}
\par\penalty0

The pair left output assertion computes the indicated ordered boundary coordinate in the assembled graph.
For the proof, the remaining natural-number index bounds follow by linear arithmetic; the defined maps and incidence equations are substituted.
\lean{Copying/ParametricityTensorGraft}{81}{85}{Parametricity.TensorGraft.pair_left_output}
\par\penalty0

The pair right output assertion computes the indicated ordered boundary coordinate in the assembled graph.
For the proof, the remaining natural-number index bounds follow by linear arithmetic; the defined maps and incidence equations are substituted.
\lean{Copying/ParametricityTensorGraft}{87}{91}{Parametricity.TensorGraft.pair_right_output}
\par\penalty0

The graft left input assertion computes the indicated ordered boundary coordinate in the assembled graph.
For the proof, the remaining natural-number index bounds follow by linear arithmetic; the defined maps and incidence equations are substituted.
\lean{Copying/ParametricityTensorGraft}{93}{99}{Parametricity.TensorGraft.graft_left_input}
\par\penalty0

The graft right input assertion computes the indicated ordered boundary coordinate in the assembled graph.
For the proof, the remaining natural-number index bounds follow by linear arithmetic; the defined maps and incidence equations are substituted.
\lean{Copying/ParametricityTensorGraft}{101}{114}{Parametricity.TensorGraft.graft_right_input}
\par\penalty0

The boundary assertion is the compatibility law for the indicated component of this construction. It relates the actual finite coordinates or certificates, with the hypotheses stated in the linked declaration.
For the proof, constructor and boundary cases are separated explicitly; the remaining natural-number index bounds follow by linear arithmetic; the defined maps and incidence equations are substituted.
\lean{Copying/ParametricityTensorGraft}{116}{129}{Parametricity.TensorGraft.boundary}
\par\penalty0

The pair graft disjoint assertion excludes a shared edge or coordinate between the indicated independently inserted components.
For the proof, constructor and boundary cases are separated explicitly.
\lean{Copying/ParametricityTensorGraft}{131}{136}{Parametricity.TensorGraft.pair_graft_disjoint}
\par\penalty0

The occurrence map is a complete occurrence, preserving original labels and every ordered source and target attachment.
\lean{Copying/ParametricityTensorGraft}{138}{140}{Parametricity.TensorGraft.occurrence}
\par\penalty0

The x input assertion computes the indicated ordered boundary coordinate in the assembled graph.
For the proof, constructor and boundary cases are separated explicitly; the defined maps and incidence equations are substituted.
\lean{Copying/ParametricityTensorGraft}{142}{151}{Parametricity.TensorGraft.x_input}
\par\penalty0

The rest input assertion computes the indicated ordered boundary coordinate in the assembled graph.
For the proof, constructor and boundary cases are separated explicitly; the defined maps and incidence equations are substituted.
\lean{Copying/ParametricityTensorGraft}{153}{158}{Parametricity.TensorGraft.rest_input}
\par\penalty0

The input assertion computes the indicated ordered boundary coordinate in the assembled graph.
For the proof, constructor and boundary cases are separated explicitly; the defined maps and incidence equations are substituted.
\lean{Copying/ParametricityTensorGraft}{160}{171}{Parametricity.TensorGraft.input}
\par\penalty0

The output assertion computes the indicated ordered boundary coordinate in the assembled graph.
For the proof, the remaining natural-number index bounds follow by linear arithmetic; the defined maps and incidence equations are substituted.
\lean{Copying/ParametricityTensorGraft}{173}{182}{Parametricity.TensorGraft.output}
\par\penalty0

The cover assertion states that every target coordinate has an original preimage under the indicated map.
For the proof, constructor and boundary cases are separated explicitly; the defined maps and incidence equations are substituted.
\lean{Copying/ParametricityTensorGraft}{184}{196}{Parametricity.TensorGraft.cover}
\par\penalty0

The entire actual two-input resource row moves into the private sharing component, preserving every original vertex attachment and external order.
\lean{Copying/ParametricityTensorGraft}{200}{201}{Parametricity.TensorGraft.iso}
\par\penalty0

\subsection{Retained Replacement}
Replacing a retained primitive edge preserves every original shallow and deep template component. The paired successor is a genuine recursive filling, pointed at the unchanged original host wire.

Replacing the unique edge of an actual corolla leaves precisely the replacement graph, including all of its ordered boundary attachments.
\lean{Copying/RetainedReplacement}{15}{34}{LabelCopy.plugIso}
\par\penalty0

The selected source corolla is replaced by the whole original allocation prefix. Every other component is literally the original supplied filler.
\lean{Copying/RetainedReplacement}{47}{47}{Allocation.Retained.family}
\par\penalty0

Canonical allocation of a retained ref corolla has the same successor as changing only its member of the original filling family.
\lean{Copying/RetainedReplacement}{51}{80}{Allocation.Retained.canonicalIso}
\par\penalty0

A rooted ref request executes to the shared changed-family representation. Only the selected original ref corolla is replaced; the operand remains external.
For the proof, the defined maps and incidence equations are substituted.
\lean{Copying/RetainedReplacement}{85}{114}{Allocation.Retained.rooted_rewrite}
\par\penalty0

Actual allocation prefixes carry no focus, including their canonical atom contraction and weakening.
For the proof, constructor and boundary cases are separated explicitly.
\lean{Copying/RetainedReplacement}{126}{129}{Allocation.Retained.filler_free}
\par\penalty0

Replace only the selected ref component by the same genuine allocation prefix on both sides; every other original recursive certificate is retained.
\lean{Copying/RetainedReplacement}{133}{143}{Allocation.Retained.witnesses}
\par\penalty0

The residual keeps the selected old host token, even when the stable argument is a lambda inside a multi-input shallow family.
\lean{Copying/RetainedReplacement}{147}{155}{Allocation.Retained.pointed}
\par\penalty0

\subsection{Tensor Association}
Association of the literal disjoint tensor assembly. It changes only parenthesization, preserving the complete ordered external boundary.

The first map is a complete occurrence, preserving original labels and every ordered source and target attachment.
\lean{Copying/TensorAssociation}{11}{12}{TensorAssociation.first}
\par\penalty0

The second map is a complete occurrence, preserving original labels and every ordered source and target attachment.
\lean{Copying/TensorAssociation}{13}{15}{TensorAssociation.second}
\par\penalty0

The third map is a complete occurrence, preserving original labels and every ordered source and target attachment.
\lean{Copying/TensorAssociation}{16}{18}{TensorAssociation.third}
\par\penalty0

The first second assertion is the compatibility law for the indicated component of this construction. It relates the actual finite coordinates or certificates, with the hypotheses stated in the linked declaration.
For the proof, constructor and boundary cases are separated explicitly.
\lean{Copying/TensorAssociation}{20}{21}{TensorAssociation.first_second}
\par\penalty0

The first Two map is a complete occurrence, preserving original labels and every ordered source and target attachment.
\lean{Copying/TensorAssociation}{22}{23}{TensorAssociation.firstTwo}
\par\penalty0

The first Two third assertion is the compatibility law for the indicated component of this construction. It relates the actual finite coordinates or certificates, with the hypotheses stated in the linked declaration.
For the proof, constructor and boundary cases are separated explicitly.
\lean{Copying/TensorAssociation}{24}{25}{TensorAssociation.firstTwo_third}
\par\penalty0

The occurrence map is a complete occurrence, preserving original labels and every ordered source and target attachment.
\lean{Copying/TensorAssociation}{26}{28}{TensorAssociation.occurrence}
\par\penalty0

The inputs assertion computes the indicated ordered boundary coordinate in the assembled graph.
For the proof, the preceding component lemmas are composed or specialized at the stated parameters.
\lean{Copying/TensorAssociation}{30}{33}{TensorAssociation.inputs}
\par\penalty0

The outputs assertion computes the indicated ordered boundary coordinate in the assembled graph.
For the proof, the preceding component lemmas are composed or specialized at the stated parameters.
\lean{Copying/TensorAssociation}{35}{38}{TensorAssociation.outputs}
\par\penalty0

The cover assertion states that every target coordinate has an original preimage under the indicated map.
For the proof, constructor and boundary cases are separated explicitly; the defined maps and incidence equations are substituted.
\lean{Copying/TensorAssociation}{40}{47}{TensorAssociation.cover}
\par\penalty0

The isomorphism retains each original component and its ordered ports.
\lean{Copying/TensorAssociation}{50}{52}{TensorAssociation.iso}
\par\penalty0

The input assertion computes the indicated ordered boundary coordinate in the assembled graph.
For the proof, constructor and boundary cases are separated explicitly; the remaining natural-number index bounds follow by linear arithmetic; the defined maps and incidence equations are substituted.
\lean{Copying/TensorAssociation}{54}{84}{TensorAssociation.input}
\par\penalty0

The output assertion computes the indicated ordered boundary coordinate in the assembled graph.
For the proof, constructor and boundary cases are separated explicitly; the remaining natural-number index bounds follow by linear arithmetic; the defined maps and incidence equations are substituted.
\lean{Copying/TensorAssociation}{86}{116}{TensorAssociation.output}
\par\penalty0

